% Logique, argumentation et démonstration — Philosophie Tle Tech — Cours signé OnBuch+
% Programme officiel MINESEC. Compiler avec : tectonic N13-logique-argumentation-demonstration.tex
\newcommand{\DOCMATIERE}{Philosophie}
\newcommand{\DOCNIVEAU}{Tle Tech}
\newcommand{\DOCMODULE}{Philosophie – classe de Terminale technique}
\newcommand{\DOCLECON}{N13}
\newcommand{\DOCTITRE}{Logique, argumentation et démonstration}
% OnBuch+ — préambule « cours signés » ENRICHI (compilé avec tectonic / XeLaTeX)
% Système visuel « Pop » d'OnBuch+ : fond crème (jamais blanc), bordures encre
% épaisses, ombres « dures » décalées, titres Archivo Black en capitales.
% Version MATHÉMATIQUES : graphes (pgfplots/tikz), boîtes pédagogiques
% (définition, propriété, méthode, attention, le savais-tu, exercice résolu,
% exercice, TP, bilan), page de garde et signature OnBuch+ ancrée en bas.
%
% Macros attendues dans main.tex AVANT \input{preamble} :
%   \DOCMATIERE  \DOCNIVEAU  \DOCMODULE  \DOCLECON (numéro)  \DOCTITRE
\documentclass[11pt]{article}

\usepackage{fontspec}
\usepackage[french]{babel}
\usepackage[a4paper,margin=2cm,top=3.4cm,bottom=2.8cm,headheight=1.4cm,footskip=1.3cm]{geometry}
\usepackage{amsmath,amssymb,amsfonts}
\usepackage{mathtools} % \xmapsto, \coloneqq... souvent utilisés par les modèles
\usepackage{cancel} % simplifications barrées (bilans de forces)
% \abs{z} (module, valeur absolue) et \norm{u} (norme), utilisés par les modèles
\providecommand{\abs}{}\let\abs\relax\DeclarePairedDelimiter{\abs}{\lvert}{\rvert}
\providecommand{\norm}{}\let\norm\relax\DeclarePairedDelimiter{\norm}{\lVert}{\rVert}
\usepackage[table]{xcolor} % [table] : fournit aussi \rowcolors (alternance de lignes)
\usepackage{pagecolor}
\usepackage{tikz}
\usetikzlibrary{arrows.meta,positioning,calc,shapes.geometric,shapes.misc,shapes.arrows,decorations.pathmorphing,decorations.pathreplacing,decorations.markings,patterns,shadows,mindmap,trees,fit,backgrounds,matrix,intersections,angles,quotes}
\usepackage{pgfplots}
\usepackage[version=4]{mhchem} % équations nucléaires \ce{^{235}_{92}U}
\usepackage[european,siunitx]{circuitikz} % schémas électriques
\pgfplotsset{compat=1.18}
\usepgfplotslibrary{fillbetween,groupplots} % aires entre courbes (calcul intégral)
\usepackage{enumitem}
\usepackage{fancyhdr}
% [most] charge toutes les bibliothèques tcolorbox (listings, minted, raster,
% documentation...) et gonfle inutilement la mémoire TeX sur un document long
% (~300 boîtes) : on ne charge que ce qui est réellement utilisé.
\usepackage[skins,breakable]{tcolorbox}
\usepackage{graphicx}
\usepackage{colortbl}
\usepackage{booktabs}
\usepackage{tabularx}
\usepackage{multirow}
\usepackage{diagbox} % case d'en-tête diagonale des tableaux à double entrée
% Colonnes extensibles centrée / gauche / droite (C, L, R), souvent utilisées par les modèles
\newcolumntype{C}{>{\centering\arraybackslash}X}
\newcolumntype{L}{>{\raggedright\arraybackslash}X}
\newcolumntype{R}{>{\raggedleft\arraybackslash}X}
\usepackage{array}
\usepackage{multicol}
\usepackage{siunitx}
% parse-numbers=false : accepte aussi \SI{2,5 \times 10^{-3}}{...}
\sisetup{locale=FR,output-decimal-marker={,},group-minimum-digits=4,parse-numbers=false}
\DeclareSIUnit\ohm{\ensuremath{\symup{\Omega}}} % U+2126 absent de la police
\DeclareSIUnit\division{div} % réglages d'oscilloscope (ms/div, V/div)
\DeclareSIUnit\tr{tr} % tours (vitesse de rotation, ex. \SI{1500}{\tr\per\minute})
\DeclareSIUnit\molar{M} % concentration molaire (ex. \SI{3}{\molar}, \SI{1}{\micro\molar})
\DeclareSIUnit\an{an} % année (le modèle confond parfois avec \year, un compteur TeX) — ex. \SI{5}{\kilo\meter\per\an}
\DeclareSIUnit\FCFA{FCFA} % franc CFA (ex. \SI{500}{\FCFA\per\kilo\gram})
\DeclareSIUnit\degreeBrix{\textdegree Brix} % teneur en sucre d'un sirop/jus (réfractométrie)
\DeclareSIUnit\month{mois} % le modèle utilise parfois \month, un compteur TeX, comme unité de durée
\DeclareSIUnit\microliter{\micro\liter} % le modèle écrit parfois \microliter en un seul token
\DeclareSIUnit\million{millions} % dénombrement (ex. \SI{15}{\million\per\mL} spermatozoïdes)
\usepackage{qrcode}
% \degree : commande fréquemment utilisée par les modèles pour un angle ou une
% température en dehors de \SI{}{} (non fournie par les paquets chargés).
\providecommand{\degree}{^{\circ}}
% \qed (fin de démonstration), utilisé par les modèles sans amsthm
\providecommand{\qed}{\hfill$\square$}
% \barre : double barre des valeurs interdites dans un tableau de variations (tkz-tab)
\providecommand{\barre}{\ensuremath{\|}}
% environnement proof (amsthm non chargé) utilisé par les modèles
\ifdefined\proof\else\newenvironment{proof}[1][Démonstration]{\par\noindent\textit{#1.}\ }{\qed\par}\fi
\usepackage{lastpage}
\usepackage{hyperref}
\hypersetup{hidelinks}

% ============================================================
% Palette « Pop » OnBuch+ — mêmes hex que la classe P (plus_ui.dart)
% ============================================================
\definecolor{poporange}{HTML}{F59321}
\definecolor{popdark}{HTML}{E07A0C}
\definecolor{popcream}{HTML}{FFF6E8}
\definecolor{popink}{HTML}{1C1714}
\definecolor{poppurple}{HTML}{7A5AE0}
\definecolor{popgreen}{HTML}{2E9E6A}
\definecolor{poppink}{HTML}{E0457B}
\definecolor{popblue}{HTML}{2F7DE1}
\definecolor{popgold}{HTML}{C98A12}
\definecolor{popmuted}{HTML}{8A7F76}
\definecolor{popline}{HTML}{EADFD2}
\definecolor{popgray}{HTML}{8A7F76}
\colorlet{popgrey}{popgray}
% Alias tolérants (le modèle nomme parfois les couleurs différemment)
\colorlet{popwhite}{white}
\colorlet{popblack}{popink}
\colorlet{popred}{poppink}
\colorlet{popyellow}{popgold}
% Teintes claires (fonds de tableaux/encadrés) — le modèle nomme parfois
% les couleurs avec un suffixe L (ex. popblueL) sans les avoir définies.
\colorlet{poporangeL}{poporange!20}
\colorlet{popdarkL}{popdark!20}
\colorlet{poppurpleL}{poppurple!20}
\colorlet{popgreenL}{popgreen!20}
\colorlet{poppinkL}{poppink!20}
\colorlet{popblueL}{popblue!20}
\colorlet{popgoldL}{popgold!20}
\colorlet{popredL}{popred!20}
\colorlet{popgrayL}{popgray!20}
% Teintes claires pour fonds de tableaux / zones
\colorlet{popblueL}{popblue!10!white}
\colorlet{poporangeL}{poporange!14!white}
\colorlet{popgreenL}{popgreen!12!white}
\colorlet{poppurpleL}{poppurple!12!white}
\colorlet{poppinkL}{poppink!12!white}
\colorlet{popgoldL}{popgold!14!white}

\pagecolor{popcream}

% ============================================================
% Typographie
% ============================================================
\defaultfontfeatures{Ligatures=TeX}
\setmainfont{PlusJakartaSans-Medium.ttf}[Path=../../../fonts/,BoldFont=PlusJakartaSans-Bold.ttf]
% Maths en sans-serif (Fira Math) avec chiffres et lettres droites de Plus
% Jakarta, pour que les formules se fondent dans le texte.
\usepackage[math-style=ISO,bold-style=ISO]{unicode-math}
\setmathfont{FiraMath-Regular.otf}[Path=../../../fonts/]
\setmathfont{PlusJakartaSans-Medium.ttf}[Path=../../../fonts/,range={up/num,up/{Latin,latin},\mathup}]
\usepackage{icomma} % virgule décimale sans espace en mode math
% Caractères mathématiques Unicode absents des polices de texte (Plus Jakarta,
% Archivo Black) : rendus par leur équivalent mathématique, en texte comme en
% formule — sinon ils s'affichent en carré vide (ex. « Arithmétique dans ℤ »).
\usepackage{newunicodechar}
\newunicodechar{ℝ}{\ensuremath{\mathbb{R}}}
\newunicodechar{ℤ}{\ensuremath{\mathbb{Z}}}
\newunicodechar{ℂ}{\ensuremath{\mathbb{C}}}
\newunicodechar{ℕ}{\ensuremath{\mathbb{N}}}
\newunicodechar{ℚ}{\ensuremath{\mathbb{Q}}}
\newunicodechar{𝔻}{\ensuremath{\mathbb{D}}}
\newunicodechar{α}{\ensuremath{\alpha}}
\newunicodechar{β}{\ensuremath{\beta}}
\newunicodechar{γ}{\ensuremath{\gamma}}
\newunicodechar{δ}{\ensuremath{\delta}}
\newunicodechar{ε}{\ensuremath{\varepsilon}}
\newunicodechar{θ}{\ensuremath{\theta}}
\newunicodechar{λ}{\ensuremath{\lambda}}
\newunicodechar{μ}{\ensuremath{\mu}}
\newunicodechar{σ}{\ensuremath{\sigma}}
\newunicodechar{τ}{\ensuremath{\tau}}
\newunicodechar{φ}{\ensuremath{\varphi}}
\newunicodechar{ψ}{\ensuremath{\psi}}
\newunicodechar{ω}{\ensuremath{\omega}}
\newunicodechar{Σ}{\ensuremath{\Sigma}}
% majuscules produites par \MakeUppercase dans les titres (α-aminés -> Α)
\newunicodechar{Α}{\ensuremath{\alpha}}
\newunicodechar{Β}{\ensuremath{\beta}}
\newunicodechar{Γ}{\ensuremath{\Gamma}}
\newunicodechar{Δ}{\ensuremath{\Delta}}
\newunicodechar{Ε}{\ensuremath{\varepsilon}}
\newunicodechar{Θ}{\ensuremath{\Theta}}
\newunicodechar{Λ}{\ensuremath{\Lambda}}
\newunicodechar{Μ}{\ensuremath{\mu}}
\newunicodechar{Τ}{\ensuremath{\tau}}
\newunicodechar{Φ}{\ensuremath{\Phi}}
\newunicodechar{Ψ}{\ensuremath{\Psi}}
\newunicodechar{Ω}{\ensuremath{\Omega}}
\newunicodechar{↦}{\ensuremath{\mapsto}}
\newunicodechar{∈}{\ensuremath{\in}}
\newunicodechar{∉}{\ensuremath{\notin}}
\newunicodechar{∧}{\ensuremath{\wedge}}
\newunicodechar{⇔}{\ensuremath{\Leftrightarrow}}
\newunicodechar{⟺}{\ensuremath{\Longleftrightarrow}}
\newunicodechar{⇒}{\ensuremath{\Rightarrow}}
\newunicodechar{⟹}{\ensuremath{\Longrightarrow}}
\newunicodechar{∘}{\ensuremath{\circ}}
\newunicodechar{∪}{\ensuremath{\cup}}
\newunicodechar{∩}{\ensuremath{\cap}}
\newunicodechar{≡}{\ensuremath{\equiv}}
\newunicodechar{⊂}{\ensuremath{\subset}}
\newunicodechar{⊥}{\ensuremath{\perp}}
\newunicodechar{∃}{\ensuremath{\exists}}
\newunicodechar{∀}{\ensuremath{\forall}}
\newunicodechar{∛}{\ensuremath{\sqrt[3]{\,}}}
\newunicodechar{∜}{\ensuremath{\sqrt[4]{\,}}}
\newunicodechar{⃗}{\ensuremath{{}^{\to}}}
\newunicodechar{ⁿ}{\ifmmode^{n}\else\textsuperscript{\ensuremath{n}}\fi}
\newunicodechar{ˣ}{\ifmmode^{x}\else\textsuperscript{\ensuremath{x}}\fi}
\newunicodechar{ᵇ}{\ifmmode^{b}\else\textsuperscript{\ensuremath{b}}\fi}
\newunicodechar{ʳ}{\ifmmode^{r}\else\textsuperscript{\ensuremath{r}}\fi}
\newunicodechar{ᵘ}{\ifmmode^{u}\else\textsuperscript{\ensuremath{u}}\fi}
\newunicodechar{ʸ}{\ifmmode^{y}\else\textsuperscript{\ensuremath{y}}\fi}
\newunicodechar{ᵃ}{\ifmmode^{a}\else\textsuperscript{\ensuremath{a}}\fi}
\newunicodechar{ᵉ}{\ifmmode^{e}\else\textsuperscript{\ensuremath{e}}\fi}
\newunicodechar{ᵏ}{\ifmmode^{k}\else\textsuperscript{\ensuremath{k}}\fi}
\newunicodechar{ᵖ}{\ifmmode^{p}\else\textsuperscript{\ensuremath{p}}\fi}
\newunicodechar{ᴺ}{\ifmmode^{N}\else\textsuperscript{\ensuremath{N}}\fi}
\newunicodechar{ᶜ}{\ifmmode^{c}\else\textsuperscript{\ensuremath{c}}\fi}
\newunicodechar{ʰ}{\ifmmode^{h}\else\textsuperscript{\ensuremath{h}}\fi}
\newunicodechar{ᵠ}{\ifmmode^{q}\else\textsuperscript{\ensuremath{q}}\fi}
\newunicodechar{ₙ}{\ifmmode_{n}\else\textsubscript{\ensuremath{n}}\fi}
\newunicodechar{ₐ}{\ifmmode_{a}\else\textsubscript{\ensuremath{a}}\fi}
\newunicodechar{ᵢ}{\ifmmode_{i}\else\textsubscript{\ensuremath{i}}\fi}
\newunicodechar{ᵣ}{\ifmmode_{r}\else\textsubscript{\ensuremath{r}}\fi}
\newunicodechar{ₖ}{\ifmmode_{k}\else\textsubscript{\ensuremath{k}}\fi}
\newunicodechar{ᵦ}{\ifmmode_{\beta}\else\textsubscript{\ensuremath{\beta}}\fi}
\newunicodechar{ₓ}{\ifmmode_{x}\else\textsubscript{\ensuremath{x}}\fi}
\newfontfamily\popdisplay{ArchivoBlack-Regular.ttf}[Path=../../../fonts/]
\newfontfamily\poplogo{SpaceGrotesk-Bold.ttf}[Path=../../../fonts/]
\newfontfamily\popsemi{PlusJakartaSans-SemiBold.ttf}[Path=../../../fonts/]
\newfontfamily\popxbold{PlusJakartaSans-ExtraBold.ttf}[Path=../../../fonts/]
\newfontfamily\popxboldls{PlusJakartaSans-ExtraBold.ttf}[Path=../../../fonts/,LetterSpace=3.0]

\setlist[enumerate]{leftmargin=1.7em,itemsep=5pt,topsep=4pt}
\setlist[itemize]{leftmargin=1.7em,itemsep=4pt,topsep=4pt,label={\color{poporange}\textbullet}}
\setlength{\parindent}{0pt}
\setlength{\parskip}{5pt}
\renewcommand{\arraystretch}{1.25}

% pgfplots : style Pop par défaut
\pgfplotsset{
  every axis/.append style={axis line style={popink,line width=1pt},tick style={popink},
    grid style={popline,line width=0.5pt},label style={font=\small},tick label style={font=\footnotesize}},
  popcurve/.style={poporange,line width=1.8pt,no markers,smooth},
}
% Même style disponible dans un \draw TikZ simple (hors pgfplots)
\tikzset{popcurve/.style={poporange,line width=1.8pt}}
\tikzset{popgrid/.style={popline,very thin}}
\colorlet{popcurve}{poporange} % le nom du style est parfois utilisé comme couleur

% ============================================================
% Pill / squiggle
% ============================================================
\newcommand{\poppill}[3]{%
  \tikz[baseline=-0.55ex]\node[draw=popink,line width=1.2pt,rounded corners=2.6pt,%
    fill=#2,inner xsep=6.5pt,inner ysep=3pt]{%
    \popxboldls\fontsize{7}{7}\selectfont\color{#3}\MakeUppercase{#1}};%
}
\newcommand{\popsquiggle}[1]{%
  \tikz[baseline=-0.4ex]{
    \draw[#1,line width=2.6pt,line cap=round]
      plot [smooth] coordinates {(0,0.09) (0.34,0.01) (0.68,0.21) (1.02,0.01) (1.36,0.10)};
  }%
}
% Mot-clé surligné (marqueur orange sous le texte)
\newcommand{\cle}[1]{{\popxbold\color{popdark}#1}}

% ============================================================
% En-tête / pied
% ============================================================
\pagestyle{fancy}
\fancyhf{}
\renewcommand{\headrulewidth}{0pt}
\renewcommand{\footrulewidth}{0pt}
\lhead{\raisebox{-0.15ex}{\poplogo\fontsize{13}{13}\selectfont\color{popink}OnBuch\color{poporange}+}}
\rhead{\poppill{\DOCMATIERE\ \textbullet\ \DOCNIVEAU\ \textbullet\ Leçon \DOCLECON}{white}{popink}}
\lfoot{\popxbold\fontsize{7.6}{7.6}\selectfont\color{popmuted}\MakeUppercase{Cours signé OnBuch+}\ \textbullet\ {\popsemi\DOCTITRE}}
\rfoot{\tikz[baseline=-0.6ex]\node[rounded corners=8.5pt,fill=popink,minimum height=17pt,minimum width=17pt,inner xsep=5pt,inner sep=0]{\popxbold\fontsize{8}{8}\selectfont\color{popcream}\thepage\,/\,\pageref*{LastPage}};}

% ============================================================
% Titres
% ============================================================
\newcommand{\chaptitre}[2]{%
  \begin{tcolorbox}[enhanced,colback=poporange,colframe=popink,boxrule=2.6pt,arc=11pt,
      drop shadow=popink,left=15pt,right=15pt,top=12pt,bottom=11pt,before skip=3pt,after skip=18pt]
    \poppill{#2}{popink}{popcream}\par\vspace{9pt}
    {\raggedright\hyphenpenalty=10000\exhyphenpenalty=10000\popdisplay\fontsize{22}{24}\selectfont\color{popcream}\MakeUppercase{#1}\par}
  \end{tcolorbox}
}
% Section numérotée : pastille orange avec le numéro + titre Archivo
\newcounter{popsec}
\newcommand{\coursec}[1]{%
  \stepcounter{popsec}\setcounter{popsub}{0}%
  \par\vspace{16pt}\noindent
  \begin{minipage}{\linewidth}
  \tikz[baseline=-0.7ex]\node[draw=popink,line width=1.6pt,fill=poporange,rounded corners=4pt,minimum size=22pt,inner sep=0]{\popdisplay\fontsize{12}{12}\selectfont\color{popcream}\thepopsec};%
  \hspace{8pt}{\popdisplay\fontsize{15}{17}\selectfont\color{popink}\MakeUppercase{#1}}\par
  \vspace{4pt}\noindent{\color{poporange}\rule{36pt}{3.6pt}}
  \end{minipage}\par\nopagebreak\vspace{8pt}
}
\newcounter{popsub}
\newcommand{\courssub}[1]{%
  \stepcounter{popsub}%
  \par\vspace{9pt}\noindent{\popxbold\fontsize{11.5}{13}\selectfont\color{popink}{\color{poporange}\thepopsec.\thepopsub}\hspace{6pt}#1}\par\nopagebreak\vspace{4pt}
}

% ============================================================
% Boîtes pédagogiques — même recette : fond blanc, bordure épaisse colorée,
% ombre dure encre, badge plein. Titre optionnel [#1].
% ============================================================
\tcbset{popbase/.style={breakable,enhanced,colback=white,boxrule=2.3pt,arc=9pt,
  shadow={3pt}{-3pt}{0pt}{popink},left=14pt,right=14pt,top=11pt,bottom=11pt,before skip=11pt,after skip=15pt,
  fontupper=\popsemi\fontsize{10.6}{14.5}\selectfont\color{popink},
  fontlower=\popsemi\fontsize{10.6}{14.5}\selectfont\color{popink}}}
% Coche dessinée (le glyphe ✓ n'existe pas dans les polices du document)
\newcommand{\popcheck}[1]{\tikz[baseline=-0.5ex]\draw[#1,line width=1.6pt,line cap=round,line join=round](-0.13,0.0)--(-0.04,-0.09)--(0.13,0.1);}
\providecommand{\coche}{\popcheck{popgreen}}
\providecommand{\resultat}[1]{\par\smallskip\noindent\fcolorbox{popgreen}{popgreenL}{\parbox{\dimexpr\linewidth-2\fboxsep-2\fboxrule\relax}{\textbf{Résultat.} #1}}\par\smallskip}
\newcommand{\popboxhead}[3]{\poppill{#1}{#2}{white}\ifx\relax#3\relax\else\hspace{6pt}{\popxbold\fontsize{10.6}{12}\selectfont\color{popink}#3}\fi\par\vspace{7pt}}

\newtcolorbox{aretenir}[1][]{popbase,colframe=popgreen,before upper={\popboxhead{À retenir}{popgreen}{#1}}}
\newtcolorbox{exemplebox}[1][]{popbase,colframe=poporange,before upper={\popboxhead{Exemple}{poporange}{#1}}}
\newtcolorbox{definition}[1][]{popbase,colframe=popblue,before upper={\popboxhead{Définition}{popblue}{#1}}}
\newtcolorbox{propriete}[1][]{popbase,colframe=poppurple,before upper={\popboxhead{Propriété}{poppurple}{#1}}}
\newtcolorbox{methode}[1][]{popbase,colframe=popgold,before upper={\popboxhead{Méthode}{popgold}{#1}}}
\newtcolorbox{attention}[1][]{popbase,colframe=poppink,before upper={\popboxhead{Attention}{poppink}{#1}}}
\newtcolorbox{experience}[1][]{popbase,colframe=popblue,colback=popblueL,before upper={\popboxhead{Expérience}{popblue}{#1}}}
% « Le savais-tu ? » : carte violette pleine (contexte, culture, Cameroun)
\newtcolorbox{savaistu}[1][]{popbase,colback=poppurple,colframe=popink,
  fontupper=\popsemi\fontsize{10.4}{14.2}\selectfont\color{white},
  before upper={\poppill{Le savais-tu\,?}{white}{poppurple}\ifx\relax#1\relax\else\hspace{6pt}{\popxbold\color{white}#1}\fi\par\vspace{7pt}}}
% Exercice résolu : énoncé en haut, \tcblower puis la solution sur fond crème
\newcounter{popexo}\newcounter{popexr}
\newtcolorbox{exoresolu}[1][]{popbase,colframe=poporange,colbacklower=popcream,
  segmentation style={popink,line width=1.2pt,dashed},
  before upper={\stepcounter{popexr}\popboxhead{Exercice résolu \thepopexr}{poporange}{#1}},
  before lower={\poppill{Solution}{popink}{popcream}\par\vspace{6pt}}}
% Exercice (non résolu en place) — la correction est dans la partie Corrigés
\newtcolorbox{exercice}[1][]{popbase,colframe=popink,
  before upper={\stepcounter{popexo}\popboxhead{Exercice \thepopexo}{popink}{#1}}}
% Corrigé
\newtcolorbox{corrige}[1][]{popbase,colframe=popgreen,colback=popgreenL,
  before upper={\popboxhead{Corrigé}{popgreen}{#1}}}
% Objectifs / prérequis / bilan
\newtcolorbox{objectifs}{popbase,colframe=popink,colback=poporangeL,
  before upper={\popboxhead{Objectifs de la leçon}{popink}{}}}
\newtcolorbox{prerequis}{popbase,colframe=popink,colback=popblueL,
  before upper={\popboxhead{Prérequis}{popblue}{}}}
\newtcolorbox{bilan}{popbase,colframe=popink,colback=white,boxrule=2.8pt,
  before upper={\popboxhead{Fiche bilan}{poporange}{L'essentiel à connaître pour l'examen}}}
% Figure encadrée avec légende
\newtcolorbox{popfigure}[1][]{enhanced,breakable=false,colback=white,colframe=popline,boxrule=1.4pt,arc=8pt,
  left=10pt,right=10pt,top=10pt,bottom=8pt,before skip=10pt,after skip=12pt,halign=center,
  lower separated=false,halign lower=center,
  fontlower=\popsemi\fontsize{9}{11.5}\selectfont\color{popmuted},
  #1}
\newcommand{\legende}[1]{\par\vspace{6pt}{\popsemi\fontsize{9}{11.5}\selectfont\color{popmuted}#1}}

% ============================================================
% Signature OnBuch+
% ============================================================
\newcommand{\poplogoinline}[1]{{\poplogo\fontsize{#1}{#1}\selectfont\color{popink}OnBuch\color{poporange}+}}
\newcommand{\signatureonbuch}{%
  \begin{tcolorbox}[enhanced,colback=popink,colframe=popink,boxrule=2.2pt,arc=10pt,
      drop shadow=poporange,left=14pt,right=14pt,top=10pt,bottom=10pt,before skip=0pt,after skip=0pt]
    \tikz[baseline=-0.7ex]\node[circle,fill=popgreen,draw=popcream,line width=1.3pt,minimum size=22pt,inner sep=0]{\popcheck{white}};%
    \hspace{9pt}\begin{minipage}[c]{0.62\linewidth}
      {\popxbold\fontsize{12}{13}\selectfont\color{popcream}Signé\ \textbullet\ {\poplogo OnBuch\color{poporange}+}}\par\vspace{2pt}
      {\raggedright\popsemi\fontsize{8}{10.5}\selectfont\color{popline}\MakeUppercase{Équipe pédagogique OnBuch+}\\\MakeUppercase{Conforme au programme officiel MINESEC}\par}
    \end{minipage}\hfill
    {\poplogo\fontsize{22}{22}\selectfont\color{popcream}OnBuch\color{poporange}+}
  \end{tcolorbox}%
}

% ============================================================
% Page de garde
% #1 titre · #2 matière · #3 niveau · #4 module · #5 n° leçon · #6 durée ·
% #7 description / objectifs (texte ou itemize)
% ============================================================
\newcommand{\pagegarde}[7]{%
  \thispagestyle{empty}
  \newgeometry{a4paper,margin=1.7cm,top=1.6cm,bottom=1.4cm}%
  \begin{tikzpicture}[remember picture,overlay]
    \fill[poporange!18] ([xshift=-3.2cm,yshift=-3.6cm]current page.north east) circle (4.6cm);
    \fill[poppurple!14] ([xshift=2.4cm,yshift=7.6cm]current page.south west) circle (3.8cm);
  \end{tikzpicture}%
  {\poplogo\fontsize{30}{30}\selectfont\color{popink}OnBuch\color{poporange}+}\hfill
  \raisebox{0.6ex}{\poppill{Cours signé \textbullet\ Premium}{popink}{popcream}}\par
  \vspace{1pt}\popsquiggle{poppurple}\par
  \vspace{8pt}
  \begin{tcolorbox}[enhanced,colback=poporange,colframe=popink,boxrule=2.8pt,arc=13pt,
      drop shadow=popink,left=18pt,right=18pt,top=12pt,bottom=13pt,after skip=10pt]
    \poppill{#2 \textbullet\ #3}{popink}{popcream}\hspace{5pt}\poppill{#4}{white}{popink}\par\vspace{8pt}
    {\popdisplay\fontsize{14}{14}\selectfont\color{popink}LEÇON #5}\par\vspace{4pt}
    {\raggedright\hyphenpenalty=10000\exhyphenpenalty=10000\popdisplay\fontsize{25}{27}\selectfont\color{popcream}\MakeUppercase{#1}\par}
    \vspace{8pt}
    {\popxbold\fontsize{9.5}{9.5}\selectfont\color{popink}\MakeUppercase{Durée conseillée : #6}}
  \end{tcolorbox}
  \begin{tcolorbox}[enhanced,colback=white,colframe=popink,boxrule=2.2pt,arc=10pt,
      drop shadow=popink,left=16pt,right=16pt,top=10pt,bottom=10pt,after skip=8pt]
    \poppill{Dans cette leçon}{poppurple}{white}\par\vspace{5pt}
    \popsemi\fontsize{9.3}{12.6}\selectfont\color{popink}#7
  \end{tcolorbox}
  \noindent\begin{minipage}[t]{0.487\linewidth}
    \begin{tcolorbox}[enhanced,colback=popgreen,colframe=popink,boxrule=2.2pt,arc=10pt,
        drop shadow=popink,left=11pt,right=11pt,top=10pt,bottom=10pt,height=3.1cm,valign=center]
      \tcbox[colback=white,colframe=popink,boxrule=1.3pt,arc=5pt,boxsep=0pt,nobeforeafter,
        left=3pt,right=3pt,top=3pt,bottom=3pt,valign=center]{\qrcode[height=1.9cm]{https://play.google.com/store/apps/details?id=com.luvvix.onbuch&hl=fr}}%
      \hspace{8pt}\begin{minipage}[c]{0.5\linewidth}
        {\popdisplay\fontsize{11}{12.5}\selectfont\color{white}\MakeUppercase{Télécharge OnBuch}}\par\vspace{4pt}
        {\raggedright\popsemi\fontsize{8.4}{11}\selectfont\color{white}Gratuit \textbullet\ Google Play\\Tuteur IA, annales, concours, résultats\par}
      \end{minipage}
    \end{tcolorbox}
  \end{minipage}\hfill
  \begin{minipage}[t]{0.487\linewidth}
    \begin{tcolorbox}[enhanced,colback=poppurple,colframe=popink,boxrule=2.2pt,arc=10pt,
        drop shadow=popink,left=11pt,right=11pt,top=10pt,bottom=10pt,height=3.1cm,valign=center]
      \tikz[baseline=-0.7ex]\node[circle,fill=white,draw=popink,line width=1.3pt,minimum size=1.6cm,inner sep=0]{\popdisplay\fontsize{24}{24}\selectfont\color{poppurple}+};%
      \hspace{8pt}\begin{minipage}[c]{0.6\linewidth}
        {\popdisplay\fontsize{11}{12.5}\selectfont\color{white}\MakeUppercase{Passe à OnBuch+}}\par\vspace{4pt}
        {\raggedright\popsemi\fontsize{8.4}{11}\selectfont\color{white}Tous les cours signés, Tuteur IA illimité, fiches PDF — onglet OnBuch+ de l'app\par}
      \end{minipage}
    \end{tcolorbox}
  \end{minipage}
  \vspace{8pt}
  \popcontenu
  \vfill
  \signatureonbuch
  \restoregeometry
  \newpage
}

% Rangée « Ce cours contient » (page de garde)
\newcommand{\poptile}[3]{% #1 couleur #2 gros texte #3 légende
  \begin{tcolorbox}[enhanced,colback=#1,colframe=popink,boxrule=2pt,arc=9pt,shadow={2.5pt}{-2.5pt}{0pt}{popink},
      left=6pt,right=6pt,top=9pt,bottom=9pt,halign=center,valign=center,height=1.75cm,width=\linewidth,nobeforeafter]
    {\popdisplay\fontsize{12.5}{13}\selectfont\color{white}\MakeUppercase{#2}}\par\vspace{3pt}
    {\centering\popsemi\fontsize{7.8}{9.5}\selectfont\color{white}#3\par}
  \end{tcolorbox}}
\newcommand{\popcontenu}{%
  \par\noindent{\popxboldls\fontsize{7.5}{7.5}\selectfont\color{popmuted}CE COURS SIGNÉ CONTIENT}\par\vspace{7pt}
  \noindent\begin{minipage}[t]{0.235\linewidth}\poptile{popblue}{Cours}{complet, illustré, pas à pas}\end{minipage}\hfill
  \begin{minipage}[t]{0.235\linewidth}\poptile{poporange}{Méthodes}{et exercices résolus}\end{minipage}\hfill
  \begin{minipage}[t]{0.235\linewidth}\poptile{poppink}{Exercices}{type examen + corrigés}\end{minipage}\hfill
  \begin{minipage}[t]{0.235\linewidth}\poptile{popgreen}{Fiche bilan}{l'essentiel à retenir}\end{minipage}\par}

% Sommaire
\newtcolorbox{sommaire}{enhanced,colback=white,colframe=popink,boxrule=2.2pt,arc=10pt,
  drop shadow=popink,left=18pt,right=18pt,top=13pt,bottom=13pt,before skip=6pt,after skip=20pt,
  before upper={\poppill{Sommaire}{poppurple}{white}\par\vspace{9pt}\popsemi\fontsize{10.6}{14.5}\selectfont\color{popink}}}

% Signature de fin de cours — poussée tout en bas de la dernière page
\newcommand{\findecours}{%
  \par\vfill\nopagebreak
  \begin{center}\popsquiggle{poporange}\end{center}\vspace{4pt}
  \signatureonbuch
}
\AtBeginDocument{\def\llbracket{\mathopen{[\mkern-3.5mu[}}\def\rrbracket{\mathclose{]\mkern-3.5mu]}}} % intervalles d'entiers (glyphe ⟦ absent de la police)
% Glyphes absents de Fira Math : redéfinis à partir de glyphes disponibles
\AtBeginDocument{%
  \def\setminus{\mathbin{\backslash}}\let\smallsetminus\setminus
  \def\vdots{\mathord{\vbox{\baselineskip3.2pt\lineskiplimit0pt\kern4pt\hbox{.}\hbox{.}\hbox{.}}}}%
  \def\ddots{\mathinner{\mkern1mu\raise6.5pt\hbox{.}\mkern2mu\raise3.6pt\hbox{.}\mkern2mu\raise0.7pt\hbox{.}\mkern1mu}}%
  \def\triangle{\mathord{\scalebox{0.9}{$\symup{\Delta}$}}}%
  \def\nabla{\mathord{\raisebox{\depth}{\scalebox{1}[-1]{$\symup{\Delta}$}}}}%
  \def\checkmark{\popcheck{popdark}}%
  \def\star{\mathbin{\ast}}\def\bigstar{\mathord{\ast}}%
  \def\diamond{\mathbin{\scalebox{0.6}{$\Diamond$}}}%
  \def\bigcup{\mathop{\mathchoice{\raisebox{-0.3em}{\scalebox{1.7}{$\cup$}}}{\scalebox{1.25}{$\cup$}}{\cup}{\cup}}\displaylimits}%
  \def\bigcap{\mathop{\mathchoice{\raisebox{-0.3em}{\scalebox{1.7}{$\cap$}}}{\scalebox{1.25}{$\cap$}}{\cap}{\cap}}\displaylimits}%
  \def\lesseqqgtr{\mathrel{\lessgtr}}\def\bigoplus{\mathop{\oplus}\displaylimits}%
}
\providecommand{\ssi}{si et seulement si} % abréviation fréquente
\AtBeginDocument{\providecommand{\wideparen}{\overparen}} % arc de cercle

\begin{document}

\pagegarde{Logique, argumentation et démonstration}{Philosophie}{Tle Tech}{Philosophie – classe de Terminale technique}{N13}{5 séances (fiche de progression harmonisée)}{Cette leçon initie les élèves de Terminale technique aux opérations logiques fondamentales, à l'art de l'argumentation et à la démonstration. Elle leur fournit les instruments pour construire des raisonnements valides, déjouer les sophismes et rédiger une démarche justifiée, compétence clé au Baccalauréat technique.
\vspace{4pt}
\begin{multicols}{2}\small
\begin{itemize}
\item À la fin de cette leçon, je sais définir la logique et distinguer la logique formelle de la logique matérielle
\item À la fin de cette leçon, je sais identifier et utiliser les opérations logiques fondamentales (négation, conjonction, disjonction, implication, équivalence)
\item À la fin de cette leçon, je sais construire et analyser un syllogisme et vérifier sa validité
\item À la fin de cette leçon, je sais distinguer argumentation, persuasion et démonstration
\item À la fin de cette leçon, je sais reconnaître les principaux types d'arguments (autorité, analogie, exemple, cause à effet)
\item À la fin de cette leçon, je sais repérer et dénoncer les sophismes et paralogismes dans un discours
\end{itemize}\end{multicols}}

\begin{sommaire}
\begin{multicols}{2}
\begin{enumerate}
\item La logique : nature et enjeux
\item Les opérations logiques fondamentales
\item Le raisonnement déductif et le syllogisme
\item L'argumentation : convaincre et persuader
\item La démonstration
\item Applications : rédiger et justifier une démarche
\item Activité d'intégration
\item Méthodes et exercices résolus
\item Exercices
\item Corrigés des exercices
\item Fiche bilan
\end{enumerate}
\end{multicols}
\end{sommaire}

\begin{objectifs}
\begin{itemize}
\item À la fin de cette leçon, je sais définir la logique et distinguer la logique formelle de la logique matérielle
\item À la fin de cette leçon, je sais identifier et utiliser les opérations logiques fondamentales (négation, conjonction, disjonction, implication, équivalence)
\item À la fin de cette leçon, je sais construire et analyser un syllogisme et vérifier sa validité
\item À la fin de cette leçon, je sais distinguer argumentation, persuasion et démonstration
\item À la fin de cette leçon, je sais reconnaître les principaux types d'arguments (autorité, analogie, exemple, cause à effet)
\item À la fin de cette leçon, je sais repérer et dénoncer les sophismes et paralogismes dans un discours
\item À la fin de cette leçon, je sais construire une démonstration directe, par l'absurde ou par contraposée
\item À la fin de cette leçon, je sais appliquer ces notions à des situations concrètes de la vie courante camerounaise
\item À la fin de cette leçon, je sais rédiger et justifier une démarche, une réponse ou une conclusion dans une dissertation philosophique
\end{itemize}
\end{objectifs}

\begin{prerequis}
\begin{itemize}
\item Notion de vérité et d'opinion (doxa) vue en classe de Première
\item Notions de concept, jugement et raisonnement (introduction à la philosophie)
\item Capacité à analyser un texte argumentatif (français, classes antérieures)
\item Notion de discours et de rhétorique (Première)
\end{itemize}
\end{prerequis}

\newpage

\coursec{La logique : nature et enjeux}

\courssub{Situation d'accroche et problématique}

Lors d'une réunion de famille à Yaoundé, un oncle affirme : « Tous les jeunes qui réussissent au Bac ont travaillé ; or mon fils a travaillé ; donc il réussira. » Tout le monde acquiesce, mais un neveu élève en Terminale objecte que le raisonnement est faux. Qui a raison ? Comment distinguer un argument solide d'un discours simplement persuasif ? Dans les débats politiques, les marchés de Mokolo ou les réseaux sociaux, savoir raisonner juste est une arme de citoyen. Cette leçon nous donne les outils de la logique pour argumenter et démontrer rigoureusement.

La problématique de cette première partie est la suivante : qu'est-ce que la logique, que juge-t-elle exactement dans un raisonnement, et pourquoi est-il indispensable de distinguer la \emph{validité} d'un raisonnement de la \emph{vérité} de ce qu'il affirme ?

\courssub{Définition et étymologie de la logique}

Commençons par une intuition simple. Quand vous assemblez les pièces d'un moteur d'engin agricole, l'ordre des opérations compte : si vous montez les pièces dans le désordre, la machine ne fonctionne pas, même si chaque pièce est en parfait état. Il en va de même de la pensée : nos idées peuvent être « bonnes » une à une, mais si elles sont mal enchaînées, le raisonnement s'effondre. La logique est précisément la discipline qui étudie les règles de cet enchaînement.

Le mot « logique » vient du grec \emph{logos}, terme riche qui signifie à la fois la \cle{parole}, le \cle{discours}, la \cle{raison} et le \cle{calcul}. La logique est donc, à l'origine, la science de ce que la raison dit quand elle parle correctement.

\begin{definition}[La logique]
La \cle{logique} est la science des lois formelles de la pensée correcte, c'est-à-dire l'ensemble des règles qui garantissent qu'un raisonnement est bien construit et que sa conclusion découle nécessairement de ses prémisses.
\end{definition}

Deux précisions sont essentielles dans cette définition.

\begin{itemize}
\item La logique est une science des \textbf{lois} : elle ne décrit pas comment les gens pensent effectivement (c'est le travail de la psychologie), mais comment ils \emph{devraient} penser pour raisonner juste. Elle est une discipline \emph{normative} : elle pose des normes, comme le code de la route pose des règles de conduite.
\item Ces lois sont \textbf{formelles} : elles concernent la \emph{forme} du raisonnement, sa structure, indépendamment du sujet dont on parle. Qu'on raisonne sur les lions, les prix du manioc ou les élections, les mêmes règles d'enchaînement s'appliquent.
\end{itemize}

\begin{savaistu}[Aristote, père de la logique]
Aristote (384--322 av. J.-C.), élève de Platon et maître d'Alexandre le Grand, est considéré comme le père de la logique formelle. Dans un ensemble de traités réunis sous le titre d'\emph{Organon} (« l'instrument »), il fut le premier à analyser les formes du raisonnement, notamment le \emph{syllogisme}. Pendant plus de deux mille ans, sa logique est restée la référence en Europe et dans le monde arabe (Averroès, Avicenne). Au XIX\textsuperscript{e}--XX\textsuperscript{e} siècle, une \emph{logique moderne symbolique} (Boole, Frege, Russell) a traduit les raisonnements en langage mathématique, ouvrant la voie à l'informatique : chaque ordinateur applique, dans ses circuits, les lois de la logique.
\end{savaistu}

\courssub{Logique formelle et logique matérielle}

Pour bien juger un raisonnement, il faut distinguer deux questions différentes : le raisonnement est-il \emph{bien construit} ? Et ce qu'il affirme est-il \emph{vrai} ? Ces deux questions relèvent de deux branches de la logique.

\begin{definition}[Logique formelle et logique matérielle]
La \cle{logique formelle} étudie la validité de la \emph{structure} du raisonnement : la conclusion découle-t-elle nécessairement des prémisses, quel que soit le contenu ? La \cle{logique matérielle} s'intéresse à la \emph{vérité du contenu} : les prémisses et la conclusion sont-elles conformes à la réalité ?
\end{definition}

Prenons une image : la logique formelle ressemble au contrôle technique d'un véhicule. Le contrôleur vérifie que le moteur, les freins et la direction fonctionnent correctement ; il ne vérifie pas si le chauffeur se rend à la bonne destination. De même, la logique formelle vérifie que la « mécanique » du raisonnement fonctionne, sans se prononcer sur la vérité des affirmations de départ.

\begin{exemplebox}[Forme et contenu]
Considérons la forme suivante : « Tous les A sont des B ; or x est un A ; donc x est un B. » Cette forme est toujours valide, quel que soit le contenu :
\begin{itemize}
\item Contenu vrai : « Tous les mammifères ont un cœur ; or la chèvre est un mammifère ; donc la chèvre a un cœur. »
\item Contenu absurde : « Tous les élèves de Terminale sont des ministres ; or Amina est élève de Terminale ; donc Amina est ministre. »
\end{itemize}
Dans les deux cas, la \emph{forme} est correcte ; seul le contenu change. C'est précisément ce que la logique formelle étudie.
\end{exemplebox}

\courssub{Les trois opérations de l'esprit}

D'où viennent les éléments que la logique met en ordre ? La tradition philosophique, depuis Aristote, distingue trois opérations fondamentales de l'esprit, qui s'enchaînent comme les étapes d'une construction.

\textbf{Première opération : concevoir.} Concevoir, c'est former des \cle{concepts}, c'est-à-dire des idées générales qui saisissent l'essence d'une chose. Quand je vois un baobab, un palmier, un eucalyptus, mon esprit forme le concept d'« arbre », qui retient ce que ces réalités ont de commun. Le concept s'exprime par un mot ou un terme. Il a une \emph{compréhension} (les caractères qu'il contient : être végétal, avoir un tronc, des racines...) et une \emph{extension} (l'ensemble des êtres auxquels il s'applique : tous les arbres du monde).

\textbf{Deuxième opération : juger.} Juger, c'est unir ou séparer deux concepts en affirmant ou en niant quelque chose : « le baobab est un arbre », « le camerounais moyen n'est pas riche ». Le jugement s'exprime par une \cle{proposition} (une phrase qui affirme ou nie). C'est au niveau du jugement que naît la question de la \emph{vérité} : un jugement est vrai s'il correspond à la réalité, faux sinon.

\textbf{Troisième opération : raisonner.} Raisonner, c'est tirer d'un ou plusieurs jugements un jugement nouveau qui en découle : « Tous les hommes sont mortels ; or Socrate est un homme ; donc Socrate est mortel. » Le raisonnement s'exprime par un enchaînement de propositions : les \emph{prémisses} (ce qu'on donne au départ) et la \emph{conclusion} (ce qu'on en tire). C'est au niveau du raisonnement que naît la question de la \emph{validité}.

\begin{popfigure}
\begin{center}
\begin{tikzpicture}[
  node distance=0.9cm and 1.6cm,
  op/.style={rectangle, rounded corners, draw=popblue, fill=popblueL, thick, minimum width=3.1cm, minimum height=0.9cm, align=center, font=\bfseries},
  prod/.style={rectangle, rounded corners, draw=poporange, fill=poporangeL, thick, minimum width=3.4cm, minimum height=0.9cm, align=center},
  fleche/.style={-{Stealth[length=3mm]}, thick, popdark}
]
\node[op] (c) {Concevoir};
\node[prod, right=of c, align=center] (cp){Le \textbf{concept}\\ (le terme : « arbre »)};
\node[op, below=of c] (j) {Juger};
\node[prod, right=of j, align=center] (jp){Le \textbf{jugement}\\ (la proposition : « le baobab est un arbre »)};
\node[op, below=of j] (r) {Raisonner};
\node[prod, right=of r, align=center] (rp){Le \textbf{raisonnement}\\ (prémisses $\rightarrow$ conclusion)};
\draw[fleche] (c) -- (cp);
\draw[fleche] (j) -- (jp);
\draw[fleche] (r) -- (rp);
\draw[fleche, popgreen] (c) -- node[left, font=\itshape]{fournit les concepts} (j);
\draw[fleche, popgreen] (j) -- node[left, font=\itshape]{fournit les jugements} (r);
\end{tikzpicture}
\end{center}
\legende{Les trois opérations de l'esprit et leurs produits. Chaque opération s'appuie sur la précédente : on ne peut juger sans concepts, ni raisonner sans jugements.}
\end{popfigure}

\begin{aretenir}[Concept, jugement, raisonnement]
\begin{itemize}
\item \textbf{Concevoir} produit le \cle{concept}, exprimé par le \emph{terme}.
\item \textbf{Juger} produit le \cle{jugement}, exprimé par la \emph{proposition} (vraie ou fausse).
\item \textbf{Raisonner} produit le \cle{raisonnement}, exprimé par l'enchaînement prémisses--conclusion (valide ou invalide).
\end{itemize}
La logique s'occupe des trois opérations, mais son cœur est le raisonnement.
\end{aretenir}

\courssub{Validité et vérité : la distinction capitale}

Voici le point le plus important de cette section, et le plus délicat pour les candidats au Baccalauréat. Revenons à l'intuition : un raisonnement est comme un tuyau. Si le tuyau est bien raccordé, tout ce qu'on y verse arrive à l'autre bout ; si on y verse de l'eau propre, on recueille de l'eau propre ; si on y verse de l'eau sale, on recueille de l'eau sale — mais le tuyau, lui, fonctionne parfaitement. La \emph{validité} concerne le tuyau ; la \emph{vérité} concerne l'eau.

\begin{definition}[Validité]
La \cle{validité} est la propriété d'un raisonnement dont la conclusion découle nécessairement des prémisses : si les prémisses sont supposées vraies, la conclusion ne peut pas être fausse. Un raisonnement qui ne respecte pas cette condition est dit \emph{invalide}.
\end{definition}

Il faut donc retenir ceci : la vérité qualifie les \emph{propositions} (les jugements), tandis que la validité qualifie les \emph{raisonnements} (les enchaînements). Dire d'un raisonnement qu'il est « vrai » ou d'une proposition qu'elle est « valide » est un abus de langage à éviter dans une copie.

\begin{attention}[Ne jamais confondre « vrai » et « valide »]
Erreur fréquente au Bac : croire qu'un raisonnement valide garantit une conclusion vraie. C'est faux ! Un raisonnement peut être \textbf{valide avec des prémisses fausses} : la conclusion sera alors fausse, mais la faute vient des prémisses, non de la structure. Inversement, un raisonnement peut aboutir à une conclusion vraie tout en étant \textbf{invalide} : la conclusion est vraie par hasard, non par la force de l'argument. En logique, on examine toujours séparément : (1) la forme est-elle valide ? (2) les prémisses sont-elles vraies ?
\end{attention}

\begin{exemplebox}[Un raisonnement valide à conclusion fausse]
« Tous les Camerounais sont des lions ; or Paul Biya est camerounais ; donc Paul Biya est un lion. »
\begin{itemize}
\item La \emph{forme} est valide : c'est le schéma correct « Tous les A sont B ; x est A ; donc x est B ». Si les prémisses étaient vraies, la conclusion le serait nécessairement.
\item Mais la première prémisse est \emph{fausse} (les Camerounais sont des hommes, non des lions, même si l'équipe nationale porte fièrement le nom de Lions Indomptables !). La conclusion est donc fausse.
\end{itemize}
Le raisonnement est valide, mais non \emph{concluant} au sens fort : pour convaincre rationnellement, il faut une forme valide \textbf{et} des prémisses vraies.
\end{exemplebox}

Le tableau suivant croise les deux critères et donne un exemple de chaque cas.

\begin{center}
\begin{tabularx}{\linewidth}{|l|X|X|}
\hline
\rowcolor{popblueL} \textbf{Cas} & \textbf{Raisonnement VALIDE} & \textbf{Raisonnement INVALIDE} \\
\hline
\textbf{Prémisses vraies} & « Tous les élèves de la classe sont inscrits à l'examen ; or Ngo Bassa est élève de la classe ; donc elle est inscrite. » Conclusion vraie, garantie. & « S'il pleut, la route est glissante ; or la route est glissante ; donc il pleut. » Faux raisonnement : la route peut être glissante à cause d'une nappe d'huile. \\
\hline
\textbf{Prémisses fausses} & « Tous les Camerounais sont des lions ; or Paul Biya est camerounais ; donc Paul Biya est un lion. » Forme correcte, conclusion fausse. & « Tous les riches sont heureux ; or mon voisin est heureux ; donc il est riche. » Prémisse douteuse \emph{et} forme fautive : rien ne va. \\
\hline
\end{tabularx}
\end{center}

\begin{exoresolu}[Analyser le raisonnement de l'oncle]
Énoncé. Reprenons la situation d'accroche : « Tous les jeunes qui réussissent au Bac ont travaillé ; or mon fils a travaillé ; donc il réussira. » Le neveu affirme que ce raisonnement est faux. Analysez ce raisonnement en distinguant validité et vérité, puis tranchez le débat familial.
\tcblower
Solution détaillée. Traduisons le raisonnement en forme abstraite : « Tous les A (ceux qui réussissent) sont B (ceux qui ont travaillé) ; or x (le fils) est B ; donc x est A. »
\begin{enumerate}
\item \textbf{Examen de la forme.} La prémisse dit que tous les A sont B, c'est-à-dire que le travail est une condition \emph{nécessaire} de la réussite. Elle ne dit pas que tous les B sont A, c'est-à-dire que le travail est une condition \emph{suffisante}. Or le raisonnement conclut comme si tout travailleur réussissait automatiquement. C'est l'erreur classique dite de \emph{conversion illégale} : on a retourné « tous les A sont B » en « tous les B sont A ». La forme est donc \textbf{invalide}.
\item \textbf{Examen du contenu.} La première prémisse est plausible (la réussite sans travail est rarissime), et la seconde est supposée vraie (le fils a travaillé). Mais même avec des prémisses vraies, la conclusion ne s'impose pas : le fils peut travailler et échouer (maladie le jour de l'épreuve, mauvaise méthode, sujet mal traité).
\item \textbf{Conclusion.} Le neveu élève de Terminale a raison : le raisonnement de l'oncle est invalide, car il confond condition nécessaire et condition suffisante. Le travail rend la réussite possible et probable, mais ne la garantit pas logiquement. Cet exemple montre l'utilité pratique de la logique : elle protège des faux espoirs... et des fausses promesses.
\end{enumerate}
\end{exoresolu}

\courssub{Repères historiques : d'Aristote à la logique moderne}

La logique naît en Grèce, au IV\textsuperscript{e} siècle av. J.-C., dans un contexte précis : celui des débats de l'agora, où les sophistes enseignaient l'art de persuader par tous les moyens, y compris par des raisonnements trompeurs (les \emph{sophismes}). Contre eux, Platon puis surtout Aristote cherchent à fixer les règles du discours vrai. Dans l'\emph{Organon}, Aristote définit le \cle{syllogisme} : « un discours dans lequel, certaines choses étant posées, autre chose que ce qui a été posé s'ensuit nécessairement ». C'est la première théorie de la validité formelle de l'histoire.

La \emph{logique classique} ainsi fondée est transmise et enrichie par les philosophes arabes (Al-Fârâbî, Avicenne, Averroès), puis par la scolastique médiévale (saint Thomas d'Aquin) et les logiciens de Port-Royal au XVII\textsuperscript{e} siècle, qui rattachent la logique à une « art de penser ».

Au XIX\textsuperscript{e} siècle, une révolution s'accomplit : George Boole (1815--1864) montre que les raisonnements peuvent se traduire en calcul algébrique ; Gottlob Frege (1848--1925) fonde la \cle{logique symbolique} (ou logique mathématique), qui remplace le langage ordinaire par des signes : le connecteur « et », « ou », « si... alors », les quantificateurs « pour tout x », « il existe un x ». Cette logique moderne est devenue le fondement des mathématiques contemporaines et de l'informatique : les circuits électroniques des ordinateurs et des téléphones portables appliquent littéralement l'algèbre de Boole. Au Bac technique, il suffit de savoir que cette logique symbolique existe et qu'elle prolonge, en les formalisant, les intuitions d'Aristote.

\courssub{L'intérêt pratique de la logique}

Pourquoi un élève de Terminale technique, futur mécanicien, électrotechnicien, comptable, infirmier ou informaticien, devrait-il étudier la logique ? Pour au moins quatre raisons.

\begin{itemize}
\item \textbf{Dans les débats publics.} En période électorale, les discours abondent en sophismes : « Ce candidat a construit une école, donc il est honnête », « Vous critiquez le gouvernement, donc vous êtes un ennemi du pays ». La logique permet de démonter ces raisonnements et de voter en citoyen éclairé.
\item \textbf{Aux examens.} Au Probatoire et au Baccalauréat technique, la dissertation exige une démonstration rigoureuse : des prémisses claires, des enchaînements valides, une conclusion qui découle. Un correcteur sanctionne les conclusions qui « tombent du ciel ».
\item \textbf{Dans la vie professionnelle.} Le diagnostic d'une panne est un raisonnement logique : « Si le fusible est bon et que le courant arrive, alors la panne vient du moteur. » Le technicien qui raisonne mal change des pièces au hasard et fait perdre du temps et de l'argent à son atelier.
\item \textbf{Dans la vie quotidienne.} Au marché de Mokolo, sur WhatsApp ou Facebook, celui qui sait distinguer le vrai du valide ne se laisse ni escroquer ni manipuler par les rumeurs.
\end{itemize}

\begin{methode}[Analyser un raisonnement en quatre étapes]
Face à n'importe quel argument (dans un sujet de dissertation, un discours, une publicité) :
\begin{enumerate}
\item \textbf{Dégager la structure} : identifier les prémisses et la conclusion (repères : « donc », « par conséquent », « c'est pourquoi »).
\item \textbf{Traduire en forme abstraite} : remplacer les contenus par des lettres (A, B, x) pour voir la structure.
\item \textbf{Tester la validité} : la conclusion découle-t-elle nécessairement ? Chercher un contre-exemple où les prémisses seraient vraies et la conclusion fausse.
\item \textbf{Vérifier la vérité des prémisses} : sont-elles conformes aux faits, prouvables, admissibles ?
\end{enumerate}
Un argument n'est pleinement convaincant que s'il passe les étapes 3 et 4.
\end{methode}

\begin{exercice}[À vous de raisonner]
Analysez les trois raisonnements suivants en précisant, pour chacun, s'il est valide ou invalide, et si ses prémisses sont vraies ou fausses :
\begin{enumerate}
\item « Tous les métaux conduisent l'électricité ; or le cuivre est un métal ; donc le cuivre conduit l'électricité. »
\item « Tous les footballeurs professionnels sont bien payés ; or Samuel Eto'o est bien payé ; donc Samuel Eto'o est footballeur professionnel. »
\item « Tous les ordinateurs sont des êtres vivants ; or cet appareil est un ordinateur ; donc cet appareil est un être vivant. »
\end{enumerate}
\end{exercice}

\begin{corrige}[Exercice : analyser trois raisonnements]
\begin{enumerate}
\item Forme « Tous les A sont B ; x est A ; donc x est B » : \textbf{valide}. Les prémisses sont \textbf{vraies} (le cuivre est bien un métal conducteur). Raisonnement valide et conclusion vraie : argument parfaitement concluant.
\item Forme « Tous les A sont B ; x est B ; donc x est A » : \textbf{invalide} (conversion illégale), même si la conclusion est vraie dans les faits. On peut être bien payé sans être footballeur (un grand chirurgien, un chef d'entreprise). La conclusion est vraie par hasard, non par la logique.
\item Forme « Tous les A sont B ; x est A ; donc x est B » : \textbf{valide}. Mais la première prémisse est \textbf{fausse} : la conclusion est donc fausse. C'est le cas symétrique de l'exemple des lions : la structure fonctionne, c'est le contenu qui trompe.
\end{enumerate}
\end{corrige}

\begin{aretenir}[L'essentiel de la section]
\begin{itemize}
\item La \cle{logique} (du grec \emph{logos}) est la science des lois formelles de la pensée correcte ; son fondateur est Aristote (\emph{Organon}, IV\textsuperscript{e} s. av. J.-C.).
\item La \cle{logique formelle} juge la structure du raisonnement ; la \cle{logique matérielle} juge la vérité du contenu.
\item L'esprit accomplit trois opérations : \emph{concevoir} (concept), \emph{juger} (jugement, vrai ou faux), \emph{raisonner} (raisonnement, valide ou invalide).
\item La \cle{validité} est la propriété d'un raisonnement dont la conclusion découle nécessairement des prémisses ; elle est indépendante de la vérité des prémisses.
\item Un bon argument exige les deux : une forme \textbf{valide} et des prémisses \textbf{vraies}.
\end{itemize}
\end{aretenir}

La suite de la leçon examinera en détail les opérations logiques fondamentales (négation, conjonction, disjonction, implication), puis le raisonnement déductif et le syllogisme, avant de distinguer argumentation et démonstration.

\coursec{Les opérations logiques fondamentales}

Après avoir précisé la nature de la logique et ses enjeux pour la pensée rigoureuse, nous entrons maintenant dans le « moteur » même du raisonnement : les \textbf{opérations logiques fondamentales}. Tout comme l'arithmétique repose sur l'addition, la soustraction, la multiplication et la division, la logique propositionnelle repose sur quelques connecteurs qui permettent de construire des énoncés complexes à partir d'énoncés simples. Maîtriser ces opérations, c'est s'outiller pour analyser la structure d'un argument, détecter les sophismes et construire des démonstrations valides.

\courssub{Les propositions et les connecteurs logiques}

Une \cle{proposition} (ou énoncé) est une phrase qui a un sens complet et qui peut être déclarée \textbf{vraie} (V) ou \textbf{fausse} (F), mais pas les deux à la fois. C'est le \emph{principe de bivalence}.
\begin{itemize}
    \item \textit{Exemples :} « Yaoundé est la capitale politique du Cameroun » (V). « 2 + 2 = 5 » (F).
    \item \textit{Non-propositions :} « Quelle heure est-il ? » (question), « Fermez la porte ! » (injonction), « Je souhaite réussir » (souhait).
\end{itemize}

Les \cle{connecteurs logiques} (ou opérateurs) permettent de lier des propositions entre elles pour former des \textbf{propositions composées}. La valeur de vérité de la proposition composée dépend \emph{uniquement} de la valeur de vérité de ses composantes et de la règle propre au connecteur utilisé. Cette règle est synthétisée dans une \cle{table de vérité}.

\begin{methode}[Construire une table de vérité]
\textbf{Étape 1 :} Identifier les propositions élémentaires (p, q, r...). Pour $n$ propositions, il y a $2^n$ lignes possibles.
\textbf{Étape 2 :} Lister systématiquement toutes les combinaisons de valeurs de vérité (V/F) pour ces propositions. On procède souvent par ordre binaire décroissant (VV, VF, FV, FF pour deux propositions).
\textbf{Étape 3 :} Évaluer la proposition composée étape par étape, en ajoutant une colonne par sous-formule, en appliquant la règle du connecteur principal en dernier.
\textbf{Étape 4 :} Lire la colonne finale : elle donne la valeur de vérité de la formule pour chaque cas de figure.
\end{methode}

\courssub{La négation (non p) : $\neg p$ ou $\overline{p}$}

La négation est une opération \textbf{unaire} (elle porte sur une seule proposition). Elle inverse la valeur de vérité.

\begin{definition}[Négation]
La négation de la proposition $p$, notée $\neg p$ (se lit « non $p$ »), est la proposition qui est \textbf{vraie quand $p$ est fausse} et \textbf{fausse quand $p$ est vraie}.
\end{definition}

\begin{popfigure}
\begin{center}
\begin{tabular}{|c|c|}\hline
$p$ & $\neg p$ \\ \hline
V & F \\ \hline
F & V \\ \hline
\end{tabular}
\end{center}
\legende{Table de vérité de la négation.}
\end{popfigure}

\begin{exemplebox}[Négation dans le langage courant]
\begin{itemize}
    \item $p$ : « L'élève a obtenu le Probatoire. » $\rightarrow$ $\neg p$ : « L'élève n'a \textbf{pas} obtenu le Probatoire. »
    \item $p$ : « Tous les élèves sont présents. » $\rightarrow$ $\neg p$ : « \textbf{Pas tous} les élèves sont présents » (au moins un est absent). \textbf{Attention :} La négation de « Tous » n'est pas « Aucun », c'est « Pas tous » (existence d'au moins un contre-exemple).
\end{itemize}
\end{exemplebox}

\courssub{La conjonction (p et q) : $p \land q$}

La conjonction est une opération \textbf{binaire}. Elle exprime la simultanéité ou l'addition de deux faits.

\begin{definition}[Conjonction]
La conjonction de $p$ et $q$, notée $p \land q$ (se lit « $p$ \textbf{et} $q$ »), est \textbf{vraie seulement si les deux propositions sont vraies}. Elle est fausse dans tous les autres cas (au moins l'une est fausse).
\end{definition}

\begin{popfigure}
\begin{center}
\begin{tabular}{|c|c|c|}\hline
$p$ & $q$ & $p \land q$ \\ \hline
V & V & V \\ \hline
V & F & F \\ \hline
F & V & F \\ \hline
F & F & F \\ \hline
\end{tabular}
\end{center}
\legende{Table de vérité de la conjonction.}
\end{popfigure}

\begin{popfigure}
\begin{tikzpicture}[scale=1.2]
    % Ensemble U
    \draw (0,0) rectangle (5,3.5) node[above left] {$U$ : Élèves de la classe};
    % Ensemble P (Philo)
    \draw[fill=popblue!20, draw=popblue, thick] (1,1) circle (1.2) node[above left] {$P$ : Réussissent en Philo};
    % Ensemble M (Maths)
    \draw[fill=poporange!20, draw=poporange, thick] (3,1) circle (1.2) node[above right] {$M$ : Réussissent en Maths};
    % Intersection hachurée
    \begin{scope}
        \clip (1,1) circle (1.2);
        \clip (3,1) circle (1.2);
        \fill[pattern=north east lines, pattern color=popgreen] (1,1) circle (1.2);
    \end{scope}
    \node[popgreen, font=\bfseries] at (2,1) {$P \land M$};
    % Légende
    \node[align=left, anchor=north west] at (0,3.7) {
        \textbf{Diagramme de Venn : Conjonction}\\
        Zone hachurée = Élèves réussissant \textbf{à la fois} en Philo \textbf{ET} en Maths.
    };
\end{tikzpicture}
\legende{Illustration ensembliste de la conjonction (intersection).}
\end{popfigure}

\begin{exemplebox}[Conjonction dans un règlement scolaire]
« L'accès à la salle informatique est autorisé \textbf{si l'élève a une carte d'identité} \textbf{ET} \textbf{s'il a signé le règlement}. »
$p$ : « L'élève a une carte d'identité. »
$q$ : « L'élève a signé le règlement. »
Condition d'accès : $p \land q$. Si l'un des deux manque ($\neg p$ ou $\neg q$), l'accès est refusé ($p \land q$ est faux).
\end{exemplebox}

\courssub{La disjonction : « ou » inclusif ($p \lor q$) et « ou » exclusif ($p \veebar q$)}

En français courant, « ou » est ambigu. La logique le distingue en deux connecteurs formels.

\textbf{La disjonction inclusive (ou faible) : $p \lor q$}\par
C'est le « ou » mathématique et logique standard : « $p$ ou $q$ (ou les deux) ».

\begin{definition}[Disjonction inclusive]
$p \lor q$ est \textbf{fausse seulement si les deux propositions sont fausses}. Elle est vraie dans les trois autres cas.
\end{definition}

\begin{popfigure}
\begin{center}
\begin{tabular}{|c|c|c|}\hline
$p$ & $q$ & $p \lor q$ \\ \hline
V & V & V \\ \hline
V & F & V \\ \hline
F & V & V \\ \hline
F & F & F \\ \hline
\end{tabular}
\end{center}
\legende{Table de vérité de la disjonction inclusive.}
\end{popfigure}

\begin{popfigure}
\begin{tikzpicture}[scale=1.2]
    \draw (0,0) rectangle (5,3.5) node[above left] {$U$};
    % Cercles avec fond coloré
    \draw[fill=popblue!20, draw=popblue, thick] (1,1) circle (1.2) node[above left] {$P$};
    \draw[fill=poporange!20, draw=poporange, thick] (3,1) circle (1.2) node[above right] {$M$};
    % Union hachurée : on remplit les deux cercles par-dessus
    \fill[pattern=north west lines, pattern color=poppurple, opacity=0.6] (1,1) circle (1.2);
    \fill[pattern=north west lines, pattern color=poppurple, opacity=0.6] (3,1) circle (1.2);
    \node[poppurple, font=\bfseries] at (2,1) {$P \lor M$};
    \node[align=left, anchor=north west] at (0,3.7) {
        \textbf{Diagramme de Venn : Disjonction inclusive}\\
        Zone hachurée = Élèves réussissant en Philo \textbf{OU} en Maths \textbf{(ou les deux)}.
    };
\end{tikzpicture}
\legende{Illustration ensembliste de la disjonction inclusive (réunion).}
\end{popfigure}

\textbf{La disjonction exclusive (ou fort) : $p \veebar q$ ou $p \oplus q$}\par
C'est le « ou bien... ou bien... » du langage courant quand les alternatives s'excluent mutuellement : « $p$ ou $q$, mais \textbf{pas les deux} ».

\begin{definition}[Disjonction exclusive]
$p \veebar q$ est \textbf{vraie si l'une des deux propositions est vraie et l'autre fausse}. Elle est fausse si les deux sont vraies ou si les deux sont fausses.
\end{definition}

\begin{popfigure}
\begin{center}
\begin{tabular}{|c|c|c|}\hline
$p$ & $q$ & $p \veebar q$ \\ \hline
V & V & F \\ \hline
V & F & V \\ \hline
F & V & V \\ \hline
F & F & F \\ \hline
\end{tabular}
\end{center}
\legende{Table de vérité de la disjonction exclusive.}
\end{popfigure}

\begin{attention}[Piège : « Ou » inclusif vs exclusif dans les énoncés]
\begin{itemize}
    \item \textbf{Inclusif (souvent par défaut en logique/maths) :} « Vous pouvez payer par carte \textbf{ou} par mobile money. » (Les deux sont acceptés, voire cumulables).
    \item \textbf{Exclusif (souvent dans les choix de vie, menus, alternatives) :} « Pour le dessert, ce sera glace \textbf{ou bien} tarte. » (Il faut choisir, pas les deux).
    \item \textbf{Indice linguistique :} « Soit... soit... » marque souvent l'exclusif. « Et/ou » marque explicitement l'inclusif.
\end{itemize}
\end{attention}

\courssub{L'implication (si p alors q) : $p \Rightarrow q$}

C'est le connecteur le plus important pour le raisonnement (théorèmes, lois, règlements, promesses), mais aussi le plus contre-intuitif.

\begin{definition}[Implication / Conditionnelle]
L'implication $p \Rightarrow q$ (se lit « si $p$ alors $q$ », « $p$ implique $q$ ») est \textbf{fausse seulement dans un cas : lorsque l'antécédent $p$ est VRAI et le conséquent $q$ est FAUX}. Dans les trois autres cas, elle est VRAIE.
\end{definition}

\begin{popfigure}
\begin{center}
\begin{tabular}{|c|c|c|}\hline
$p$ (antécédent) & $q$ (conséquent) & $p \Rightarrow q$ \\ \hline
V & V & V \\ \hline
V & F & \textbf{F} \\ \hline
F & V & V \\ \hline
F & F & V \\ \hline
\end{tabular}
\end{center}
\legende{Table de vérité de l'implication. La seule case fausse est soulignée.}
\end{popfigure}

\begin{propriete}[Condition suffisante et condition nécessaire]
Dans l'implication $p \Rightarrow q$ :
\begin{itemize}
    \item $p$ est une \textbf{condition suffisante} de $q$ : la vérité de $p$ \emph{garantit} (suffit pour) la vérité de $q$.
    \item $q$ est une \textbf{condition nécessaire} à $p$ : on ne peut pas avoir $p$ vrai sans avoir $q$ vrai ($p$ vrai $\Rightarrow$ $q$ vrai).
\end{itemize}
\textbf{Attention :} $p \Rightarrow q$ n'est \textbf{PAS} équivalent à $q \Rightarrow p$ (réciproque). Confondre les deux est une erreur logique majeure (affirmation du conséquent).
\end{propriete}

\begin{exemplebox}[Implication dans la vie scolaire (Règlement / Contrat)]
Énoncé : \textbf{« Si un élève a une moyenne $\ge$ 10/20, alors il est admis au conseil de classe. »}
\begin{itemize}
    \item $p$ : « Moyenne $\ge$ 10 ». (Condition \textbf{suffisante} pour être admis).
    \item $q$ : « Admis au conseil ». (Condition \textbf{nécessaire} pour avoir eu $\ge$ 10).
\end{itemize}
Analyse des cas (Table de vérité) :
\begin{enumerate}
    \item Moyenne 12/20 \textbf{ET} Admis : $p$=V, $q$=V $\rightarrow$ Implication \textbf{VRAIE} (Règle respectée).
    \item Moyenne 12/20 \textbf{ET} Refusé : $p$=V, $q$=F $\rightarrow$ Implication \textbf{FAUSSE} (Règle violée ! Injustice).
    \item Moyenne 8/20 \textbf{ET} Admis : $p$=F, $q$=V $\rightarrow$ Implication \textbf{VRAIE}. \emph{Pourquoi ?} La règle ne dit \textbf{pas} « Seuls les $\ge$ 10 sont admis ». Elle fixe un plancher, pas un plafond. Un élève à 8/20 peut être admis pour d'autres motifs (avis favorable, dossier). La règle n'est pas contredite.
    \item Moyenne 8/20 \textbf{ET} Refusé : $p$=F, $q$=F $\rightarrow$ Implication \textbf{VRAIE}. Cas normal.
\end{enumerate}
\end{exemplebox}

\begin{attention}[Erreur fréquente : Confondre implication et équivalence]
\begin{itemize}
    \item Dire « Si j'ai 10, je passe » ($p \Rightarrow q$) n'interdit \textbf{pas} qu'un élève à 8 passe aussi.
    \item Dire « Je passe \textbf{si et seulement si} j'ai 10 » ($p \Leftrightarrow q$) \textbf{interdit} qu'un élève à 8 passe (et qu'un élève à 10 rate).
    \item \textbf{Piège rhétorique :} « Si tu m'aimes, tu feras ça. » $\not\Rightarrow$ « Si tu fais ça, c'est que tu m'aimes. » (Réciproque non garantie).
\end{itemize}
\end{attention}

\courssub{L'équivalence (p si et seulement si q) : $p \Leftrightarrow q$}

L'équivalence exprime une double implication : les deux propositions ont \emph{exactement} la même valeur de vérité. Elles sont « logiquement indiscernables ».

\begin{definition}[Équivalence]
$p \Leftrightarrow q$ (se lit « $p$ \textbf{si et seulement si} $q$ », ou « $p$ est équivalent à $q$ ») est \textbf{vraie si $p$ et $q$ ont la même valeur de vérité} (tous deux vrais ou tous deux faux). Elle est fausse si elles ont des valeurs différentes.
\end{definition}
C'est la conjonction de deux implications : $(p \Rightarrow q) \land (q \Rightarrow p)$.

\begin{popfigure}
\begin{center}
\begin{tabular}{|c|c|c|}\hline
$p$ & $q$ & $p \Leftrightarrow q$ \\ \hline
V & V & V \\ \hline
V & F & F \\ \hline
F & V & F \\ \hline
F & F & V \\ \hline
\end{tabular}
\end{center}
\legende{Table de vérité de l'équivalence.}
\end{popfigure}

\begin{exemplebox}[Équivalence dans une définition ou un contrat]
« Un triangle est rectangle \textbf{si et seulement si} le carré de son plus grand côté est égal à la somme des carrés des deux autres côtés. » (Théorème de Pythagore + Réciproque).
\begin{itemize}
    \item Si triangle rectangle $\Rightarrow$ relation vraie.
    \item Si relation vraie $\Rightarrow$ triangle rectangle.
    \item Les deux propositions sont \textbf{inséparables} : on ne peut pas avoir l'une sans l'autre.
\end{itemize}
\end{exemplebox}

\courssub{Synthèse : Tables de vérité des cinq connecteurs fondamentaux}

\begin{popfigure}
\begin{center}
\small
\begin{tabular}{|c|c|c|c|c|c|c|}\hline
$p$ & $q$ & $\neg p$ & $p \land q$ & $p \lor q$ & $p \Rightarrow q$ & $p \Leftrightarrow q$ \\ \hline
V & V & F & V & V & V & V \\ \hline
V & F & F & F & V & F & F \\ \hline
F & V & V & F & V & V & F \\ \hline
F & F & V & F & F & V & V \\ \hline
\end{tabular}
\end{center}
\legende{Table de vérité récapitulative des 5 connecteurs pour 2 propositions.}
\end{popfigure}

\courssub{Les trois grands principes logiques (Lois de la pensée)}

Au-delà des connecteurs, la logique classique repose sur trois principes axiomatisés par Aristote, fondements de toute argumentation rigoureuse.

\begin{propriete}[Principes aristotéliciens fondamentaux]
\begin{enumerate}
    \item \textbf{Principe d'identité :} $p \Rightarrow p$ (ou « A est A »). Une proposition est identique à elle-même. Ce qui est, est. Ce qui n'est pas, n'est pas. \emph{Exemple :} « Un élève admis est un élève admis. »
    \item \textbf{Principe de non-contradiction :} $\neg (p \land \neg p)$. Une proposition ne peut être \textbf{à la fois vraie et fausse} au même moment, sous le même rapport, dans le même sens. \emph{Exemple :} On ne peut pas dire simultanément « Paul est présent » et « Paul n'est pas présent » (au même cours, à la même heure).
    \item \textbf{Principe du tiers exclu :} $p \lor \neg p$. Entre une proposition et sa négation, il n'y a pas de troisième possibilité : l'une des deux est nécessairement vraie. \emph{Exemple :} « Soit Paul a réussi le Probatoire, soit il ne l'a pas réussi. » (Pas de « demi-réussite » logique).
\end{enumerate}
\end{propriete}

\begin{aretenir}[Pourquoi ces principes sont-ils cruciaux ?]
\begin{itemize}
    \item Sans \textbf{Identité} : les mots changent de sens au fil du discours (équivoque), le dialogue devient impossible.
    \item Sans \textbf{Non-contradiction} : on peut tout prouver et son contraire (explosion logique : \emph{ex contradictione quodlibet}). La pensée s'effondre.
    \item Sans \textbf{Tiers exclu} : on ne peut pas trancher par le raisonnement (dichotomie), on reste dans le flou, l'indécidable.
\end{itemize}
Ces principes ne sont pas « prouvés » (ils sont la condition de la preuve), mais on peut montrer qu'on ne peut pas les nier sans les utiliser (argument \emph{ad hominem} ou rétorsion).
\end{aretenir}

\courssub{Application : Formaliser des énoncés de la vie courante}

La formalisation consiste à traduire une phrase en langage naturel vers le langage symbolique de la logique propositionnelle. Cela permet de faire apparaître la structure logique, masquée par les mots.

\begin{methode}[Formaliser un énoncé]
\textbf{1. Repérer les propositions élémentaires} et leur donner une lettre (p, q, r...).
\textbf{2. Identifier les connecteurs} (mots-clés : \textit{et, ou, si... alors, si et seulement si, non, ni... ni...}).
\textbf{3. Attention aux pièges linguistiques} :
\begin{itemize}
    \item « Mais », « cependant », « pourtant » $\rightarrow$ souvent « et » (conjonction) avec nuance d'opposition.
    \item « Ni p ni q » $\rightarrow$ $\neg p \land \neg q$ (ni l'un ni l'autre = non p ET non q).
    \item « Soit p soit q » $\rightarrow$ souvent exclusif ($p \veebar q$), sauf contexte maths (inclusif).
    \item « p sauf si q » $\rightarrow$ $\neg q \Rightarrow p$ (ou $p \lor q$).
    \item « p seulement si q » $\rightarrow$ $p \Rightarrow q$ (q est nécessaire).
    \item « p si q » $\rightarrow$ $q \Rightarrow p$ (q est suffisant).
\end{itemize}
\textbf{4. Écrire la formule} en respectant la priorité des connecteurs (par défaut : $\neg$ > $\land$ > $\lor$ > $\Rightarrow$ > $\Leftrightarrow$). Utiliser des parenthèses pour lever toute ambiguïté.
\end{methode}

\begin{exoresolu}[Exercice 1 : Règlement scolaire]
\textbf{Énoncé :} « L'élève sera exclu \textbf{s'il} commet une faute grave \textbf{ou s'}il récidive après avertissement. »
\tcblower
\textbf{Solution :}
\begin{itemize}
    \item $p$ : « L'élève commet une faute grave. »
    \item $q$ : « L'élève récidive après avertissement. »
    \item $r$ : « L'élève sera exclu. »
    \item Structure : « Si ($p$ ou $q$) alors $r$ ». $\rightarrow$ \textbf{$(p \lor q) \Rightarrow r$}.
\end{itemize}
\end{exoresolu}

\begin{exoresolu}[Exercice 2 : Conditions d'obtention d'une bourse]
\textbf{Énoncé :} « Pour obtenir la bourse, il faut \textbf{avoir une moyenne $\ge$ 12/20} \textbf{ET} \textbf{être inscrit dans un établissement public}. »
\tcblower
\textbf{Solution :}
\begin{itemize}
    \item $p$ : « Moyenne $\ge$ 12/20. »
    \item $q$ : « Inscrit établissement public. »
    \item $r$ : « Obtient la bourse. »
    \item « Il faut p et q pour r » signifie : Si r alors (p et q). $\rightarrow$ \textbf{$r \Rightarrow (p \land q)$}.
    \item \emph{Note :} Cela ne dit pas que si on a p et q, on a la bourse (conditions nécessaires, pas forcément suffisantes s'il y a d'autres critères).
\end{itemize}
\end{exoresolu}

\begin{exoresolu}[Exercice 3 : Validité d'un contrat]
\textbf{Énoncé :} « Le contrat est valide \textbf{si et seulement si} les deux parties signent \textbf{ET} la loi ne l'interdit pas. »
\tcblower
\textbf{Solution :}
\begin{itemize}
    \item $p$ : « Contrat valide. »
    \item $q$ : « Les deux parties signent. »
    \item $r$ : « La loi l'interdit. » (Donc $\neg r$ : « La loi ne l'interdit pas »).
    \item $\rightarrow$ \textbf{$p \Leftrightarrow (q \land \neg r)$}.
\end{itemize}
\end{exoresolu}

\begin{exoresolu}[Exercice 4 : Piège « seulement si »]
\textbf{Énoncé :} « Tu réussiras \textbf{seulement si} tu travailles. »
\tcblower
\textbf{Solution :}
\begin{itemize}
    \item $p$ : « Tu réussis. »
    \item $q$ : « Tu travailles. »
    \item « Seulement si » introduit la condition \textbf{nécessaire}. $\rightarrow$ \textbf{$p \Rightarrow q$}.
    \item (Si tu réussis, c'est que tu as nécessairement travaillé. Travailler ne garantit pas la réussite, mais c'est indispensable).
\end{itemize}
\end{exoresolu}

\begin{exercice}[À vous de formaliser]
Formalisez les énoncés suivants en définissant vos lettres propositionnelles.
\begin{enumerate}
    \item « L'accès au laboratoire est interdit aux élèves qui n'ont pas de blouse \textbf{ou} qui n'ont pas de lunettes de protection. »
    \item « Le jury délibère \textbf{si et seulement si} tous les membres sont présents. »
    \item « Si le temps est au beau fixe, \textbf{alors} nous irons en excursion, \textbf{mais} si il pleut, nous resterons en classe. »
    \item « Vous ne pouvez pas avoir le permis \textbf{si} vous n'avez pas 18 ans. »
    \item « Soit le projet est accepté, soit il est refusé, \textbf{mais} il ne peut pas être les deux à la fois. »
\end{enumerate}
\end{exercice}

\begin{corrige}[Exercices de formalisation]
\begin{enumerate}
    \item $p$ : « A une blouse », $q$ : « A des lunettes », $r$ : « Accès autorisé ».
        « Interdit si (pas blouse OU pas lunettes) » $\equiv$ « Autorisé seulement si (blouse ET lunettes) ».
        $\rightarrow$ $r \Rightarrow (p \land q)$  (ou équivalent : $(\neg p \lor \neg q) \Rightarrow \neg r$).
    \item $p$ : « Jury délibère », $q$ : « Tous membres présents ».
        $\rightarrow$ $p \Leftrightarrow q$.
    \item $p$ : « Beau temps », $q$ : « Excursion », $r$ : « Pluie », $s$ : « Reste en classe ».
        $\rightarrow$ $(p \Rightarrow q) \land (r \Rightarrow s)$.
    \item $p$ : « A le permis », $q$ : « A 18 ans ».
        « Pas permis si pas 18 ans » $\equiv$ « Si pas 18 ans alors pas permis » $\equiv$ $\neg q \Rightarrow \neg p$.
        Par contraposition (équivalent) : $p \Rightarrow q$ (Avoir le permis $\Rightarrow$ Avoir 18 ans).
    \item $p$ : « Projet accepté », $q$ : « Projet refusé ».
        « Soit p soit q (exclusif) » ET « Pas (p et q) ».
        $\rightarrow$ $(p \veebar q)$  (la disjonction exclusive inclut déjà l'exclusion des deux vrais).
        Ou explicitement : $(p \lor q) \land \neg(p \land q)$.
\end{enumerate}
\end{corrige}

\begin{savaistu}[Logique et informatique : Les portes logiques]
Les opérations que nous venons d'étudier sont \textbf{physiquement réalisées} dans les processeurs de vos téléphones et ordinateurs par des \textbf{portes logiques} (transistors).
\begin{itemize}
    \item Porte \textbf{NOT} (Inverseur) = Négation.
    \item Porte \textbf{AND} = Conjonction (Sortie 1 si Entrée1=1 \textbf{ET} Entrée2=1).
    \item Porte \textbf{OR} = Disjonction inclusive.
    \item Porte \textbf{XOR} = Disjonction exclusive.
    \item L'implication $p \Rightarrow q$ se programme souvent par $\neg p \lor q$ (soit p est faux, soit q est vrai).
\end{itemize}
Tout algorithme, aussi complexe soit-il (IA, cryptage, jeu vidéo), se réduit in fine à des milliards de combinaisons de ces 5 opérations de base. La logique n'est pas une abstraction vaine : c'est l'architecture matérielle du numérique.
\end{savaistu}

\coursec{Le raisonnement déductif et le syllogisme}

Dans la section précédente, nous avons étudié les opérations logiques fondamentales : la négation, la conjonction, la disjonction et l'implication. Ces opérations sont comme les briques élémentaires du raisonnement. Mais un raisonnement n'est pas une simple juxtaposition de briques : c'est un édifice organisé, où l'on passe d'énoncés admis (les prémisses) à un énoncé nouveau (la conclusion). Comment ce passage peut-il être garanti ? C'est toute la question de la \cle{déduction}, que le philosophe grec Aristote (IV\textsuperscript{e} siècle av. J.-C.) a été le premier à analyser rigoureusement dans son \emph{Organon}, en inventant la théorie du \cle{syllogisme}. Cette section vous donne les outils pour construire des raisonnements valides et pour démasquer les faux raisonnements qui encombrent nos débats quotidiens.

\courssub{Le raisonnement déductif : du général au particulier}

\textbf{Intuition.} Imaginez un élève de Terminale qui sait que « tous les élèves inscrits au Baccalauréat technique ont passé le Probatoire ». S'il apprend que « Marie est inscrite au Baccalauréat technique », il conclut immédiatement : « donc Marie a passé le Probatoire ». Personne ne discute cette conclusion : elle s'impose avec la même nécessité qu'un résultat de calcul. Voilà l'essentiel de la déduction.

\begin{definition}[Raisonnement déductif]
Le raisonnement \cle{déductif} est un raisonnement qui va \textbf{du général au particulier} : à partir d'énoncés généraux admis comme vrais (les \cle{prémisses}), il tire une conclusion particulière qui s'impose avec \textbf{nécessité}. Si les prémisses sont vraies et si la forme du raisonnement est correcte, alors la conclusion est \emph{nécessairement} vraie : il est impossible qu'elle soit fausse.
\end{definition}

La propriété décisive de la déduction est donc la \cle{nécessité} de la conclusion. Contrairement à une simple opinion, une conclusion déductive ne laisse aucune place au doute, \emph{à condition} que deux exigences soient remplies : la vérité des prémisses et la validité de la forme.

\begin{attention}[Validité et vérité : ne pas confondre]
Un raisonnement peut être \textbf{valide} (correct dans sa forme) avec des prémisses fausses : « Tous les élèves de Terminale ont 40 ans ; or Jean est en Terminale ; donc Jean a 40 ans ». La forme est parfaite, la conclusion est fausse parce que la majeure est fausse. Inversement, un raisonnement peut aboutir à une conclusion vraie par une forme invalide. La logique étudie la \emph{validité formelle}, c'est-à-dire la correction de l'enchaînement, indépendamment du contenu.
\end{attention}

\courssub{La structure du syllogisme catégorique}

Le syllogisme est la forme la plus célèbre du raisonnement déductif.

\begin{definition}[Syllogisme]
Le \cle{syllogisme} est un raisonnement déductif composé de \textbf{deux prémisses} et d'\textbf{une conclusion}, reliant \textbf{trois termes} : le \cle{grand terme} (attribut de la conclusion), le \cle{petit terme} (sujet de la conclusion) et le \cle{moyen terme}, qui figure dans les deux prémisses mais \emph{disparaît} de la conclusion. La prémisse contenant le grand terme s'appelle la \cle{majeure} ; celle qui contient le petit terme s'appelle la \cle{mineure}.
\end{definition}

Prenons l'exemple canonique, légué par toute la tradition philosophique :

\begin{center}
\begin{tabularx}{\linewidth}{|l|X|}\hline
\rowcolor{popblueL} \textbf{Élément} & \textbf{Énoncé} \\ \hline
Majeure & Tous les hommes sont mortels. \\ \hline
Mineure & Or Socrate est un homme. \\ \hline
Conclusion & Donc Socrate est mortel. \\ \hline
\end{tabularx}
\end{center}

Analysons les trois termes : le \textbf{moyen terme} est « homme » (il sert d'intermédiaire et disparaît de la conclusion) ; le \textbf{grand terme} est « mortel » (attribut de la conclusion, présent dans la majeure) ; le \textbf{petit terme} est « Socrate » (sujet de la conclusion, présent dans la mineure). Le moyen terme fonctionne comme un pont : il relie le petit terme au grand terme, puis s'efface, comme l'échafaudage qu'on retire une fois la maison construite.

\begin{popfigure}
\begin{center}
\begin{tikzpicture}[box/.style={draw=popblue, thick, rounded corners, fill=popblueL, text width=3.6cm, align=center, minimum height=1.2cm}, conc/.style={draw=poporange, thick, rounded corners, fill=poporangeL, text width=4.2cm, align=center, minimum height=1.2cm}, arr/.style={-{Stealth[length=3mm]}, thick, popdark}]
\node[box, align=center] (maj) at (0,2.2){\textbf{Majeure}\\ Tous les hommes (M) sont mortels (G)};
\node[box, align=center] (min) at (5.4,2.2){\textbf{Mineure}\\ Socrate (P) est un homme (M)};
\node[conc, align=center] (ccl) at (2.7,0){\textbf{Conclusion}\\ Socrate (P) est mortel (G)};
\draw[arr] (maj.south) -- (ccl.north west);
\draw[arr] (min.south) -- (ccl.north east);
\node[popmuted, align=center] at (8.9,1.2) {M = moyen terme\\ G = grand terme\\ P = petit terme};
\end{tikzpicture}
\end{center}
\legende{Structure du syllogisme : la majeure et la mineure, unies par le moyen terme M, produisent nécessairement la conclusion.}
\end{popfigure}

\courssub{Vérifier la validité : les diagrammes de Venn}

Comment s'assurer qu'un syllogisme est valide ? Une méthode visuelle consiste à représenter les ensembles par des cercles, selon la technique des \cle{diagrammes de Venn} (logicien anglais du XIX\textsuperscript{e} siècle). La majeure « tous les hommes sont mortels » signifie que le cercle des hommes est entièrement inclus dans celui des mortels. La mineure « Socrate est un homme » place Socrate dans le cercle des hommes. On \emph{voit} alors que Socrate se trouve forcément dans le cercle des mortels : la conclusion est inévitable.

\begin{popfigure}
\begin{center}
\begin{tikzpicture}[scale=1]
\draw[thick, poporange, fill=poporangeL, opacity=0.85] (0,0) circle (2.6);
\draw[thick, popblue, fill=popblueL, opacity=0.85] (0,-0.2) circle (1.4);
\node[poporange, font=\bfseries] at (0,2.1) {Mortels (G)};
\node[popblue, font=\bfseries] at (0,0.85) {Hommes (M)};
\fill[popdark] (0,-0.35) circle (0.07);
\node[popdark, below] at (0,-0.45) {Socrate (P)};
\end{tikzpicture}
\end{center}
\legende{Diagramme de Venn du syllogisme de Socrate : le cercle des hommes est inclus dans celui des mortels ; le point « Socrate », placé dans le cercle des hommes, appartient nécessairement au cercle des mortels.}
\end{popfigure}

\begin{methode}[Vérifier un syllogisme en quatre étapes]
\begin{enumerate}
\item \textbf{Identifier la conclusion} : elle est introduite par « donc », « par conséquent », « c'est pourquoi ».
\item \textbf{Repérer les trois termes} : le petit terme (sujet de la conclusion), le grand terme (attribut de la conclusion) et le moyen terme (présent deux fois dans les prémisses).
\item \textbf{Vérifier que le moyen terme disparaît} de la conclusion : s'il y figure encore, le raisonnement est mal construit.
\item \textbf{Contrôler la structure} : le moyen terme doit être pris dans toute son extension au moins une fois (on dit qu'il doit être \emph{distribué}), et on ne peut rien affirmer dans la conclusion qui ne soit déjà contenu dans les prémisses (règles de non-distribution illicite du grand ou du petit terme). En cas de doute, dessiner les cercles de Venn.
\end{enumerate}
\end{methode}

La règle de la \cle{distribution du moyen terme} est la plus importante. Dire que le moyen terme doit être distribué au moins une fois signifie qu'il doit être pris, dans l'une des prémisses au moins, dans la \emph{totalité} de son extension (par exemple « \emph{tous} les hommes »). Sinon, rien ne garantit que le petit terme et le grand terme parlent des \emph{mêmes} éléments du moyen terme : le pont est rompu.

\courssub{Le paralogisme : l'apparence trompeuse de la rigueur}

\begin{definition}[Paralogisme]
Un \cle{paralogisme} est un raisonnement \textbf{faux} qui présente l'\textbf{apparence} de la rigueur logique. Il est généralement construit de bonne foi, par inadvertance, contrairement au \cle{sophisme} qui est un raisonnement fallacieux délibérément destiné à tromper.
\end{definition}

\begin{attention}[Erreur fréquente : l'affirmation du conséquent]
Le schéma « Tous les A sont B ; or C est B ; donc C est A » est \textbf{invalide}. Exemple : « Tous les fonctionnaires sont des employés ; or Paul est un employé ; donc Paul est fonctionnaire ». Faux : Paul peut être employé dans une entreprise privée. Le moyen terme « employé » n'est jamais pris dans toute son extension : le raisonnement confond « tous les A sont B » avec « tous les B sont A ». Sur un diagramme de Venn, le point C est dans le grand cercle B, mais rien ne le place dans le petit cercle A.
\end{attention}

\begin{exemplebox}[Le raisonnement de l'oncle passé au crible]
Reprenons le raisonnement entendu dans l'accroche de la leçon : « Tous les supporters des Lions Indomptables portent un maillot vert ; or mon neveu porte un maillot vert ; donc mon neveu est un supporter des Lions Indomptables. » Formalisons : A = supporter des Lions, B = porteur d'un maillot vert, C = le neveu. La structure est : « Tous les A sont B ; or C est B ; donc C est A ». C'est exactement l'affirmation du conséquent : le moyen terme B (« porteur d'un maillot vert ») n'est jamais distribué. Le neveu peut porter un maillot vert par simple goût de la couleur, ou pour supporter une équipe de quartier. Le raisonnement de l'oncle est un \textbf{paralogisme} : sa conclusion pourrait être vraie par hasard, mais elle n'est pas \emph{nécessaire}.
\end{exemplebox}

\courssub{Induction et analogie : les autres formes de raisonnement}

La déduction n'est pas la seule manière de raisonner. Deux autres formes méritent d'être distinguées.

\begin{definition}[Induction et analogie]
L'\cle{induction} est un raisonnement qui va \textbf{du particulier au général} : d'un certain nombre de cas observés, on conclut à une loi générale (« tous les corbeaux observés sont noirs, donc tous les corbeaux sont noirs »). Sa conclusion est seulement \textbf{probable}, jamais certaine. Le raisonnement par \cle{analogie} va \textbf{du particulier au particulier} : de la ressemblance de deux situations sur certains points, on conclut qu'elles se ressemblent sur un autre point (« ce quartier ressemble à celui que je connais, donc les prix y sont sans doute pareils »). Sa conclusion est encore plus fragile.
\end{definition}

\begin{center}
\begin{tabularx}{\linewidth}{|l|X|X|X|}\hline
\rowcolor{popblueL} \textbf{Raisonnement} & \textbf{Direction} & \textbf{Conclusion} & \textbf{Exemple} \\ \hline
Déduction & Général $\rightarrow$ particulier & Nécessaire & Syllogisme de Socrate \\ \hline
Induction & Particulier $\rightarrow$ général & Probable & Loi scientifique tirée d'expériences \\ \hline
Analogie & Particulier $\rightarrow$ particulier & Très probable au mieux & Comparaison de deux marchés \\ \hline
\end{tabularx}
\end{center}

\begin{savaistu}[Le problème de l'induction chez Hume]
L'induction est le raisonnement privilégié des sciences expérimentales : la physique, la biologie et la médecine généralisent à partir d'observations répétées. Pourtant, le philosophe écossais David Hume (1711-1776) a montré qu'aucune induction n'est jamais \emph{certaine} : avoir vu le soleil se lever chaque matin ne prouve pas logiquement qu'il se lèvera demain ; on fait seulement le pari, fondé sur l'habitude, que la nature restera constante. C'est le célèbre « problème de l'induction ». Ainsi, même la science la plus solide repose sur des conclusions probables, tandis que la déduction seule offre la nécessité — mais au prix de ne rien apprendre de vraiment nouveau sur le monde.
\end{savaistu}

\courssub{Application : démasquer les paralogismes dans les débats camerounais}

Les plateaux de la CRTV, les émissions interactives des radios privées de Douala et de Yaoundé, et les discussions au carrefour regorgent de raisonnements qui \emph{semblent} rigoureux. Savoir les formaliser est un savoir-faire civique autant que philosophique.

\begin{exoresolu}[Analyse d'un raisonnement radiophonique]
Énoncé. Lors d'un débat radiophonique sur l'éducation, un auditeur déclare : « Tous les élèves sérieux réussissent au Baccalauréat ; or ce jeune a réussi au Baccalauréat ; donc c'est un élève sérieux. » Identifiez les trois termes, formalisez le raisonnement et dites s'il est valide.
\tcblower
Solution détaillée.
\begin{enumerate}
\item \textbf{Conclusion} : « ce jeune est un élève sérieux ». Petit terme P = « ce jeune » ; grand terme G = « élève sérieux » ; moyen terme M = « celui qui réussit au Baccalauréat ».
\item \textbf{Formalisation} : Tous les G sont M (majeure) ; or P est M (mineure) ; donc P est G (conclusion).
\item \textbf{Vérification} : le moyen terme M (« ceux qui réussissent ») n'est jamais distribué : la majeure dit que les élèves sérieux réussissent, mais ne dit pas que \emph{seuls} les élèves sérieux réussissent. Le jeune a pu réussir grâce à la chance, à un examen facile ou à un travail de dernière minute.
\item \textbf{Conclusion de l'analyse} : le raisonnement est un \textbf{paralogisme} du type « affirmation du conséquent ». Pour le rendre valide, il faudrait une majeure convertie : « Seuls les élèves sérieux réussissent au Baccalauréat » (c'est-à-dire « tous ceux qui réussissent sont sérieux ») — mais cette majeure-là est contestable en fait.
\end{enumerate}
\end{exoresolu}

\begin{exercice}[À vous de raisonner]
Au marché de Mokolo, un client affirme : « Tous les commerçants honnêtes affichent leurs prix ; or cette vendeuse affiche ses prix ; donc elle est honnête. » 1) Formalisez ce raisonnement en repérant les trois termes. 2) Montrez, à l'aide d'un diagramme de Venn ou de la règle de distribution du moyen terme, qu'il est invalide. 3) Proposez une majeure qui rendrait la conclusion nécessaire, puis discutez sa vérité.
\end{exercice}

\begin{corrige}[Exercice]
1) Petit terme P = « cette vendeuse » ; grand terme G = « commerçant honnête » ; moyen terme M = « celui qui affiche ses prix ». Forme : « Tous les G sont M ; or P est M ; donc P est G ».
2) Le moyen terme M n'est distribué nulle part : la majeure affirme que les honnêtes affichent leurs prix, mais pas que tous ceux qui affichent leurs prix sont honnêtes. Sur un diagramme de Venn, le cercle des honnêtes est inclus dans celui des afficheurs de prix ; la vendeuse est dans le grand cercle, mais rien ne la place dans le petit. C'est l'affirmation du conséquent : paralogisme.
3) Il faudrait la majeure « Tous ceux qui affichent leurs prix sont honnêtes » (ou « seuls les honnêtes affichent leurs prix »). Le raisonnement deviendrait valide, mais cette majeure est fausse en réalité : un commerçant malhonnête peut afficher des prix trompeurs. La conclusion reste donc incertaine.
\end{corrige}

\begin{aretenir}[L'essentiel de la section]
\begin{itemize}
\item La \cle{déduction} va du général au particulier et produit une conclusion \textbf{nécessaire} si les prémisses sont vraies et la forme valide.
\item Le \cle{syllogisme} relie trois termes (grand, petit, moyen) à travers deux prémisses (majeure, mineure) ; le moyen terme doit être \textbf{distribué au moins une fois} et disparaître de la conclusion.
\item L'\cle{induction} (du particulier au général) et l'\cle{analogie} (du particulier au particulier) ne donnent que des conclusions \textbf{probables}.
\item Le \cle{paralogisme} imite la rigueur : le piège classique est « Tous les A sont B ; or C est B ; donc C est A » (affirmation du conséquent).
\end{itemize}
\end{aretenir}

Ces outils de la déduction et de la détection des erreurs nous préparent à la section suivante, consacrée à l'argumentation : car convaincre honnêtement suppose de savoir enchaîner correctement des raisons, tandis que persuader utilise parfois des procédés qui ne résistent pas à l'épreuve de la logique.

\coursec{L'argumentation : convaincre et persuader}

Dans la section précédente, nous avons étudié le raisonnement déductif et le syllogisme : nous avons vu comment, à partir de prémisses posées comme vraies, la logique permet de tirer une conclusion nécessaire. Mais dans la vie quotidienne — au marché, à la tribune familiale, sur les réseaux sociaux, dans les débats politiques —, personne ne parle en syllogismes parfaits. On cherche surtout à \emph{faire accepter une idée} à celui qui nous écoute : on \cle{argumente}. Cette section est consacrée à l'art d'argumenter : qu'est-ce qu'argumenter ? Quelle différence y a-t-il entre convaincre et persuader ? Quels sont les types d'arguments, leur structure, leurs pièges ? Enfin, comment discuter de manière honnête et critique ?

\courssub{Définition de l'argumentation}

Commençons par une intuition. Un élève veut obtenir de son père l'autorisation de s'inscrire à un concours d'entrée dans une école technique de Douala. Il ne se contente pas de dire : « Je veux y aller. » Il avance des \emph{raisons} : l'école est reconnue par l'État, les frais sont abordables, les débouchés sont réels, un cousin y a été formé et travaille aujourd'hui. Il construit un discours destiné à faire \emph{adhérer} son père à sa thèse : « je dois m'inscrire à ce concours ». Cette démarche, c'est l'argumentation.

\begin{definition}[L'argumentation]
L'\cle{argumentation} est une démarche discursive qui consiste à présenter un ensemble de raisons (les \emph{arguments}) pour faire accepter ou rejeter une thèse (une affirmation, une opinion, une décision) par un auditoire. Argumenter, c'est donc toujours défendre une \cle{thèse} devant quelqu'un, en mobilisant des preuves, des exemples et des raisonnements.
\end{definition}

Trois éléments sont donc indispensables à toute argumentation :
\begin{enumerate}
\item une \textbf{thèse} : ce que l'on veut faire admettre ;
\item des \textbf{arguments} : les raisons avancées pour la soutenir ;
\item un \textbf{auditoire} : celui ou ceux que l'on cherche à convaincre ou à persuader.
\end{enumerate}

\begin{attention}[Argumenter n'est pas affirmer]
Une erreur fréquente à l'examen consiste à confondre \emph{affirmer} et \emph{argumenter}. Dire « la corruption nuit au développement du Cameroun » est une affirmation. Argumenter, c'est dire \emph{pourquoi} : « la corruption détourne les fonds publics destinés aux routes et aux hôpitaux, elle décourage les investisseurs, elle ruine la confiance des citoyens dans les institutions ; elle nuit donc au développement ». Sans raisons données, il n'y a pas d'argumentation.
\end{attention}

\courssub{Convaincre et persuader : deux visées différentes}

Toute argumentation vise l'adhésion de l'auditoire, mais elle peut l'obtenir par deux voies très différentes, que les philosophes distinguent soigneusement.

\begin{definition}[Convaincre et persuader]
\cle{Convaincre}, c'est obtenir l'adhésion de l'autre par la seule force de la \emph{raison} : des arguments logiques, des preuves vérifiables, valables pour tout esprit rationnel. Convaincre vise l'\emph{universalité} : une démonstration mathématique convainc n'importe qui, partout.\\
\cle{Persuader}, c'est obtenir l'adhésion en agissant sur les \emph{émotions}, les sentiments, les préjugés, par le charme du discours, l'image, le rythme des mots. Persuader s'adresse à un public particulier, à un moment particulier : c'est le domaine de la \cle{rhétorique}.
\end{definition}

Prenons un exemple. Pour établir que « le paludisme tue », le médecin qui présente des statistiques du ministère de la Santé publique \emph{convainc}. La campagne publicitaire qui montre un enfant fiévreux dans les bras de sa mère en pleurs, sur une musique triste, \emph{persuade}. Les deux discours aboutissent peut-être à la même adhésion, mais non par le même chemin : l'un passe par l'intelligence, l'autre par la sensibilité.

Faut-il condamner la persuasion ? Pas nécessairement. Un avocat qui plaide, un enseignant qui motive ses élèves, un chef traditionnel qui apaise un conflit utilisent légitimement les émotions. Mais la persuasion devient condamnable lorsqu'elle \emph{manipule} : lorsqu'elle fait adhérer sans raisons valables, en jouant sur la peur ou la flatterie. La philosophie, elle, privilégie le convaincre : elle veut que l'on adhère à une thèse \emph{parce qu'elle est justifiée}, non parce qu'elle est séduisante.

\begin{savaistu}[Les sophistes et la naissance de la rhétorique]
Dans l'Athènes du V\textsuperscript{e} siècle avant J.-C., des maîtres appelés \emph{sophistes} (Protagoras, Gorgias) enseignaient, contre argent, l'art de persuader : la \cle{rhétorique}. Ils se vantaient de pouvoir « rendre forte la cause la plus faible » et de faire triompher n'importe quelle thèse devant un tribunal ou une assemblée. Socrate, puis Platon, les ont violemment critiqués : pour eux, le sophiste ne cherche pas la vérité, mais seulement la victoire dans la discussion. Aristote, plus nuancé, rédigea une \emph{Rhétorique} où il distingue trois moyens de persuasion : le caractère de l'orateur (\emph{ethos}), les émotions du public (\emph{pathos}) et la rigueur du raisonnement (\emph{logos}). Ce débat est d'une actualité brûlante à l'ère des réseaux sociaux.
\end{savaistu}

\courssub{Les types d'arguments}

Tous les arguments ne se valent pas. Il faut savoir les reconnaître pour en mesurer la force.

\textbf{L'argument d'autorité.} Il consiste à appuyer une thèse sur la parole d'une personne compétente, d'une institution ou d'un texte de référence : « selon l'INS, le taux de chômage des jeunes est de... ». Il est légitime lorsque l'autorité invoquée est réellement compétente dans le domaine concerné ; il devient fallacieux lorsque l'autorité est étrangère au sujet (un chanteur célèbre qui donne un avis médical) ou lorsqu'elle sert à interdire toute discussion.

\begin{exemplebox}[Discuter un argument d'autorité]
« Le ministre a dit que le pont sera construit, donc c'est vrai. » Cet argument d'autorité est faible : le ministre est certes une autorité politique, mais sa déclaration est une \emph{promesse}, non une preuve. Pour être convaincant, il faudrait des éléments vérifiables : le vote du budget, le lancement des travaux, l'entreprise retenue. L'argument d'autorité ne dispense jamais de l'examen critique des faits.
\end{exemplebox}

\textbf{L'argument par l'exemple.} On illustre une thèse générale par un cas particulier : « la jeunesse camerounaise est entreprenante : voyez ces jeunes de Bafoussam qui ont créé une start-up de transformation du manioc ». L'exemple rend la thèse concrète et vivante, mais un seul exemple ne prouve pas une généralité : il \emph{illustre} plus qu'il ne \emph{démontre}.

\textbf{L'argument par analogie.} On raisonne par ressemblance : « un pays est comme une famille ; or une famille doit vivre selon ses moyens ; donc l'État doit limiter ses dépenses ». L'analogie éclaire, mais elle est fragile : deux réalités semblables sur un point peuvent différer sur l'essentiel (un État n'est pas une famille : il peut lever l'impôt, battre monnaie...).

\textbf{L'argument par les causes et les conséquences.} On justifie une thèse en montrant ses causes (« cette crise vient de... ») ou ses effets (« si nous faisons cela, il s'ensuivra... »). C'est un argument solide \emph{si} le lien de causalité est réellement établi, et non simplement supposé.

\textbf{L'argument \emph{ad hominem}.} Au lieu de réfuter la thèse de l'adversaire, on attaque sa personne : « tu défends ce projet, mais toi-même tu as été condamné, donc ton avis ne compte pas ». C'est un procédé rhétoriquement efficace mais logiquement nul : la valeur d'une idée ne dépend pas de la moralité de celui qui la formule.

\courssub{La structure d'un texte argumentatif}

Un texte argumentatif bien construit — comme la dissertation philosophique que vous rédigerez au Baccalauréat technique — obéit à une architecture précise : on énonce la thèse, on la soutient par des arguments, on étaye chaque argument par des exemples, on examine les objections possibles et on y répond, avant de conclure.

\begin{popfigure}
\begin{center}
\begin{tikzpicture}[
box/.style={draw=popblue, thick, fill=popblueL, rounded corners=2pt, align=center, minimum height=0.85cm, text width=3.1cm, font=\small},
obj/.style={draw=poporange, thick, fill=poporangeL, rounded corners=2pt, align=center, minimum height=0.85cm, text width=3.1cm, font=\small},
fleche/.style={-{Stealth[length=3mm]}, very thick, popdark}]
\node[box, align=center] (t) at (0,0){\textbf{Thèse}\\ ce que je défends};
\node[box, align=center] (a1) at (-3.6,-2){\textbf{Argument 1}\\ + exemple};
\node[box, align=center] (a2) at (0,-2){\textbf{Argument 2}\\ + exemple};
\node[box, align=center] (a3) at (3.6,-2){\textbf{Argument 3}\\ + exemple};
\node[obj, align=center] (o) at (0,-4){\textbf{Objection}\\ + réponse};
\node[box, fill=popgreenL, draw=popgreen, align=center] (c) at (0,-6){\textbf{Conclusion}\\ thèse confirmée};
\draw[fleche] (t) -- (a1);
\draw[fleche] (t) -- (a2);
\draw[fleche] (t) -- (a3);
\draw[fleche] (a1) -- (o);
\draw[fleche] (a2) -- (o);
\draw[fleche] (a3) -- (o);
\draw[fleche] (o) -- (c);
\end{tikzpicture}
\end{center}
\legende{Structure d'un texte argumentatif : la thèse est soutenue par plusieurs arguments illustrés d'exemples ; les objections sont examinées et réfutées ; la conclusion reprend la thèse désormais justifiée.}
\end{popfigure}

L'étape des \textbf{objections} est la marque d'un esprit honnête : au lieu de cacher les difficultés, on les formule soi-même (« on pourrait objecter que... ») et on y répond. C'est exactement la démarche de la dissertation philosophique, où l'on confronte les thèses des auteurs avant de trancher.

\courssub{Les sophismes : les pièges du discours}

\begin{definition}[Le sophisme]
Un \cle{sophisme} (ou \emph{paralogisme}) est un raisonnement fallacieux, c'est-à-dire faux, construit avec l'intention de tromper, ou du moins présenté comme valide alors qu'il ne l'est pas. Le sophisme a l'\emph{apparence} de la rigueur logique, mais sa conclusion ne découle pas réellement de ses prémisses.
\end{definition}

\begin{center}
\begin{tabularx}{\linewidth}{|l|X|X|}
\hline
\rowcolor{popblueL} \textbf{Sophisme} & \textbf{Mécanisme} & \textbf{Exemple camerounais} \\
\hline
Pente glissante & On prétend qu'un premier pas entraîne fatalement une catastrophe finale, sans le prouver. & « Si on autorise les élèves à contester une note, bientôt ils contesteront tout, et ce sera l'anarchie dans tous les lycées. » \\
\hline
Fausse analogie & On compare deux situations profondément différentes. & « On dirige bien une famille sans consulter les enfants ; on doit donc diriger une commune sans consulter les habitants. » \\
\hline
Généralisation hâtive & On tire une loi générale d'un ou deux cas particuliers. & « Mon voisin de quartier est malhonnête ; tous les gens de son village sont donc des voleurs. » \\
\hline
Appel à la peur & On fait adhérer en effrayant, sans argument rationnel. & « Si vous ne votez pas pour moi, ce sera la guerre et la misère pour vos enfants. » \\
\hline
Homme de paille & On déforme la thèse adverse pour la réfuter plus facilement. & « Tu demandes plus de moyens pour les écoles rurales ? Tu veux donc supprimer les lycées de Yaoundé ! » \\
\hline
\end{tabularx}
\end{center}

\begin{attention}[L'erreur la plus fréquente : l'attaque \emph{ad hominem}]
Dans les débats — y compris dans les copies —, l'erreur la plus répandue consiste à attaquer la \emph{personne} au lieu de réfuter sa \emph{thèse}. Dire « ton avis est nul parce que tu es jeune / pauvre / de telle région » ne réfute rien : même une personne méprisable peut dire une vérité, et une personne respectable peut se tromper. En philosophie, on discute les \emph{idées}, jamais les personnes.
\end{attention}

\courssub{Analyser un discours réel et discuter avec éthique}

\begin{methode}[Analyser un discours argumentatif en quatre étapes]
\begin{enumerate}
\item \textbf{Repérer la thèse} : que veut-on me faire accepter ? (Souvent exprimée au début ou à la fin du discours.)
\item \textbf{Classer les arguments} : s'agit-il d'arguments d'autorité, d'exemples, d'analogies, de causes et conséquences ?
\item \textbf{Évaluer leur pertinence} : les arguments sont-ils vrais ? Soutiennent-ils réellement la thèse ? Sont-ils suffisants ?
\item \textbf{Débusquer les sophismes} : y a-t-il appel à la peur, généralisation hâtive, attaque personnelle, fausse analogie ?
\end{enumerate}
\end{methode}

Appliquons cette méthode à des discours que vous rencontrez chaque jour. Dans un débat politique télévisé camerounais, un candidat déclare : « Mes adversaires n'ont jamais rien construit ; moi, j'ai l'expérience ; sans moi, c'est le chaos. » Analyse : la thèse est « votez pour moi » ; le premier argument est \emph{ad hominem} (attaque des adversaires, non de leur programme) ; le second est un argument d'autorité (sa propre expérience) ; le troisième est un appel à la peur. Aucun argument ne porte sur le \emph{programme} : le discours persuade, mais ne convainc pas.

Sur Facebook ou WhatsApp, une publication affirme : « Un guérisseur de l'Ouest soigne le diabète ; mon oncle en est guéri ; partagez massivement ! » Analyse : argument par l'exemple unique (généralisation hâtive), absence de toute preuve médicale, appel à l'émotion. La démarche critique consiste à demander : où sont les études ? Quel est le taux de réussite ? Que dit le corps médical ?

De même, une publicité qui montre une star du football buvant une boisson ne prouve rien sur la qualité du produit : c'est un argument d'autorité déplacé, doublé d'une séduction par l'image.

\begin{exoresolu}[Analyse critique d'un discours]
Énoncé. Analysez le discours suivant : « Le proviseur a annoncé que le port de l'uniforme sera obligatoire. Il a raison : d'abord, l'uniforme efface les différences entre riches et pauvres ; ensuite, dans les pays développés, les élèves sérieux portent l'uniforme ; enfin, ceux qui refusent l'uniforme sont des élèves indisciplinés qui veulent le désordre. » Identifiez la thèse, classez les arguments et repérez les sophismes éventuels.
\tcblower
\textbf{Solution détaillée.}
\begin{enumerate}
\item \emph{Thèse} : « le port de l'uniforme doit être obligatoire » (le proviseur « a raison »).
\item \emph{Premier argument} : « l'uniforme efface les différences sociales ». C'est un argument par les conséquences, pertinent et recevable, même s'il resterait à vérifier (le coût de l'uniforme peut lui-même créer une inégalité).
\item \emph{Deuxième argument} : « dans les pays développés, les élèves sérieux portent l'uniforme ». C'est une fausse analogie doublée d'une généralisation hâtive : tous les pays développés n'imposent pas l'uniforme, et la réussite scolaire dépend de nombreux autres facteurs (qualité des enseignants, équipements, encadrement familial).
\item \emph{Troisième argument} : « ceux qui refusent sont des indisciplinés qui veulent le désordre ». C'est une attaque \emph{ad hominem} combinée à un homme de paille : on déforme la position des opposants (qui peuvent avoir des raisons économiques légitimes) et on attaque leur personne au lieu de répondre à leurs arguments.
\end{enumerate}
\emph{Conclusion} : le discours contient un argument recevable et deux sophismes ; il persuade plus qu'il ne convainc.
\end{exoresolu}

Terminons par l'\textbf{éthique de la discussion}. Argumenter n'est pas faire la guerre : c'est chercher ensemble le vrai. Trois règles s'imposent. Premièrement, \emph{écouter} réellement l'autre avant de répondre : on ne peut réfuter que ce que l'on a compris. Deuxièmement, \emph{réfuter l'argument et non la personne} : la critique doit porter sur les idées, leurs preuves et leur logique. Troisièmement, \emph{accepter la contradiction} : celui qui refuse tout désaccord renonce à la rationalité ; comme le montre Socrate dans les dialogues de Platon, c'est souvent en étant contredit que l'on progresse vers la vérité. Un débat où l'on peut changer d'avis est un débat réussi, non perdu.

\begin{aretenir}[L'essentiel]
\begin{itemize}
\item Argumenter, c'est présenter des raisons pour faire adhérer un auditoire à une thèse.
\item \emph{Convaincre} s'adresse à la raison et vise l'universel ; \emph{persuader} s'adresse aux émotions et vise un public particulier.
\item Les principaux types d'arguments : autorité, exemple, analogie, causes et conséquences ; l'attaque \emph{ad hominem} est un procédé, non un argument valide.
\item Un texte argumentatif s'organise : thèse, arguments, exemples, objections et réponses, conclusion.
\item Les sophismes (pente glissante, fausse analogie, généralisation hâtive, appel à la peur, homme de paille) imitent le raisonnement pour tromper : il faut savoir les débusquer.
\item L'éthique de la discussion exige d'écouter, de réfuter les idées et non les personnes, et d'accepter la contradiction.
\end{itemize}
\end{aretenir}

\begin{exercice}[Entraînement à l'analyse de discours]
Relevez dans un débat politique télévisé, une publicité ou une publication WhatsApp de votre entourage un passage argumentatif. 1) Formulez sa thèse. 2) Classez ses arguments. 3) Identifiez au moins un sophisme et justifiez votre réponse. 4) Reformulez le passage en un argument rationnellement acceptable.
\end{exercice}

\begin{corrige}[Exercice — éléments de correction]
La correction dépend du passage choisi, mais la copie attendue doit suivre la méthode en quatre étapes. Exemple de traitement : pour la publicité « Bois X, la boisson des champions ! », la thèse implicite est « achetez X » ; l'argument est un argument d'autorité déplacé (le champion est compétent en football, non en nutrition) doublé d'une séduction par l'image ; le sophisme est donc l'autorité hors de son domaine ; reformulation honnête : « cette boisson contient tels nutriments, dont l'effet sur l'organisme est établi par telles études ». Toute copie qui distingue clairement thèse, arguments et sophismes, et qui propose une reformulation rationnelle, est recevable.
\end{corrige}

\coursec{La démonstration}

Nous avons vu, dans la section précédente, que l'argumentation vise à \emph{convaincre} ou à \emph{persuader} un auditoire : ses conclusions restent toujours, en principe, discutables, car elles s'appuient sur des valeurs, des opinions ou des faits contingents. Mais existe-t-il des discours dont la conclusion ne peut pas être refusée ? Lorsqu'un élève affirme que la somme de deux nombres pairs est toujours paire, il ne demande pas l'adhésion de son auditoire : il prétend établir une vérité qui s'impose à tous, partout et toujours. Un tel discours relève non plus de l'argumentation, mais de la \cle{démonstration}. C'est l'objet de cette section : comprendre ce qu'est démontrer, connaître les trois grands types de démonstration, et en mesurer la portée comme les limites.

\courssub{Définition et enjeux de la démonstration}

\begin{definition}[La démonstration]
La \cle{démonstration} est un raisonnement déductif qui établit la \textbf{vérité nécessaire} d'une proposition (la conclusion) à partir de prémisses \textbf{vraies ou admises} (axiomes, définitions, théorèmes déjà établis). Si les prémisses sont vraies et si l'enchaînement logique est correct, alors la conclusion est \emph{nécessairement} vraie : il est impossible de la nier sans se contredire.
\end{definition}

L'intuition est simple : démontrer, c'est \emph{contraindre la pensée}. Dans un débat politique à la télévision camerounaise, chaque intervenant peut contester les arguments de l'autre sans contradiction logique. En revanche, celui qui a démontré un théorème de géométrie ne laisse à son interlocuteur que deux choix : accepter la conclusion, ou refuser les prémisses. La démonstration possède donc une \cle{visée cognitive} (établir le vrai) et non une visée persuasive (gagner l'adhésion).

\begin{center}
\begin{tabularx}{\linewidth}{|l|X|X|}
\hline
\rowcolor{popblueL} \textbf{Critère} & \textbf{Argumentation} & \textbf{Démonstration} \\
\hline
Visée & Persuader ou convaincre & Établir une vérité \\
\hline
Prémisses & Opinions, valeurs, faits & Axiomes, définitions, théorèmes \\
\hline
Conclusion & Probable, discutable & Nécessaire, certaine \\
\hline
Domaine & Morale, politique, rhétorique & Mathématiques, logique, sciences \\
\hline
Réfutation possible & Oui, sans contradiction & Non, sauf à nier les prémisses \\
\hline
\end{tabularx}
\end{center}

\begin{aretenir}[La distinction fondamentale]
L'argumentation produit du \emph{vraisemblable} ; la démonstration produit du \emph{nécessaire}. Aristote distinguait déjà le \emph{syllogisme dialectique} (à partir de prémisses probables) du \emph{syllogisme apodictique} ou démonstratif (à partir de prémisses vraies et premières).
\end{aretenir}

\courssub{Les trois types de démonstration}

\textbf{1. La démonstration directe.}\par
C'est la voie la plus naturelle : on part des axiomes, des définitions et des hypothèses, et l'on progresse par enchaînements déductifs jusqu'à la conclusion. Schématiquement : $p \Rightarrow r \Rightarrow s \Rightarrow q$. Chaque étape découle nécessairement de la précédente.

\begin{exemplebox}[Démonstration directe : somme de deux nombres pairs]
\textbf{Thèse :} la somme de deux nombres pairs est un nombre pair.\par
\textbf{Démonstration :} soit $m$ et $n$ deux nombres pairs. Par définition d'un nombre pair, il existe deux entiers $a$ et $b$ tels que $m = 2a$ et $n = 2b$. Alors :
\[ m + n = 2a + 2b = 2(a+b). \]
Or $a+b$ est un entier, donc $m+n$ s'écrit $2 \times (\text{un entier})$ : c'est, par définition, un nombre pair. \textbf{La thèse est démontrée.}
\end{exemplebox}

Remarquons qu'aucun exemple numérique ($2+4=6$, $8+10=18$...) ne suffirait : seule la lettre, qui représente \emph{n'importe quel} entier, confère à la preuve son caractère universel.

\textbf{2. La démonstration par l'absurde.}\par
L'idée est audacieuse : pour prouver qu'une thèse $p$ est vraie, on suppose qu'elle est \emph{fausse} (on pose $\neg p$), on en tire rigoureusement les conséquences, et l'on aboutit à une \textbf{contradiction}. Or une contradiction est logiquement inacceptable : c'est donc l'hypothèse de départ, $\neg p$, qui est fausse, et par conséquent $p$ est vraie.

\begin{methode}[Démontrer par l'absurde en quatre étapes]
\begin{enumerate}
\item \textbf{Poser la négation} de la thèse à démontrer (hypothèse provisoire).
\item \textbf{En tirer les conséquences} par des déductions rigoureuses.
\item \textbf{Exhiber une contradiction} : une proposition à la fois vraie et fausse.
\item \textbf{Conclure} : l'hypothèse de départ est fausse, donc la thèse est vraie.
\end{enumerate}
\end{methode}

\begin{exoresolu}[L'irrationalité de $\sqrt{2}$ (complément utile au Bac)]
\textbf{Énoncé.} Démontrer que $\sqrt{2}$ n'est pas un nombre rationnel, c'est-à-dire qu'il ne peut pas s'écrire comme une fraction $\dfrac{a}{b}$ de deux entiers.
\tcblower
\textbf{Solution.} Raisonnons par l'absurde.
\begin{enumerate}
\item \emph{Négation de la thèse :} supposons que $\sqrt{2}$ est rationnel. Alors il existe deux entiers $a$ et $b$ ($b \neq 0$), \textbf{premiers entre eux} (fraction irréductible), tels que $\sqrt{2} = \dfrac{a}{b}$.
\item \emph{Conséquences :} en élevant au carré, $2 = \dfrac{a^2}{b^2}$, donc $a^2 = 2b^2$. Ainsi $a^2$ est pair, donc $a$ est pair (le carré d'un impair est impair). Posons $a = 2k$. Alors $4k^2 = 2b^2$, d'où $b^2 = 2k^2$ : $b^2$ est pair, donc $b$ est pair.
\item \emph{Contradiction :} $a$ et $b$ sont tous deux pairs, donc divisibles par 2, alors qu'ils étaient supposés premiers entre eux. Contradiction.
\item \emph{Conclusion :} l'hypothèse est fausse : $\sqrt{2}$ est \textbf{irrationnel}.
\end{enumerate}
\end{exoresolu}

\textbf{3. La démonstration par contraposée.}\par
Il est parfois plus facile de prouver une implication « à l'envers ».

\begin{propriete}[Équivalence de la contraposée]
L'implication $p \Rightarrow q$ est logiquement équivalente à sa \cle{contraposée} $\neg q \Rightarrow \neg p$. Démontrer l'une, c'est démontrer l'autre.
\end{propriete}

\begin{exemplebox}[Exemple de contraposée]
Pour démontrer « si $n^2$ est pair, alors $n$ est pair », on démontre la contraposée : « si $n$ est impair, alors $n^2$ est impair ». Or si $n = 2k+1$, alors $n^2 = 4k^2 + 4k + 1 = 2(2k^2+2k) + 1$, qui est impair. La contraposée est établie, donc l'implication initiale aussi.
\end{exemplebox}

\begin{popfigure}
\begin{center}
\begin{tikzpicture}[
  box/.style={draw=popblue, thick, rounded corners, fill=popblueL, text width=3.1cm, align=center, minimum height=1.6cm, font=\small},
  arr/.style={-{Stealth[length=2.5mm]}, thick, popdark}
]
\node[box, align=center] (p1) at (0,2.2){\textbf{Prémisses}\\ axiomes, définitions, hypothèses};
\node[box, align=center] (c1) at (6.5,2.2){\textbf{Conclusion $q$}\\ (nécessaire)};
\draw[arr] (p1) -- (c1) node[midway, above, font=\small\itshape]{déductions successives};
\node[font=\small\bfseries, popblue] at (3.25,3.4) {Démonstration directe};

\node[box, draw=poporange, fill=poporangeL, align=center] (p2) at (0,-0.6){\textbf{Hypothèse : $\neg q$}\\ (négation de la thèse)};
\node[box, draw=poporange, fill=poporangeL, align=center] (c2) at (6.5,-0.6){\textbf{Contradiction}\\ donc $q$ est vraie};
\draw[arr, poporange] (p2) -- (c2) node[midway, above, font=\small\itshape]{conséquences de $\neg q$};
\node[font=\small\bfseries, poporange] at (3.25,0.6) {Démonstration par l'absurde};

\node[box, draw=popgreen, fill=popgreenL] (p3) at (0,-3.4) {\textbf{Hypothèse : $\neg q$}};
\node[box, draw=popgreen, fill=popgreenL, align=center] (c3) at (6.5,-3.4){\textbf{Conclusion : $\neg p$}\\ donc $p \Rightarrow q$};
\draw[arr, popgreen!70!black] (p3) -- (c3) node[midway, above, font=\small\itshape]{déduction directe};
\node[font=\small\bfseries, popgreen!60!black] at (3.25,-2.2) {Démonstration par contraposée};
\end{tikzpicture}
\end{center}
\legende{Les trois types de démonstration. La voie directe va des prémisses à la conclusion ; la voie par l'absurde suppose la négation de la thèse pour aboutir à une contradiction ; la contraposée démontre $\neg q \Rightarrow \neg p$ à la place de $p \Rightarrow q$.}
\end{popfigure}

\courssub{Étude de cas : l'arbre logique de l'irrationalité de $\sqrt{2}$}

Visualisons la structure complète du raisonnement par l'absurde sous forme d'arbre logique : chaque nœud est une étape, et l'arbre se referme sur la contradiction qui invalide l'hypothèse initiale.

\begin{popfigure}
\begin{center}
\begin{tikzpicture}[
  every node/.style={font=\small},
  etape/.style={draw=popblue, thick, rounded corners, fill=popblueL, align=center, text width=4.6cm},
  contra/.style={draw=poporange, very thick, rounded corners, fill=poporangeL, align=center, text width=4.6cm},
  concl/.style={draw=popgreen!70!black, very thick, rounded corners, fill=popgreenL, align=center, text width=4.6cm},
  arr/.style={-{Stealth[length=2.5mm]}, thick, popdark}
]
\node[etape] (h) at (0,0) {\textbf{Hypothèse} : $\sqrt{2}$ est rationnel, donc $\sqrt{2}=\dfrac{a}{b}$ avec $a$, $b$ premiers entre eux};
\node[etape] (e1) at (-3.4,-2.2) {$a^2 = 2b^2$ donc $a^2$ pair, donc $a$ pair : $a=2k$};
\node[etape] (e2) at (3.4,-2.2) {$4k^2 = 2b^2$ donc $b^2 = 2k^2$, donc $b$ pair};
\node[contra] (c) at (0,-4.4) {\textbf{Contradiction} : $a$ et $b$ sont pairs, donc non premiers entre eux};
\node[concl] (f) at (0,-6.4) {\textbf{Conclusion} : l'hypothèse est fausse, $\sqrt{2}$ est irrationnel};
\draw[arr] (h) -- (e1);
\draw[arr] (h) -- (e2);
\draw[arr] (e1) -- (c);
\draw[arr] (e2) -- (c);
\draw[arr, popgreen!70!black] (c) -- (f);
\end{tikzpicture}
\end{center}
\legende{Arbre logique de la démonstration par l'absurde de l'irrationalité de $\sqrt{2}$. De l'hypothèse initiale découlent deux conséquences convergentes ($a$ pair et $b$ pair) qui se contredisent mutuellement : l'hypothèse doit être rejetée.}
\end{popfigure}

\courssub{Le rôle des axiomes, des postulats et des définitions}

Toute démonstration suppose un \emph{point de départ}. Mais si chaque vérité devait être démontrée à partir d'une autre, on remonterait à l'infini. Il faut donc des \textbf{vérités premières}, admises sans démonstration :

\begin{itemize}
\item les \cle{axiomes} : propositions évidentes par elles-mêmes, communes à toutes les sciences (ex. : « deux grandeurs égales à une même troisième sont égales entre elles ») ;
\item les \cle{postulats} : propositions admises à la base d'une théorie particulière (ex. : « par un point extérieur à une droite, on peut mener une et une seule parallèle ») ;
\item les \cle{définitions} : conventions qui fixent le sens précis des termes employés.
\end{itemize}

\begin{savaistu}[Euclide et les \emph{Éléments}]
Vers 300 avant J.-C., le mathématicien grec \textbf{Euclide} construit, dans ses \emph{Éléments}, toute la géométrie à partir de seulement \textbf{cinq postulats}, quelques axiomes et des définitions. De ces bases minimales, il déduit des centaines de théorèmes, dont le célèbre théorème attribué à Pythagore. Cet ouvrage est resté, pendant plus de 2000 ans, le modèle absolu de la rigueur scientifique : Spinoza rédigera même son \emph{Éthique} « selon l'ordre géométrique », avec axiomes, définitions et démonstrations.
\end{savaistu}

\courssub{Les limites de la démonstration}

La démonstration est le sommet de la rigueur, mais elle a des \textbf{limites} qu'il faut connaître pour ne pas lui demander plus qu'elle ne peut donner.

\begin{itemize}
\item \textbf{Les principes premiers sont indémontrables.} Aristote l'a montré : si l'on exigeait de tout démontrer, il faudrait démontrer les prémisses de chaque démonstration, et ainsi de suite à l'infini. La démonstration doit donc s'arrêter à des principes admis : le principe de non-contradiction, par exemple, ne se démontre pas, car toute démonstration le présuppose déjà.
\item \textbf{La démonstration ne vaut que dans son système.} Une vérité démontrée l'est relativement aux axiomes choisis : en changeant le postulat des parallèles, on obtient des géométries non euclidiennes, tout aussi cohérentes.
\item \textbf{Tout ne relève pas de la démonstration.} Les questions morales, esthétiques ou existentielles (« Faut-il toujours dire la vérité ? », « Cette œuvre est-elle belle ? ») ne se démontrent pas : elles s'\emph{argumentent}. Exiger une démonstration là où seule l'argumentation est possible, c'est se condamner au scepticisme.
\end{itemize}

\begin{attention}[Piège fréquent : confondre induction et démonstration]
Accumuler des exemples ne constitue \textbf{jamais} une démonstration. Vérifier que $2+4=6$, $6+8=14$, $10+22=32$ sont pairs ne prouve pas que la somme de deux pairs est \emph{toujours} paire : c'est une simple \cle{induction}, qui donne une probabilité, non une nécessité. Au Baccalauréat, un candidat qui « prouve » une thèse philosophique ou mathématique par une liste d'exemples commet cette confusion. Seul un raisonnement portant sur le cas \emph{général} (lettres, définitions, axiomes) a valeur de démonstration.
\end{attention}

\courssub{Application : rédiger une démonstration courte et justifiée}

\begin{methode}[Rédiger une démonstration en cinq points]
\begin{enumerate}
\item \textbf{Énoncer clairement la thèse} à démontrer (« Nous allons montrer que... »).
\item \textbf{Choisir le type de démonstration} : directe, par l'absurde ou par contraposée.
\item \textbf{Rappeler les prémisses} : définitions, axiomes ou théorèmes utilisés.
\item \textbf{Dérouler les étapes} en justifiant chacune (« par définition », « d'après le théorème de... », « donc »).
\item \textbf{Conclure explicitement} : « La thèse est démontrée » (on écrit parfois CQFD).
\end{enumerate}
\end{methode}

\begin{exoresolu}[Démonstration en philosophie]
\textbf{Énoncé.} Démontrer, par l'absurde, la thèse suivante : « On ne peut pas douter de tout sans affirmer au moins une vérité. »
\tcblower
\textbf{Solution.}
\begin{enumerate}
\item \emph{Négation de la thèse :} supposons qu'il soit possible de douter absolument de tout, sans rien affirmer.
\item \emph{Conséquences :} celui qui doute de tout doute aussi qu'il doute... mais douter qu'il doute est encore un doute : il est donc certain, au moins, qu'il doute. (C'est l'intuition de Descartes : « Je pense, donc je suis. »)
\item \emph{Contradiction :} l'hypothèse affirmait qu'aucune vérité n'est certaine ; or nous venons d'en trouver une : l'existence même du doute.
\item \emph{Conclusion :} il est impossible de douter de tout sans affirmer au moins une vérité. La thèse est démontrée.
\end{enumerate}
\end{exoresolu}

\begin{exercice}[À faire]
\begin{enumerate}
\item Démontrer, par contraposée, que si $n^2$ est impair, alors $n$ est impair.
\item Démontrer, par l'absurde, qu'il n'existe pas de plus grand nombre entier.
\item Dans un sujet de dissertation (« La vérité se démontre-t-elle toujours ? »), rédiger un paragraphe argumenté distinguant démonstration et argumentation.
\end{enumerate}
\end{exercice}

\begin{corrige}[Exercice]
\textbf{1.} Contraposée : « si $n$ est pair, alors $n^2$ est pair ». Si $n = 2k$, alors $n^2 = 4k^2 = 2(2k^2)$, qui est pair. La contraposée est vraie, donc l'implication initiale aussi.\par
\textbf{2.} Par l'absurde : supposons qu'il existe un plus grand entier $N$. Alors $N+1$ est un entier et $N+1 > N$ : contradiction avec la définition de $N$. Donc il n'existe pas de plus grand entier.\par
\textbf{3.} Le paragraphe doit montrer que la démonstration n'atteint que les vérités formelles (mathématiques, logique) à partir d'axiomes indémontrables, tandis que les vérités morales ou métaphysiques relèvent de l'argumentation, dont les conclusions restent probables : tout ne se démontre donc pas.
\end{corrige}

\begin{aretenir}[L'essentiel de la section]
Démontrer, c'est établir une conclusion \textbf{nécessaire} à partir de prémisses vraies ou admises. Trois voies : la démonstration \textbf{directe} (des prémisses à la conclusion), la démonstration \textbf{par l'absurde} (la négation de la thèse conduit à une contradiction) et la démonstration \textbf{par contraposée} ($\neg q \Rightarrow \neg p$ équivaut à $p \Rightarrow q$). Toute démonstration repose sur des axiomes et des définitions indémontrables (Aristote, Euclide) : la démonstration est le sommet de la rigueur, mais elle n'épuise pas le champ du vrai.
\end{aretenir}

\coursec{Applications : rédiger et justifier une démarche}

Dans la section précédente, nous avons vu que la démonstration est la forme la plus rigoureuse de l'argumentation : elle établit la vérité d'une conclusion à partir de prémisses, selon des règles de validité précises (modus ponens, raisonnement par l'absurde, tables de vérité). Mais à quoi sert tout cet apprentissage ? La logique n'est pas un jeu abstrait réservé aux mathématiciens : elle est l'outil quotidien du candidat au Probatoire et au Baccalauréat technique. Cette dernière section montre comment \cle{appliquer} les notions de l'unité à des situations concrètes : rédiger une dissertation, justifier une réponse, vérifier la cohérence d'un devoir, formaliser des phrases de la vie courante et débattre de manière contradictoire.

\courssub{Mobiliser la logique dans la dissertation philosophique}

\textbf{L'intuition.} Une dissertation n'est pas un récital de souvenirs : c'est un \emph{raisonnement organisé}. Le correcteur ne demande pas « que sais-tu ? » mais « sais-tu penser de manière ordonnée ? ». Or penser de manière ordonnée, c'est exactement ce que la logique nous a appris à faire : poser des prémisses, enchaîner valablement, conclure.

\textbf{La problématique comme question directrice.} Tout devoir commence par l'analyse du sujet : on repère les mots-clés, on définit les termes, puis on formule la \cle{problématique}, c'est-à-dire la question précise à laquelle tout le devoir devra répondre. Une problématique bien posée fonctionne comme l'hypothèse d'une démonstration : elle donne sa direction à l'ensemble du travail.

\textbf{Le plan dialectique.} Le plan attendu en Terminale technique est le plan \cle{dialectique} : thèse, antithèse, synthèse.

\begin{itemize}
\item \textbf{Thèse} : première réponse à la problématique, soutenue par des arguments et des exemples.
\item \textbf{Antithèse} : réponse opposée ou critique de la thèse, elle aussi argumentée (et non simplement niée).
\item \textbf{Synthèse} : dépassement qui tranche la question en tenant compte du débat, sans simple compromis mou.
\end{itemize}

\textbf{Les transitions argumentées.} Entre les parties, la \cle{transition} joue le rôle du connecteur logique : elle rappelle ce qui vient d'être établi, montre sa limite et annonce l'étape suivante. Exemple : « Nous avons montré que la liberté semble incompatible avec l'obéissance aux lois. Mais cette opposition n'est-elle pas apparente ? C'est ce qu'il faut maintenant examiner. »

\begin{popfigure}
\begin{center}
\begin{tikzpicture}[
node distance=0.55cm,
etape/.style={rectangle, rounded corners, draw=popblue, fill=popblueL, text width=6.2cm, align=center, minimum height=0.9cm, font=\small},
fleche/.style={-{Stealth[length=2.5mm]}, thick, popdark}]
\node[etape, fill=popgoldL, draw=popgold, align=center] (intro){\textbf{Introduction}\\ accroche, définitions, problématique, annonce du plan};
\node[etape, below=of intro, align=center] (these){\textbf{Thèse}\\ argument 1, argument 2, exemple};
\node[etape, below=of these, align=center] (anti){\textbf{Antithèse}\\ objection argumentée, contre-exemples};
\node[etape, below=of anti, align=center] (synth){\textbf{Synthèse}\\ dépassement, réponse nuancée à la problématique};
\node[etape, fill=popgreenL, draw=popgreen, below=of synth, align=center] (concl){\textbf{Conclusion}\\ bilan, réponse claire, ouverture};
\draw[fleche] (intro) -- (these);
\draw[fleche] (these) -- node[right, font=\footnotesize, popmuted]{transition} (anti);
\draw[fleche] (anti) -- node[right, font=\footnotesize, popmuted]{transition} (synth);
\draw[fleche] (synth) -- (concl);
\end{tikzpicture}
\end{center}
\legende{Organigramme de la méthode de la dissertation dialectique : chaque étape s'enchaîne logiquement à la précédente grâce aux transitions.}
\end{popfigure}

\begin{methode}[Rédiger une réponse justifiée]
Pour chaque paragraphe du développement, suivre quatre étapes :
\begin{enumerate}
\item \textbf{Annoncer} : formuler clairement la thèse ou l'idée défendue dans le paragraphe.
\item \textbf{Argumenter} : donner la raison qui la fonde, en la rattachant à une notion du cours ou à un auteur.
\item \textbf{Illustrer} : appuyer par un exemple concret (vie courante, réalité camerounaise, référence culturelle).
\item \textbf{Conclure partiellement} : tirer la conséquence et la relier explicitement à la problématique par un connecteur logique (\emph{donc}, \emph{par conséquent}, \emph{il s'ensuit que}).
\end{enumerate}
\end{methode}

\begin{attention}[Erreur fréquente : enchaîner les idées sans connecteurs]
Beaucoup de copies juxtaposent les phrases : « L'homme est libre. La société impose des lois. La liberté est limitée. » Sans connecteurs, le correcteur ne voit pas le raisonnement. Il faut écrire : « L'homme se croit libre ; \emph{or} la société impose des lois ; \emph{par conséquent} sa liberté est limitée. » Les mots \emph{donc}, \emph{car}, \emph{or}, \emph{cependant}, \emph{par conséquent}, \emph{en effet} sont les articulations visibles de la logique du devoir.
\end{attention}

\courssub{Vérifier la cohérence d'un devoir}

Avant de rendre une copie, il faut la relire comme un logicien, selon trois contrôles :

\begin{enumerate}
\item \textbf{Absence de contradiction} : je ne dois pas affirmer $p$ dans la thèse et $\neg p$ dans la synthèse sans l'avoir justifié. Une contradiction non assumée détruit la démonstration (on l'a vu avec le raisonnement par l'absurde).
\item \textbf{Validité des enchaînements} : chaque conclusion partielle suit-elle réellement des arguments donnés ? Attention au sophisme du type « $p \to q$ ; or $q$ ; donc $p$ » (affirmation du conséquent), qui est invalide.
\item \textbf{Pertinence} : chaque exemple illustre-t-il bien l'argument, et chaque argument sert-il bien la problématique ? Un hors-sujet est un raisonnement valable... qui répond à une autre question.
\end{enumerate}

\begin{aretenir}[La grille de relecture logique]
Un bon devoir est un devoir où l'on peut, à chaque paragraphe, répondre à trois questions : \emph{qu'affirmes-tu ? pourquoi l'affirmes-tu ? en quoi cela répond-il à la question posée ?}
\end{aretenir}

\courssub{Formaliser : traduire la langue courante en langage logique}

Formaliser, c'est traduire une phrase ordinaire en symboles pour tester ensuite sa validité de manière mécanique, sans se laisser piéger par les mots.

\begin{methode}[Méthode de formalisation]
\begin{enumerate}
\item \textbf{Repérer les propositions simples} : souligner les énoncés qui peuvent être vrais ou faux.
\item \textbf{Les symboliser} : attribuer à chacune une lettre ($p$, $q$, $r$...).
\item \textbf{Identifier les connecteurs} : « et » $\to \wedge$ ; « ou » $\to \vee$ ; « si... alors » $\to \to$ ; « non » $\to \neg$ ; « si et seulement si » $\to \leftrightarrow$.
\item \textbf{Assembler} la formule complète en respectant les parenthèses.
\end{enumerate}
\end{methode}

\begin{exemplebox}[Formalisation guidée]
Phrase : « Si tu révises \emph{et} que tu dors bien, \emph{alors} tu réussiras. »
\begin{itemize}
\item Propositions simples : $p$ = « tu révises » ; $q$ = « tu dors bien » ; $r$ = « tu réussiras ».
\item Connecteurs : « et » entre $p$ et $q$ ; « si... alors » entre le bloc $(p \wedge q)$ et $r$.
\item Formule : $(p \wedge q) \to r$.
\end{itemize}
Autre exemple : « Tu paies tes droits ou tu ne passes pas l'examen » se formalise $p \vee \neg q$, avec $p$ = « tu paies tes droits » et $q$ = « tu passes l'examen ».
\end{exemplebox}

\begin{exoresolu}[Formaliser puis tester un raisonnement]
Énoncé. Un élève raisonne ainsi : « Si je travaille régulièrement, je réussis au Baccalauréat technique. Or je travaille régulièrement. Donc je réussirai. » Formaliser ce raisonnement et dire s'il est valide.
\tcblower
Solution détaillée. Posons $p$ = « je travaille régulièrement » et $q$ = « je réussis au Baccalauréat technique ». Le raisonnement s'écrit : prémisse 1 : $p \to q$ ; prémisse 2 : $p$ ; conclusion : $q$. C'est exactement le schéma du \emph{modus ponens} : $[(p \to q) \wedge p] \to q$. La table de vérité de cette formule ne contient que des V (c'est une tautologie) : le raisonnement est donc \textbf{valide}. Attention : si l'élève avait dit « or je réussis, donc je travaille régulièrement », ce serait l'affirmation du conséquent, un raisonnement \textbf{invalide}.
\end{exoresolu}

\courssub{TP guidé 1 : table de vérité d'un raisonnement complexe}

\textbf{Objectif.} Construire la table de vérité complète du raisonnement à deux prémisses : « Si le candidat est admis, il se réjouit. Or le candidat est admis. Donc il se réjouit », soit $[(p \to q) \wedge p] \to q$, et conclure sur sa validité.

\begin{experience}[Protocole du TP]
\textbf{Matériel} : cahier d'exercices, règle, stylo. \textbf{Protocole} :
\begin{enumerate}
\item Identifier les propositions simples ($p$, $q$) : il y en a 2, donc $2^2 = 4$ lignes.
\item Remplir les colonnes de $p$ et $q$ avec toutes les combinaisons (V/V, V/F, F/V, F/F).
\item Calculer la colonne $p \to q$ (fausse uniquement quand $p$ est vrai et $q$ faux).
\item Calculer $(p \to q) \wedge p$.
\item Calculer la formule complète $[(p \to q) \wedge p] \to q$.
\item \textbf{Conclure} : si la dernière colonne ne contient que des V, le raisonnement est valide (tautologie).
\end{enumerate}
\end{experience}

\begin{popfigure}
\begin{center}
\begin{tabular}{|c|c|c|c|c|}
\hline
\rowcolor{popblueL} $p$ & $q$ & $p \to q$ & $(p \to q) \wedge p$ & $[(p \to q) \wedge p] \to q$ \\
\hline
V & V & V & V & \textbf{V} \\
\hline
V & F & F & F & \textbf{V} \\
\hline
F & V & V & F & \textbf{V} \\
\hline
F & F & V & F & \textbf{V} \\
\hline
\end{tabular}
\end{center}
\legende{Table de vérité du raisonnement $[(p \to q) \wedge p] \to q$ (modus ponens) : la dernière colonne ne contient que des V, donc le raisonnement est valide.}
\end{popfigure}

\textbf{Interprétation.} La dernière colonne est uniformément vraie : quelle que soit la situation, il est impossible que les prémisses soient vraies et la conclusion fausse. Le raisonnement est donc \cle{valide}. Ce même procédé permet de démasquer les sophismes : pour l'affirmation du conséquent $[(p \to q) \wedge q] \to p$, la ligne $p$ = F, $q$ = V donnerait F dans la dernière colonne, preuve d'invalidité.

\courssub{TP guidé 2 : rédiger et présenter une démonstration par l'absurde}

\textbf{Objectif.} Rédiger intégralement une démonstration par l'absurde et la présenter au tableau.

\begin{experience}[Démonstration guidée pas à pas]
\textbf{Thèse à démontrer} : « Celui qui ne remplit pas les conditions d'inscription ne peut pas être admis à composer. »
\begin{enumerate}
\item \textbf{Poser la négation de la thèse} (hypothèse de l'absurde) : supposons qu'un candidat ne remplit pas les conditions d'inscription \emph{et} soit admis à composer.
\item \textbf{Tirer les conséquences} : or le règlement de l'examen stipule que seuls les candidats régulièrement inscrits figurent sur les listes d'émargement ; notre candidat devrait donc figurer sur une liste où, par hypothèse, il ne peut pas figurer.
\item \textbf{Aboutir à une contradiction} : le candidat serait à la fois inscrit (puisqu'il compose) et non inscrit (par hypothèse) : $p \wedge \neg p$, ce qui est impossible.
\item \textbf{Rejeter l'hypothèse} : la supposition de départ est donc fausse.
\item \textbf{Conclure} : la thèse initiale est démontrée — qui ne remplit pas les conditions d'inscription ne peut pas composer.
\end{enumerate}
\textbf{Présentation au tableau} : un élève rédige les cinq étapes, la classe vérifie que la contradiction est bien explicite et que la conclusion reprend exactement la thèse annoncée.
\end{experience}

\begin{attention}[Piège de l'absurde]
L'hypothèse à poser est la \emph{négation exacte} de la thèse, pas une thèse différente. Et la contradiction obtenue doit être \emph{explicitée} : écrire simplement « c'est absurde » sans montrer le $p \wedge \neg p$ ne vaut pas démonstration.
\end{attention}

\courssub{Débat contradictoire encadré : « Faut-il rendre le port de l'uniforme obligatoire à l'université ? »}

\textbf{Organisation.} La classe est divisée en deux camps : le camp \emph{pour} et le camp \emph{contre}. Chaque camp prépare deux arguments et un exemple, puis anticipe les objections adverses (c'est l'esprit du plan dialectique appliqué à l'oral). Un arbitre (l'enseignant) veille au respect des règles logiques.

\textbf{Règles du débat.} (1) Toute affirmation doit être suivie d'un argument introduit par \emph{car} ou \emph{parce que}. (2) Toute objection doit viser l'argument, non la personne (interdiction du sophisme \emph{ad hominem}). (3) Chaque prise de parole se conclut par un connecteur de conséquence (\emph{donc}, \emph{par conséquent}).

\textbf{Exemples d'arguments attendus.} Pour : l'uniforme efface les inégalités sociales visibles entre étudiants et renforce le sentiment d'appartenance. Contre : l'université forme des adultes responsables, et la liberté de se vêtir fait partie de l'autonomie que l'Université doit respecter ; de plus, le coût de l'uniforme pèserait sur les familles modestes. La synthèse collective peut proposer une solution nuancée : un code vestimentaire minimal plutôt qu'un uniforme strict.

\begin{savaistu}[L'argumentation, clé de la réussite au Baccalauréat technique]
Au Baccalauréat technique comme au Probatoire, les barèmes de correction de la dissertation de philosophie accordent une part essentielle de la note à la qualité de l'argumentation et de la justification : pertinence de la problématique, solidité des arguments, exactitude des exemples, cohérence des enchaînements et clarté des transitions. Deux copies citant les mêmes auteurs peuvent obtenir des notes très différentes selon que le raisonnement est logiquement construit ou non. Travailler la logique, c'est donc travailler directement sa note.
\end{savaistu}

\begin{exercice}[Formalisation et validité]
\begin{enumerate}
\item Formaliser : « Si le gouvernement baisse les prix et augmente les salaires, alors le mécontentement diminuera. »
\item Construire la table de vérité de $[(p \to q) \wedge \neg q] \to \neg p$ et dire si ce raisonnement (modus tollens) est valide.
\item Repérer le sophisme : « Ce jeune n'a pas de diplôme, donc son avis sur la politique ne compte pas. »
\end{enumerate}
\end{exercice}

\begin{corrige}[Exercice]
\begin{enumerate}
\item $p$ = « le gouvernement baisse les prix », $q$ = « le gouvernement augmente les salaires », $r$ = « le mécontentement diminuera ». Formule : $(p \wedge q) \to r$.
\item Quatre lignes ; la formule $[(p \to q) \wedge \neg q] \to \neg p$ ne prend que la valeur V : le modus tollens est une tautologie, donc un raisonnement valide.
\item C'est un sophisme \emph{ad hominem} : on attaque la personne (son absence de diplôme) au lieu d'examiner la validité de son argument. Or la valeur d'un raisonnement ne dépend pas du statut de celui qui l'énonce.
\end{enumerate}
\end{corrige}

\begin{aretenir}[L'essentiel de la section]
La logique de l'unité se met au service de la copie et de la parole : problématique claire, plan dialectique, transitions explicites, paragraphes construits selon le schéma « annoncer, argumenter, illustrer, conclure partiellement » ; relecture attentive à l'absence de contradiction et à la validité des enchaînements ; formalisation des phrases courantes pour tester mécaniquement les raisonnements par les tables de vérité ; enfin, pratique encadrée du débat contradictoire et de la démonstration par l'absurde. Celui qui maîtrise ces outils sait \emph{rédiger et justifier une démarche, une réponse ou une conclusion} — exactement le savoir-faire exigé par le programme.
\end{aretenir}

\newpage

\coursec{Activité d'intégration}

\coursec{Logique, argumentation et démonstration}

\courssub{Objectifs de l'activité}
\begin{itemize}
    \item \textbf{Mobiliser} les opérations logiques fondamentales (négation, conjonction, disjonction, implication, équivalence) et les tables de vérité pour formaliser des arguments du langage courant.
    \item \textbf{Analyser} la structure d'un raisonnement déductif (syllogisme) pour en vérifier la validité formelle et la vérité des prémisses.
    \item \textbf{Identifier} et \textbf{nommer} les principaux sophismes (argument \emph{ad hominem}, homme de paille, faux dilemme, appel à l'émotion, \emph{petitio principii}, etc.) dans un discours authentique.
    \item \textbf{Compléter} une démonstration par l'absurde (réduction à l'absurde) en explicitant les étapes logiques manquantes.
    \item \textbf{Rédiger} une conclusion argumentée et justifiée (environ une demi-page) synthétisant l'analyse critique des trois documents.
    \item \textbf{Communiquer} oralement les résultats de l'analyse devant un jury (la classe) en respectant les normes de l'argumentation rationnelle.
\end{itemize}

\courssub{Situation}
\textbf{Contexte :} Dans le cadre de l'émission \emph{« Tribune Citoyenne »} sur la chaîne de radio privée \textbf{Radio Équateur} (Yaoundé), un débat animé a lieu sur l'état de la route nationale n°3 (Yaoundé–Douala), axe économique vital du Cameroun. Le professeur a constitué un corpus de trois documents authentiques (ou reconstitués fidèlement à partir d'archives sonores) soumis au « Tribunal de la logique » que forme la classe.

\bigskip
\noindent\textbf{Document 1 : Extrait du débat radiophonique (15 min 30 s -- 22 min 10 s)}
\begin{quote}
\textbf{Animateur :} « Chers auditeurs, la route Yaoundé–Douala fait encore la une. Deux poids lourds se sont encastrés hier au pk 120, bloquant le trafic six heures. Le ministre des Travaux publics a promis un plan d'urgence. Mais la population est sceptique. Au standard, M. \textbf{Fouda}, taxi-man au carrefour Nsam. Allô Fouda ? »

\textbf{Fouda (taximan, voix agitée) :} « Allô patron ! Moi je dis que ce gouvernement est incompétent, point final. \textbf{Regardez les ministres : ils roulent en V8 climatisés, ils ne sentent pas les nids-de-poule.} Eux, ils dorment dans des hôtels 5 étoiles pendant que nous, on casse nos colonnes vertébrales sur le bitume. \textbf{Si on écoute les ingénieurs du MINTP, ils vont nous noyer dans leurs formules chinoises.} \textbf{Soit on refait toute la route maintenant avec de l'argent qu'on n'a pas, soit on laisse pourrir.} Moi, je vote pour la refonte totale, même si faut vendre le pétrole ! »

\textbf{Animateur :} « Merci Fouda. Au standard, Mme \textbf{Essomba}, porte-parole du parti au pouvoir. Mme Essomba ? »

\textbf{Mme Essomba (calme, posée) :} « Merci. D'abord, \textbf{attaquer les personnes au lieu des arguments, c'est de la diversion.} Fouda parle de V8 : \textbf{lui-même conduit un taxi vétuste qui pollue l'air de nos enfants.} Ensuite, \textbf{le gouvernement a déjà bitumé 40 km sur l'axe depuis 2020.} Les chiffres du MINTP sont formels. \textbf{Nier cette réalité, c'est nier le travail des ingénieurs camerounais.} \textbf{Soit on soutient l'effort national, soit on s'allie aux destructeurs de la paix.} Nous choisissons la construction. »

\textbf{Animateur :} « Dernier intervenant, M. \textbf{Ngoua}, ingénieur des ponts et chaussées, retraité. M. Ngoua ? »

\textbf{M. Ngoua (posé, technique) :} « Bonjour. Laissons l'émotion. \textbf{Une route se conçoit selon des normes : trafic, géologie, climat.} Le tracé actuel traverse des zones marécageuses (pk 80--110) et des zones de roches dures (pk 120--150). \textbf{Si on applique la norme AASHTO 2020 sans étude géotechnique approfondie par tronçon, le revêtement craquera sous 18 mois.} \textbf{Donc, refaire "toute la route maintenant" sans études, c'est garantir l'échec.} \textbf{La solution : diagnostic par section, puis priorisation budgétaire pluriannuelle.} »
\end{quote}

\bigskip
\noindent\textbf{Document 2 : Syllogisme à trous (extrait d'un meeting électoral, Douala, 2023)}
\begin{quote}
« Camerounaises, Camerounais,
\textbf{Prémisse majeure :} \underline{\hspace{5cm}}.
\textbf{Prémisse mineure :} \underline{\hspace{5cm}}.
\textbf{Conclusion :} \textbf{Ainsi, le candidat \emph{Kamga} est le seul capable de résoudre le chômage des jeunes dans le Littoral.} »
\end{quote}
\emph{Indices fournis par le professeur :} La prémisse majeure est une universelle affirmative (type A). La prémisse mineure est une particulière affirmative (type I). La conclusion est une universelle affirmative (type A). Le terme moyen est « \emph{gestionnaire rigoureux} ». Le sujet de la conclusion est « \emph{candidat Kamga} ». L'attribut de la conclusion est « \emph{capable de résoudre le chômage des jeunes dans le Littoral} ».

\bigskip
\noindent\textbf{Document 3 : Raisonnement mathématique à compléter (Démonstration par l'absurde)}
\begin{quote}
\textbf{Proposition :} « La somme d'un nombre rationnel et d'un nombre irrationnel est toujours irrationnelle. »

\textbf{Début de démonstration :}
Soit $r \in \mathbb{Q}$ un nombre rationnel et $x \in \mathbb{R} \setminus \mathbb{Q}$ un nombre irrationnel.
On suppose, \textbf{par l'absurde}, que la somme $s = r + x$ est \underline{\hspace{3cm}}.
Or, par définition, $r = \frac{a}{b}$ avec $a,b \in \mathbb{Z}, b \neq 0$.
Si $s$ est rationnel, alors $s = \frac{c}{d}$ avec $c,d \in \mathbb{Z}, d \neq 0$.
On en déduit que $x = s - r = \frac{c}{d} - \frac{a}{b} = \underline{\hspace{3cm}}$.
Or, $\frac{cb - ad}{bd}$ est un quotient d'entiers (dénominateur non nul), donc \underline{\hspace{3cm}}.
Cela contredit l'hypothèse que $x$ est \underline{\hspace{3cm}}.
\underline{\hspace{3cm}}, la somme $r+x$ est bien irrationnelle. \hfill $\square$
\end{quote}

\courssub{Consignes}
\textbf{Travail en groupes de 4 à 5 élèves (30 min de préparation, 10 min de présentation par groupe).}

\textbf{Phase 1 : Formalisation et Tables de vérité (Document 1)}
\begin{enumerate}
    \item \textbf{1.} Extraire trois arguments distincts tenus par \textbf{Fouda} et trois par \textbf{Mme Essomba}. Les formaliser en logique propositionnelle en définissant des variables propositionnelles claires (ex. $p$ : « Le gouvernement est incompétent », $q$ : « Les ministres roulent en V8 »...).
    \item \textbf{2.} Pour chaque argument formalisé, construire la table de vérité correspondante (en ne retenant que les connecteurs principaux : $\land, \lor, \Rightarrow, \Leftrightarrow, \neg$).
    \item \textbf{3.} Déterminer si la forme logique est une tautologie (valide), une contradiction ou une contingence.
\end{enumerate}

\textbf{Phase 2 : Chasse aux sophismes (Document 1)}
\begin{enumerate}
    \item \textbf{4.} Identifier \textbf{au moins trois sophismes distincts} dans le Document 1 (répartis entre les intervenants).
    \item \textbf{5.} Pour chaque sophisme : citer la phrase exacte, nommer le sophisme (ex. \emph{Argumentum ad hominem}, \emph{Homme de paille}, \emph{Faux dilemme}, \emph{Appel à l'émotion}, \emph{Appel à la force}, \emph{Pétition de principe}), expliquer \textbf{en quoi} le raisonnement est fallacieux (rupture de la validité ou prémisse inacceptable).
\end{enumerate}

\textbf{Phase 3 : Reconstruction du syllogisme (Document 2)}
\begin{enumerate}
    \item \textbf{6.} Compléter les deux prémisses manquantes en respectant la structure standard (Majeure, Mineure, Conclusion), les quantificateurs (Tout, Certains, Aucun) et les termes fournis (Terme moyen : \emph{gestionnaire rigoureux} ; Sujet : \emph{candidat Kamga} ; Attribut : \emph{capable de résoudre le chômage...}).
    \item \textbf{7.} Identifier la \textbf{figure} et le \textbf{mode} du syllogisme obtenu (ex. AAA-1 / Barbara).
    \item \textbf{8.} Vérifier la \textbf{validité formelle} (règles de la quantité, qualité, terme moyen, distribution). Le syllogisme est-il valide ?
    \item \textbf{9.} Discuter de la \textbf{validité matérielle} (vérité des prémisses dans la réalité). La conclusion est-elle nécessairement vraie ?
\end{enumerate}

\textbf{Phase 4 : Complétion de la démonstration (Document 3)}
\begin{enumerate}
    \item \textbf{10.} Compléter les cinq trous de la démonstration par l'absurde.
    \item \textbf{11.} Justifier chaque étape par une règle de calcul (propriétés de $\mathbb{Z}$, définition de $\mathbb{Q}$) ou une règle logique (loi de non-contradiction, \emph{modus tollens}).
    \item \textbf{12.} Expliquer pourquoi cette méthode s'appelle « réduction à l'absurde » (\emph{reductio ad absurdum}).
\end{enumerate}

\textbf{Phase 5 : Synthèse et Plaidoirie (Production finale)}
\begin{enumerate}
    \item \textbf{13.} Rédiger collectivement une \textbf{conclusion justifiée (250--300 mots)} structurée ainsi :
    \begin{itemize}
        \item \textbf{Verdict global :} Quel intervenant du Doc 1 propose le raisonnement le plus rigoureux ? Pourquoi ?
        \item \textbf{Analyse du Doc 2 :} Le syllogisme électoral est-il un outil de conviction rationnelle ou de persuasion manipulatoire ?
        \item \textbf{Lien Doc 3 :} En quoi la rigueur de la démonstration mathématique (Doc 3) devrait-elle inspirer le débat public (Doc 1) ?
    \end{itemize}
    \item \textbf{14.} Désigner un \emph{avocat} par groupe pour présenter le verdict devant le « Jury » (la classe + professeur) en 3 minutes chrono. Le jury note sur une grille : \emph{Rigueur logique (5 pts), Clarté de l'oral (3 pts), Pertinence des exemples (2 pts)}.
\end{enumerate}

\courssub{Ressources à mobiliser}
\begin{definition}[Tables de vérité des connecteurs fondamentaux]
Pour deux propositions $p$ et $q$ :
\begin{center}
\begin{tabular}{|c|c|c|c|c|c|c|}\hline
$p$ & $q$ & $\neg p$ & $p \land q$ & $p \lor q$ & $p \Rightarrow q$ & $p \Leftrightarrow q$ \\ \hline
V & V & F & V & V & V & V \\
V & F & F & F & V & F & F \\
F & V & V & F & V & V & F \\
F & F & V & F & F & V & V \\ \hline
\end{tabular}
\end{center}
\emph{Rappel :} $p \Rightarrow q$ est faux \textbf{uniquement} si $p$ est Vrai et $q$ est Faux.
\end{definition}

\begin{propriete}[Règles du syllogisme catégorique (Validité formelle)]
\begin{enumerate}
    \item \textbf{Terme moyen :} Doit être distribué au moins une fois.
    \item \textbf{Distribution des termes :} Un terme non distribué dans les prémisses ne peut pas l'être dans la conclusion.
    \item \textbf{Prémisses négatives :} Aucune conclusion ne suit de deux prémisses négatives. Si une prémisse est négative, la conclusion doit l'être.
    \item \textbf{Prémisses particulières :} Aucune conclusion universelle ne suit de deux prémisses particulières. Si une prémisse est particulière, la conclusion doit l'être.
    \item \textbf{Figure :} Déterminée par la position du terme moyen (M) : Fig 1 (MP/SM), Fig 2 (PM/SM), Fig 3 (MP/MS), Fig 4 (PM/MS).
    \item \textbf{Mode :} Suite des quantificateurs (A=Universelle Affirmative, E=Universelle Négative, I=Particulaire Affirmative, O=Particulaire Négative).
\end{enumerate}
\textbf{Modes valides classiques (mnésiques) :} Barbara (AAA-1), Celarent (EAE-1), Darii (AII-1), Ferio (EIO-1), Cesare (EAE-2), Camestres (AEE-2), Festino (EIO-2), Baroco (AOO-2), Darapti (AAI-3), Disamis (IAI-3), Datisi (AII-3), Felapton (EAO-3), Bocardo (OAO-3), Ferison (EIO-3), Bramantip (AAI-4), Camenes (AEE-4), Dimaris (IAI-4), Fesapo (EAO-4), Fresison (EIO-4).
\end{propriete}

\begin{aretenir}[Principaux sophismes (liste non exhaustive pour le Bac)]
\begin{itemize}
    \item \textbf{Ad hominem (abusif) :} Attaquer la personne au lieu de son argument.
    \item \textbf{Homme de paille (Straw man) :} Déformer l'argument adverse pour le réfuter plus facilement.
    \item \textbf{Faux dilemme (Fausse alternative) :} Présenter deux extrêmes comme les seules options possibles.
    \item \textbf{Appel à l'émotion (Ad passiones) :} Remplacer la preuve par la sollicitation de la pitié, peur, colère, fierté.
    \item \textbf{Appel à l'autorité (Ad verecundiam) :} Invoquer une autorité non compétente en la matière ou non consensuelle.
    \item \textbf{Pétition de principe (Cercle vicieux) :} La conclusion est déjà contenue dans les prémisses.
    \item \textbf{Non sequitur :} La conclusion ne découle pas logiquement des prémisses.
    \item \textbf{Glissement de pente (Slippery slope) :} Une action mineure entraînerait inéluctablement une catastrophe.
\end{itemize}
\end{aretenir}

\begin{methode}[Structure de la démonstration par l'absurde (Réduction à l'absurde)]
Pour démontrer une proposition $P$ :
\begin{enumerate}
    \item Supposer $\neg P$ (hypothèse adverse / hypothèse de l'absurde).
    \item Enchaîner des déductions logiques valides à partir de $\neg P$ et des axiomes/définitions connus.
    \item Aboutir à une \textbf{contradiction} (une proposition $Q$ et sa négation $\neg Q$ sont vraies simultanément, ou violation d'un axiome, ou absurdité manifeste comme $1=0$).
    \item Conclure : « L'hypothèse $\neg P$ conduit à l'absurde, donc $\neg P$ est fausse, donc $P$ est vraie. » (Loi du tiers exclu : $P \lor \neg P$ ; $\neg(\neg P) \vdash P$).
\end{enumerate}
\end{methode}

\courssub{Démarche et corrigé commenté}
\begin{exoresolu}[Corrigé détaillé du TP « Tribunal de la logique »]
Énoncé : Voir la situation, les trois documents et les consignes (Phases 1 à 5) ci-dessus.
\tcblower
\textbf{Phase 1 -- Formalisation et Tables de vérité (Document 1)}

\medskip
\noindent\textbf{Arguments de Fouda (Taximan) :}
\begin{enumerate}
    \item \emph{« Ce gouvernement est incompétent car les ministres roulent en V8. »}
    \begin{itemize}
        \item $p$ : « Le gouvernement est incompétent. » ; $q$ : « Les ministres roulent en V8. »
        \item Formalisation : $q \Rightarrow p$ (Implication : Si V8 alors incompétent).
        \item Table de vérité : Ligne 2 ($q=V, p=F$) rend l'implication Fausse. \textbf{Contingence.} La forme n'est pas tautologique ; le lien de causalité est gratuit.
    \end{itemize}
    \item \emph{« Soit on refait toute la route maintenant, soit on laisse pourrir. »}
    \begin{itemize}
        \item $r$ : « On refait toute la route maintenant. » ; $s$ : « On laisse pourrir. »
        \item Formalisation : $r \lor s$ (Disjonction exclusive implicite, présentée comme exhaustive).
        \item Table de vérité : Vrai si $r=V$ ou $s=V$. \textbf{Contingence.} La structure logique est valide (disjonction), mais le contenu est fallacieux (exclusion du tiers).
    \end{itemize}
    \item \emph{« Si on écoute les ingénieurs, ils vont nous noyer dans leurs formules chinoises. »}
    \begin{itemize}
        \item $u$ : « On écoute les ingénieurs. » ; $v$ : « On est noyé dans des formules chinoises. »
        \item Formalisation : $u \Rightarrow v$.
        \item \textbf{Contingence.}
    \end{itemize}
\end{enumerate}

\noindent\textbf{Arguments de Mme Essomba (Porte-parole) :}
\begin{enumerate}
    \item \emph{« Le gouvernement a bitumé 40 km depuis 2020. Nier cette réalité, c'est nier le travail des ingénieurs. »}
    \begin{itemize}
        \item $a$ : « Le gouvernement a bitumé 40 km. » ; $b$ : « On nie le travail des ingénieurs. »
        \item Formalisation : $\neg a \Rightarrow b$ (Si on nie le chiffre, alors on nie le travail).
        \item \textbf{Contingence.}
    \end{itemize}
    \item \emph{« Soit on soutient l'effort national, soit on s'allie aux destructeurs de la paix. »}
    \begin{itemize}
        \item $c$ : « On soutient l'effort national. » ; $d$ : « On s'allie aux destructeurs de la paix. »
        \item Formalisation : $c \lor d$ (Disjonction exclusive présentée comme exhaustive).
        \item \textbf{Contingence.}
    \end{itemize}
    \item \emph{« Fouda conduit un taxi vétuste, donc son argument sur les V8 est invalide. »} (Implicite).
    \begin{itemize}
        \item $e$ : « Le taxi de Fouda est vétuste. » ; $f$ : « L'argument de Fouda est invalide. »
        \item Formalisation : $e \Rightarrow f$.
        \item \textbf{Contingence.}
    \end{itemize}
\end{enumerate}

\medskip
\noindent\textbf{Phase 2 -- Identification des sophismes (Document 1)}

\begin{enumerate}
    \item \textbf{Sophisme 1 : Argument \emph{ad hominem} (abusif) + \emph{Tu quoque}.}
    \begin{itemize}
        \item \textbf{Auteur :} Mme Essomba.
        \item \textbf{Citation :} « \emph{Fouda parle de V8 : lui-même conduit un taxi vétuste qui pollue l'air de nos enfants.} »
        \item \textbf{Analyse :} Au lieu de répondre sur le fond (l'incompétence alléguée liée aux V8), elle attaque la personne de Fouda (son taxi vétuste) et retourne l'accusation (\emph{tu quoque}). La validité de l'argument de Fouda sur les V8 ministériels est indépendante de l'état du taxi de Fouda. Rupture de la pertinence.
    \end{itemize}
    \item \textbf{Sophisme 2 : Faux dilemme (Fausse alternative / \emph{Bifurcation}).}
    \begin{itemize}
        \item \textbf{Auteurs :} Fouda ET Mme Essomba (symétrie frappante).
        \item \textbf{Citations :} Fouda : « \emph{Soit on refait toute la route maintenant avec de l'argent qu'on n'a pas, soit on laisse pourrir.} » --- Mme Essomba : « \emph{Soit on soutient l'effort national, soit on s'allie aux destructeurs de la paix.} »
        \item \textbf{Analyse :} Les deux présentent une disjonction exclusive ($A \lor B$) comme exhaustive, excluant les solutions intermédiaires (ex. : réhabilitation par tronçons prioritaires, financement public-privé, audit technique préalable). C'est une manipulation rhétorique pour forcer le choix.
    \end{itemize}
    \item \textbf{Sophisme 3 : Homme de paille (\emph{Straw man}) + Appel à l'émotion / au peuple.}
    \begin{itemize}
        \item \textbf{Auteur :} Fouda.
        \item \textbf{Citation :} « \emph{Si on écoute les ingénieurs du MINTP, ils vont nous noyer dans leurs formules chinoises.} »
        \item \textbf{Analyse :} Fouda caricature le discours technique (complexe, rigoureux) en « formules chinoises » (incompréhensibles, étrangères, menaçantes) pour le rejeter sans l'examiner. Il substitue un épouvantail (l'obscurantisme technique) à l'argument réel (la nécessité d'études géotechniques). L'expression « noyer » appelle la peur/émotion.
    \end{itemize}
    \item \textbf{Sophisme 4 (Bonus) : Appel à l'autorité / au patriotisme détourné.}
    \begin{itemize}
        \item \textbf{Auteur :} Mme Essomba.
        \item \textbf{Citation :} « \emph{Nier cette réalité, c'est nier le travail des ingénieurs camerounais.} »
        \item \textbf{Analyse :} Elle assimile la critique du \emph{bilan global} (40 km sur 300+ km) à une insulte à la corporation des ingénieurs. Détournement du sujet.
    \end{itemize}
\end{enumerate}

\medskip
\noindent\textbf{Phase 3 -- Reconstruction du syllogisme (Document 2)}

\noindent\textbf{Données :}
\begin{itemize}
    \item Terme Majeur (P - Attribut de la conclusion) : \emph{Capable de résoudre le chômage des jeunes dans le Littoral}.
    \item Terme Mineur (S - Sujet de la conclusion) : \emph{Candidat Kamga}.
    \item Terme Moyen (M) : \emph{Gestionnaire rigoureux}.
    \item Types : Majeure (A), Mineure (I), Conclusion (A).
\end{itemize}

\noindent\textbf{Reconstruction :}
\begin{itemize}
    \item \textbf{Prémisse Majeure (A - Universelle Affirmative) :} \textbf{Tout gestionnaire rigoureux est capable de résoudre le chômage des jeunes dans le Littoral.} (M $\subset$ P)
    \item \textbf{Prémisse Mineure (I - Particulaire Affirmative) :} \textbf{Certains candidats (dont le candidat Kamga) sont des gestionnaires rigoureux.} (S $\cap$ M $\neq \emptyset$). \emph{Note : L'énoncé impose une mineure I. On ne peut pas écrire "Kamga est un gestionnaire rigoureux" (singulier $\equiv$ A) sans violer la consigne de type I.}
    \item \textbf{Conclusion (A - Universelle Affirmative) :} \textbf{Tout candidat Kamga est capable de résoudre le chômage des jeunes dans le Littoral.}
\end{itemize}

\noindent\textbf{Analyse formelle :}
\begin{itemize}
    \item \textbf{Mode/Figure :} Figure 1 (Majeure M-P, Mineure S-M). Mode \textbf{AII-1} pour les prémisses, mais Conclusion A. \textbf{Mode inexistant parmi les valides.} Si on force Conclusion A, il faudrait Mineure A (AAA-1 Barbara), ce qui contredit l'énoncé.
    \item \textbf{Validité formelle :} \textbf{INVALIDE.} Double violation :
    \begin{enumerate}
        \item Règle 4 : « Aucune conclusion universelle ne suit d'une prémisse particulière ». La mineure est I (particulière), la conclusion est A (universelle).
        \item Règle 2 (Distribution) : Le terme mineur (Kamga / S) n'est pas distribué dans la mineure (I ne distribue pas le sujet), mais il est distribué dans la conclusion (A distribue le sujet).
    \end{enumerate}
    \item \textbf{Validité matérielle :} Même si on corrige la forme en Barbara (Majeure A, Mineure A, Conclusion A), la prémisse majeure « Tout gestionnaire rigoureux résout le chômage » est \textbf{fausse} (condition nécessaire $\neq$ suffisante ; facteurs macroéconomiques externes, conjoncture internationale). La conclusion n'est pas établie.
\end{itemize}

\medskip
\noindent\textbf{Phase 4 -- Complétion de la démonstration (Document 3)}

\noindent\textbf{Proposition :} $r \in \mathbb{Q}, x \notin \mathbb{Q} \implies r+x \notin \mathbb{Q}$.

\begin{enumerate}
    \item \textbf{Trou 1 (Hypothèse de l'absurde) :} \textbf{rationnelle} (ou $\in \mathbb{Q}$).
    \item \textbf{Trou 2 (Calcul de $x$) :} $\frac{c}{d} - \frac{a}{b} = \frac{cb - ad}{bd}$.
    \item \textbf{Trou 3 (Nature du résultat) :} \textbf{un nombre rationnel} (ou $\in \mathbb{Q}$). \emph{Justification :} $cb, ad, bd \in \mathbb{Z}$ (stabilité de $\mathbb{Z}$ par $+,\times$), $bd \neq 0$. Quotient d'entiers $\Rightarrow \mathbb{Q}$.
    \item \textbf{Trou 4 (Contradiction) :} \textbf{irrationnel} (ou $\notin \mathbb{Q}$). \emph{Justification :} Hypothèse de départ sur $x$ ($x \in \mathbb{R} \setminus \mathbb{Q}$).
    \item \textbf{Trou 5 (Conclusion) :} \textbf{Par conséquent / Donc / CQFD}.
\end{enumerate}

\noindent\textbf{Justification logique de la méthode :}
La structure est : $[\neg P \vdash (Q \land \neg Q)] \vdash P$.
Ici $P \equiv (r+x \notin \mathbb{Q})$. $\neg P \equiv (r+x \in \mathbb{Q})$.
De $\neg P$ on déduit $x \in \mathbb{Q}$ (Trou 3). Or Axiome : $x \notin \mathbb{Q}$.
On a $(x \in \mathbb{Q}) \land (x \notin \mathbb{Q})$, contradiction (Principe de non-contradiction).
Donc $\neg P$ est faux. Par loi du tiers exclu ($P \lor \neg P$) et \emph{modus tollendo ponens} (ou élimination de la disjonction), $P$ est vrai.

\medskip
\noindent\textbf{Phase 5 -- Exemple de conclusion justifiée (Synthèse attendue)}

\emph{Verdict du Tribunal de la Logique}

Au terme de l'instruction, le Tribunal déclare que le discours de \textbf{M. Ngoua (Ingénieur)} est le seul à satisfaire aux exigences de la \textbf{raison critique}. Son argumentation repose sur des prémisses techniques vérifiables (normes AASHTO, géologie des sols) et une structure déductive valide : \emph{Si normes non respectées $\Rightarrow$ échec ; Or proposition Fouda/Essomba ignore normes $\Rightarrow$ Donc proposition vouée à l'échec.} Il distingue le \emph{fait technique} de l'\emph{opinion politique}.

À l'inverse, le débat politique (Fouda / Mme Essomba) sombre dans la \textbf{rhétorique eristique} : faux dilemmes symétriques (refaire tout / rien ; soutenir / trahir), \emph{ad hominem} réciproques, hommes de paille. Ils confondent \textbf{persuasion} (emporter l'adhésion par l'affect, l'autorité, l'exclusion) et \textbf{conviction} (imposer l'assentiment par la preuve logique).

Le syllogisme électoral (Doc 2) illustre cette dérive : bien que mimant la forme déductive (Barbara tronquée), il est \textbf{formellement invalide} (mineure particulière $\nrightarrow$ conclusion universelle) et \textbf{matériellement fallacieux} (majeure fausse : rigueur gestionnaire $\not\Rightarrow$ résolution chômage structurel). C'est un \emph{argument d'autorité} déguisé en logique.

Enfin, la démonstration mathématique (Doc 3) offre le \textbf{modèle épistémologique} absent du débat public : la \emph{reductio ad absurdum} impose une rigueur où chaque pas est nécessaire, où la contradiction sanctionne l'erreur sans appel à l'émotion. Si le débat sur la route Yaoundé-Douala suivait cette exigence — poser les définitions (qu'est-ce qu'une route durable ?), énoncer les axiomes (budget, géologie), déduire les conséquences, tester par l'absurde les options « tout béton » ou « rien » — il sortirait du stade de la querelle de clocher pour entrer dans celui de la \textbf{délibération rationnelle}, seul garant de l'intérêt général dans une République technique.

\hfill \textbf{Le Greffier du Tribunal}
\end{exoresolu}

\courssub{Bilan de l'activité}
Cette activité d'intégration a permis de mettre en évidence les acquis suivants, conformément aux attendus du programme de Terminale Technique :

\begin{itemize}
    \item \textbf{Maîtrise opératoire de la logique propositionnelle :} Les élèves sont capables de traduire un discours naturel en langage formel ($p, q, \Rightarrow, \lor$), de construire des tables de vérité et d'en lire la valeur de vérité (tautologie, contingence), distinguant ainsi la \textbf{forme} (valide/invalide) du \textbf{contenu} (vrai/faux).
    \item \textbf{Esprit critique affûté (Détection des sophismes) :} La contextualisation camerounaise (route Yaoundé-Douala, meeting électoral) a rendu concrète l'identification des sophismes les plus fréquents dans l'espace public (\emph{ad hominem}, faux dilemme, homme de paille, appel à l'émotion). Les élèves comprennent qu'un sophisme n'est pas une « erreur de calcul » mais une \textbf{rupture de la rationalité communicationnelle}.
    \item \textbf{Rigueur du raisonnement déductif :} L'exercice du syllogisme « piégé » (mineure I / conclusion A) a forcé l'application stricte des règles de distribution et de quantité, montrant que la \textbf{validité formelle} est une condition nécessaire (mais non suffisante) de la vérité. Le passage de la forme au fond (validité matérielle) a été opéré.
    \item \textbf{Compréhension de la démonstration mathématique comme paradigme :} La complétion de la preuve par l'absurde (somme rationnel+irrationnel) a ancré la structure \emph{Supposer $\neg P \rightarrow$ Contradiction $\rightarrow$ Donc $P$} et l'usage des définitions ensemblistes ($\mathbb{Q}, \mathbb{Z}$). Les élèves perçoivent la démonstration comme une \textbf{chaîne de nécessités logiques}, opposée à l'argumentation probabiliste ou persuasive.
    \item \textbf{Compétence communicationnelle (Écrire et Parler) :} La rédaction de la synthèse (demi-page) et la soutenance orale devant jury ont mobilisé les savoir-faire « Rédiger et justifier une démarche » et « Appliquer les notions à des situations concrètes ». L'évaluation par les pairs (grille jury) a renforcé l'exigence de clarté, de structure et de justesse terminologique.
\end{itemize}

\noindent\textbf{Poursuites possibles :} Analyse d'un article de presse camerounais (ex. \emph{Mutations}, \emph{Cameroon Tribune}) pour repérer les sophismes ; rédaction d'une dissertation type Bac : « La logique est-elle le seul garde-fou de la démocratie ? » ; initiation à la logique des prédicats (quantificateurs $\forall, \exists$) pour lever l'ambiguïté des syllogismes singuliers.

\newpage

\coursec{Méthodes et exercices résolus}

\coursec{Logique, argumentation et démonstration}

\courssub{La logique : nature et enjeux}

\begin{definition}[Logique]
La \cle{logique} (du grec \emph{logos} : raison, discours, science) est la discipline qui étudie les règles du raisonnement correct. Elle ne s'intéresse pas au contenu des pensées (la psychologie), mais à leur \emph{forme} : comment passer de propositions admises (prémisses) à une nouvelle proposition (conclusion) avec certitude.
\end{definition}

\begin{aretenir}[Deux acceptions de la logique]
\begin{itemize}
\item \textbf{Logique formelle (ou mineure)} : étude des structures valides du raisonnement (syllogisme, calcul propositionnel), abstraction faite du contenu. Elle garantit la \emph{validité} (si les prémisses sont vraies, la conclusion l'est nécessairement).
\item \textbf{Logique matérielle (ou majeure)} : étude des conditions de vérité des prémisses (définitions, divisions, preuves empiriques). Elle vise la \emph{vérité} du contenu.
\end{itemize}
\end{aretenir}

\begin{propriete}[Validité vs Vérité]
Un raisonnement est \textbf{valide} si sa conclusion découle \emph{nécessairement} de ses prémisses (forme correcte). Il est \textbf{solide} (ou \emph{sound}) s'il est valide \emph{et} si ses prémisses sont vraies. \textbf{Attention} : une conclusion vraie peut sortir d'un raisonnement invalide (ex. sophisme) ; une conclusion fausse peut sortir d'un raisonnement valide à partir de prémisses fausses.
\end{propriete}

\begin{exemplebox}[Validité sans vérité]
\begin{quote}
« Tous les oiseaux sont des mammifères ; or la poule est un oiseau ; donc la poule est un mammifère. »
\end{quote}
La forme est valide (syllogisme Barbara), mais la prémisse majeure est fausse $\rightarrow$ conclusion fausse. La logique formelle ne détecte pas l'erreur factuelle.
\end{exemplebox}

\begin{savaistu}[Histoire : L'Organon d'Aristote]
Aristote (384--322 av. J.-C.) est le fondateur de la logique formelle. Ses six traités réunis sous le nom d'\emph{Organon} (« instrument ») codifient le syllogisme et les catégories. Au Cameroun, l'enseignement de la logique aristotélicienne a structuré la philosophie au lycée depuis l'indépendance ; elle reste la base du programme de Terminale technique pour former l'esprit critique.
\end{savaistu}

\courssub{Les opérations logiques fondamentales}

\begin{definition}[Les trois opérations de l'esprit (Port-Royal / Aristote)]
\begin{enumerate}
\item \textbf{La simple appréhension (Concept)} : saisir l'essence d'une chose. Opérations : \cle{définition} (genre + différence spécifique) et \cle{division} (répartition de l'extension en espèces disjointes et exhaustives).
\item \textbf{Le jugement (Proposition)} : affirmer ou nier quelque chose de quelque chose. Structure : Sujet -- Copule -- Attribut. Caractères : \emph{quantité} (universelle / particulière) et \emph{qualité} (affirmative / négative) $\rightarrow$ 4 formes canoniques : A, E, I, O.
\item \textbf{Le raisonnement (Inférence)} : passer du connu à l'inconnu. \emph{Immédiat} (conversion, opposition) ou \emph{médiat} (syllogisme, induction).
\end{enumerate}
\end{definition}

\begin{propriete}[Carré des oppositions logiques]
Pour deux propositions de même sujet et même attribut :
\begin{center}
\begin{tabular}{|c|c|c|c|c|}\hline
 & \textbf{A (Univ. Aff.)} & \textbf{E (Univ. Nég.)} & \textbf{I (Part. Aff.)} & \textbf{O (Part. Nég.)} \\ \hline
\textbf{Contraires} & \emph{Tout S est P} & \emph{Aucun S n'est P} &  &  \\ \hline
\textbf{Subalternes} & $\downarrow$ &  & $\uparrow$ &  \\ \hline
\textbf{Subcontraires} &  &  & \emph{Certains S sont P} & \emph{Certains S ne sont pas P} \\ \hline
\textbf{Contradictoires} & $\updownarrow$ &  &  & $\updownarrow$ \\ \hline
\end{tabular}
\end{center}
\textbf{Règles} : Contraires = ne peuvent être vraies ensemble (peuvent être fausses ensemble). Subcontraires = ne peuvent être fausses ensemble (peuvent être vraies ensemble). Subalternation : Vrai $\rightarrow$ Vrai ; Faux $\rightarrow$ Faux (sens inverse). Contradictoires = l'une vraie, l'autre fausse (exclu du tiers).
\end{propriete}

\begin{popfigure}
\begin{tikzpicture}[scale=0.9, every node/.style={font=\small}]
% Square of Opposition Diagram
\node[draw, rounded corners, fill=popblueL, minimum width=3.5cm, minimum height=1.2cm, align=center] (A) at (0,2){\textbf{A}\\\emph{Tout S est P}};
\node[draw, rounded corners, fill=poporangeL, minimum width=3.5cm, minimum height=1.2cm, align=center] (E) at (5,2){\textbf{E}\\\emph{Aucun S n'est P}};
\node[draw, rounded corners, fill=popgreenL, minimum width=3.5cm, minimum height=1.2cm, align=center] (I) at (0,-1){\textbf{I}\\\emph{Certains S sont P}};
\node[draw, rounded corners, fill=poppinkL, minimum width=3.5cm, minimum height=1.2cm, align=center] (O) at (5,-1){\textbf{O}\\\emph{Certains S ne sont pas P}};

\draw[popblue, thick, <->] (A) -- node[midway, above, font=\tiny] {Contraires} (E);
\draw[popgreen, thick, <->] (I) -- node[midway, below, font=\tiny] {Subcontraires} (O);
\draw[popdark, thick, <->] (A) -- node[midway, left, font=\tiny] {Subalternation} (I);
\draw[popdark, thick, <->] (E) -- node[midway, right, font=\tiny] {Subalternation} (O);
\draw[poppurple, thick, dashed, <->] (A) -- node[midway, above left, font=\tiny, rotate=-45] {Contradictoires} (O);
\draw[poppurple, thick, dashed, <->] (E) -- node[midway, above right, font=\tiny, rotate=45] {Contradictoires} (I);
\end{tikzpicture}
\legende{Carré des oppositions logiques (A, E, I, O). Flèches pleines : relations classiques ; flèches pointillées : contradiction.}
\end{popfigure}

\begin{exemplebox}[Définition et Division]
\textbf{Définition} : « Le triangle est un polygone à trois côtés. » Genre = polygone ; Différence spécifique = à trois côtés.
\textbf{Division dichotomique} des êtres vivants : Vivants $\rightarrow$ (Végétaux / Animaux) ; Animaux $\rightarrow$ (Invertébrés / Vertébrés) ; Vertébrés $\rightarrow$ (Poissons, Amphibiens, Reptiles, Oiseaux, Mammifères). Exhaustif et disjoint.
\end{exemplebox}

\courssub{Le raisonnement déductif et le syllogisme}

\begin{methode}[Vérifier la validité d'un raisonnement déductif]
Pour juger un raisonnement déductif (syllogisme), procéder en étapes :
\begin{enumerate}
\item \textbf{Identifier la structure} : repérer les deux prémisses et la conclusion (souvent introduite par « donc », « par conséquent », « ainsi »).
\item \textbf{Repérer les trois termes} : \emph{Majeur} (P, attribut de la conclusion), \emph{Mineur} (S, sujet de la conclusion), \emph{Moyen} (M, présent dans les deux prémisses, absent de la conclusion).
\item \textbf{Vérifier la forme (règles du syllogisme)} :
  \begin{itemize}
  \item Trois termes exactement (pas d'équivoque).
  \item Le moyen terme doit être \emph{distribué} (universel) au moins une fois.
  \item Si une prémisse est négative, la conclusion est négative ; si la conclusion est négative, une prémisse est négative.
  \item Aucune conclusion affirmative de deux prémisses négatives.
  \item Le terme distribué dans la conclusion doit l'être dans la prémisse.
  \end{itemize}
\item \textbf{Distinguer validité et vérité} : la logique juge la forme, pas le contenu.
\item \textbf{Conclure en justifiant} : valide / invalide ; nommer le défaut (moyen terme non distribué, quaternius terminorum, conclusion trop large, etc.).
\end{enumerate}
\textbf{Réflexe} : reformuler sous la forme canonique « Tout M est P ; Tout S est M ; Donc Tout S est P » (Mode Barbara, 1\textsuperscript{re} figure) pour tester la structure.
\end{methode}

\begin{popfigure}
\begin{tikzpicture}[scale=1.1, every node/.style={font=\small}]
% Euler Diagram for Socrates Syllogism (Invalid)
\begin{scope}[xshift=-4cm]
\node[draw, circle, minimum width=4cm, fill=popblue!10] (H) at (0,0) {};
\node[draw, circle, minimum width=2cm, fill=poporange!20] (P) at (0,0) {};
\node at (0,0.8) {\textbf{Hommes}};
\node at (0,-0.5) {\textbf{Philosophes}};
\node[popdark] at (1.5, 1.2) {Socrate ?};
\draw[->, popdark, thick] (1.5,1) -- (0.5, 0.3);
\node[align=center, font=\tiny] at (-3, -1.5) {Prémisse 1 : Tout P est H\\Prémisse 2 : S est H\\Conclusion : S est P\\$\rightarrow$ \textcolor{poporange}{INVALIDE}};
\end{scope}

% Euler Diagram for Valid Syllogism (Barbara)
\begin{scope}[xshift=4cm]
\node[draw, circle, minimum width=4cm, fill=popblue!10] (H2) at (0,0) {};
\node[draw, circle, minimum width=2cm, fill=poporange!20] (P2) at (0,0) {};
\node at (0,0.8) {\textbf{Hommes}};
\node at (0,-0.5) {\textbf{Philosophes}};
\node[popdark] at (-1.5, -0.2) {Socrate};
\draw[->, popdark, thick] (-1.5,-0.2) -- (-0.5, -0.2);
\node[align=center, font=\tiny] at (3, -1.5) {Prémisse 1 : Tout P est H\\Prémisse 2 : S est P\\Conclusion : S est H\\$\rightarrow$ \textcolor{popgreen}{VALIDE}};
\end{scope}
\end{tikzpicture}
\legende{Diagrammes d'Euler : à gauche, le sophisme du moyen terme non distribué (Socrate dans H mais pas nécessairement dans P) ; à droite, le syllogisme valide (Barbara).}
\end{popfigure}

\begin{exoresolu}[Analyse d'un syllogisme : le cas de Socrate]
\textbf{Énoncé.} Un élève tient le raisonnement suivant : « Tous les philosophes sont des hommes ; or Socrate est un homme ; donc Socrate est un philosophe. » Ce raisonnement est-il valide ? Justifiez rigoureusement votre réponse.
\tcblower
\textbf{Solution détaillée.}
\begin{enumerate}
\item \textbf{Structure.} Prémisse 1 (Majeure) : « Tous les philosophes sont des hommes » (Tout P est H). Prémisse 2 (Mineure) : « Socrate est un homme » (S est H). Conclusion : « Socrate est un philosophe » (S est P).
\item \textbf{Termes.} Majeur (P) = Philosophe ; Mineur (S) = Socrate ; Moyen (M) = Homme.
\item \textbf{Test de validité (Règle du moyen terme).} Le moyen terme « Homme » est le sujet d'une universelle affirmative (Tout P est H) : il n'est \emph{pas distribué} (le sujet d'une A n'est pas distribué). Il est l'attribut d'une affirmative (S est H) : il n'est pas distribué non plus. \textbf{Le moyen terme n'est distribué dans aucune prémisse.} Règle violée $\rightarrow$ syllogisme invalide.
\item \textbf{Analyse ensembliste.} P $\subset$ H. S $\in$ H. On ne peut pas déduire S $\in$ P (Socrate pourrait être un mécanicien de Douala, un cultivateur de Bafoussam).
\item \textbf{Validité et vérité.} La conclusion est vraie dans les faits (Socrate était philosophe), mais elle n'est pas \emph{démontrée} par ce raisonnement : la forme est invalide. C'est le sophisme classique du \cle{moyen terme non distribué}.
\end{enumerate}
\textbf{Conclusion.} Le raisonnement est invalide : la conclusion, bien que vraie en fait, ne découle pas nécessairement des prémisses. Un raisonnement valide exige que la vérité de la conclusion soit garantie par la seule forme de l'inférence.
\end{exoresolu}

\courssub{L'argumentation : convaincre et persuader}

\begin{definition}[Argumentation]
L'\cle{argumentation} est l'activité discursive qui vise à faire admettre une thèse par un auditoire. Elle mobilise des \emph{arguments} (raisons, preuves, exemples) liés à la thèse par des \emph{liaisons} (connecteurs logiques, topoï). Elle se distingue de la démonstration par son public (universel vs particulier) et sa certitude (probable vs nécessaire).
\end{definition}

\begin{aretenir}[Convaincre vs Persuader (Pascal, \emph{Pensées})]
\begin{center}
\begin{tabularx}{\linewidth}{|l|X|X|}\hline
\textbf{Critère} & \textbf{Convaincre (Preuve rationnelle)} & \textbf{Persuader (Séduction affective)} \\ \hline
\textbf{Visée} & L'adhésion de la \emph{raison} (entendement) & L'adhésion de la \emph{sensibilité} (cœur, passions) \\ \hline
\textbf{Moyens} & Arguments logiques, faits vérifiables, définitions, syllogismes, démonstrations & Images, émotions, flatterie, menace, humour, autorité, répétition, témoignages \\ \hline
\textbf{Portée} & Universelle : valable pour tout esprit rationnel & Particulière : dépend du public, du contexte, de la culture \\ \hline
\textbf{Valeur épistémique} & Connaissance objective, vérité & Croyance, opinion, adhésion pratique \\ \hline
\textbf{Exemple} & Démonstration mathématique, diagnostic médical & Publicité, discours politique, plaidoirie d'avocat \\ \hline
\end{tabularx}
\end{center}
\textbf{Formule} : « Convaincre s'impose à la raison ; persuader séduit la sensibilité. » (Pascal : « L'art de persuader » vs « L'art de convaincre »).
\end{aretenir}

\begin{methode}[Démasquer un argument fallacieux (Sophisme)]
Face à une argumentation (débat, discours, publicité, réseaux sociaux) :
\begin{enumerate}
\item \textbf{Dégager la thèse} défendue et les arguments avancés.
\item \textbf{Classer chaque argument} : pertinent (porte sur la thèse) ou fallacieux (déplace le débat).
\item \textbf{Nommer le sophisme} :
  \begin{itemize}
  \item \emph{Argumentum ad hominem} : attaque la personne (origine, morale) au lieu de l'idée.
  \item \emph{Argumentum ad populum} : invoque l'opinion majoritaire (« tout le monde pense que... »).
  \item \emph{Argumentum ad baculum} : utilise la menace, la force, la peur.
  \item \emph{Argumentum ad verecundiam} : invoque une autorité non compétente en la matière.
  \item \emph{Pente glissante} : enchaînement catastrophiste sans preuve.
  \item \emph{Fausse analogie} : compare des choses non comparables sur le point litigieux.
  \item \emph{Généralisation abusive} : tire une règle universelle d'un cas unique.
  \item \emph{Pétition de principe} : suppose ce qu'il faut prouver.
  \end{itemize}
\item \textbf{Expliquer pourquoi} l'argument ne prouve rien : montrer le déplacement du débat (thèse $\rightarrow$ personne, émotion, nombre...).
\item \textbf{Reformuler} ce qu'il faudrait prouver pour que la conclusion soit établie.
\end{enumerate}
\textbf{Réflexe} : se demander toujours « cet argument prouve-t-il la thèse, ou cherche-t-il seulement à me faire céder ? ».
\end{methode}

\begin{exoresolu}[Analyse d'un discours persuasif : débat radio à Yaoundé]
\textbf{Énoncé.} Lors d'un débat à la radio de Yaoundé sur l'opportunité d'introduire l'informatique dès le collège, un intervenant déclare : « Tout le monde pense que c'est une bonne idée ; d'ailleurs celui qui s'y oppose est un ennemi du progrès qui veut maintenir notre jeunesse dans l'ignorance. » Identifiez les procédés fallacieux de ce discours et expliquez pourquoi ils ne constituent pas une démonstration.
\tcblower
\textbf{Solution détaillée.}
\begin{enumerate}
\item \textbf{Thèse défendue} : « Il faut introduire l'informatique dès le collège. »
\item \textbf{Premier procédé} : « Tout le monde pense que... » est un \cle{argumentum ad populum} (appel à la foule). Or la vérité n'est pas une question de nombre : au XVII\textsuperscript{e} siècle, « tout le monde » pensait que le Soleil tourne autour de la Terre. L'accord de la majorité ne prouve pas la vérité d'une thèse.
\item \textbf{Deuxième procédé} : « celui qui s'y oppose est un ennemi du progrès » est un \cle{argumentum ad hominem} (attaque personnelle) : au lieu de répondre aux arguments de l'adversaire (coût des équipements, formation des enseignants, pertinence pédagogique), on attaque sa personne et ses intentions. C'est un déplacement du débat.
\item \textbf{Troisième procédé} : l'accusation implicite de vouloir « maintenir la jeunesse dans l'ignorance » joue sur la peur et la culpabilité : c'est une pression affective (proche de l'\emph{ad baculum}), qui vise à faire taire l'objecteur sans le réfuter.
\item \textbf{Ce qu'il faudrait prouver} : une véritable argumentation devrait montrer, par des faits et des raisons, que l'introduction précoce de l'informatique améliore les apprentissages (études comparatives), qu'elle est financièrement réalisable (budget, maintenance) et que les enseignants peuvent être formés (plan de formation).
\end{enumerate}
\textbf{Conclusion.} Ce discours persuade sans convaincre : il manipule l'auditoire par l'opinion commune, l'attaque personnelle et la peur, au lieu d'établir rationnellement sa thèse. Seule une argumentation fondée sur des preuves pertinentes mérite l'adhésion de la raison.
\end{exoresolu}

\begin{methode}[Analyser un texte argumentatif : Convaincre ou Persuader ?]
\begin{enumerate}
\item \textbf{Repérer la visée} : le texte cherche-t-il l'adhésion de la \emph{raison} (convaincre) ou celle des \emph{sentiments} (persuader) ?
\item \textbf{Relever les indices lexicaux et syntaxiques} :
  \begin{itemize}
  \item \textbf{Convaincre} : vocabulaire technique/précis, définitions, enchaînements logiques (« donc », « parce que », « or »), exemples vérifiables, données chiffrées, renvois à des autorités scientifiques.
  \item \textbf{Persuader} : vocabulaire affectif/évaluatif, images, métaphores, apostrophe directe (« vous », « nous »), impératifs, exclamations, répétitions, témoignages personnels, mise en scène dramatique.
  \end{itemize}
\item \textbf{Évaluer la valeur} : convaincre s'adresse à l'universalité (valable pour tout esprit rationnel) ; persuader s'adresse à la particularité d'un public.
\item \textbf{Nuance} : un même discours peut mêler les deux (avocat, homme politique, prédicateur) ; il faut alors peser la part de chacun.
\end{enumerate}
\end{methode}

\begin{exoresolu}[Convaincre ou persuader ? Comparaison de deux textes]
\textbf{Énoncé.} Comparez les deux textes suivants et indiquez, en le justifiant, lequel relève de la persuasion et lequel de la conviction.
\begin{quote}
\textbf{Texte A :} « La déforestation réduit l'évapotranspiration ; or celle-ci alimente les pluies ; donc la destruction de la forêt du bassin du Congo menace le régime des pluies de toute la région. »
\textbf{Texte B :} « Sauvez notre belle forêt, berceau de nos ancêtres, avant que vos enfants ne pleurent sur une terre désolée ! »
\end{quote}
\tcblower
\textbf{Solution détaillée.}
\begin{enumerate}
\item \textbf{Texte A (Convaincre).} Il présente un enchaînement déductif : Cause (déforestation) $\rightarrow$ Mécanisme (réduction évapotranspiration) $\rightarrow$ Conséquence (menace sur pluies). Les termes sont précis (évapotranspiration, bassin du Congo), le raisonnement est vérifiable par la climatologie. Il s'adresse à la raison de tout lecteur, quel que soit son pays ou ses sentiments : c'est un discours qui \cle{convainc}.
\item \textbf{Texte B (Persuader).} Il utilise des images affectives (« belle forêt », « berceau de nos ancêtres »), une apostrophe directe (« vos enfants »), une mise en scène dramatique (« pleurent sur une terre désolée ») et un impératif d'urgence (« Sauvez ! »). Aucune preuve n'est fournie : on cherche à émouvoir pour obtenir l'adhésion immédiate : c'est un discours qui \cle{persuade}.
\item \textbf{Jugement.} Le texte A a une valeur universelle : sa force vient de la nécessité du raisonnement. Le texte B dépend de la sensibilité du public : il peut toucher l'un et laisser l'autre indifférent. Comme le note Pascal, on peut « entrer dans les esprits » par la raison ou par « l'agrément » ; mais seule la conviction rationnelle est un assentiment fondé sur la vérité.
\end{enumerate}
\textbf{Conclusion.} Le texte A relève de la conviction (preuve rationnelle), le texte B de la persuasion (séduction affective). Le philosophe doit préférer la première, sans nier que la seconde soit parfois utile pour faire agir.
\end{exoresolu}

\courssub{La démonstration}

\begin{definition}[Démonstration]
La \cle{démonstration} est un raisonnement déductif qui établit la vérité d'une proposition (thèse) à partir de prémisses certaines (définitions, axiomes, principes, faits établis) par une chaîne d'inférences nécessaires. Elle produit la \emph{science} (Aristote : \emph{epistémè}) par opposition à l'opinion (\emph{doxa}).
\end{definition}

\begin{propriete}[Conditions de la démonstration rigoureuse (Aristote, \emph{Analytiques postérieurs})]
\begin{enumerate}
\item \textbf{Prémisses vraies} : non supposées, connues par elles-mêmes (premiers principes) ou déjà démontrées.
\item \textbf{Prémisses premières} : immédiates (non déduites d'autres).
\item \textbf{Prémisses causes} : la cause de l'attribut dans le sujet (la prémisse majeure contient la raison \emph{pourquoi}).
\item \textbf{Prémisses plus connues que la conclusion} : antérieures dans l'ordre de la connaissance.
\item \textbf{Forme syllogistique valide} : déduction nécessaire.
\end{enumerate}
\end{propriete}

\begin{methode}[Construire une démonstration]
\begin{enumerate}
\item \textbf{Énoncer clairement la thèse} à démontrer (une phrase, sans ambiguïté).
\item \textbf{Définir les termes} : toute démonstration repose sur des définitions précises (genre + différence).
\item \textbf{Poser les prémisses} : propositions admises comme vraies (définitions, principes, faits établis, résultats antérieurs).
\item \textbf{Enchaîner rigoureusement} : chaque étape doit découler nécessairement de la précédente ; utiliser les connecteurs logiques (« or », « donc », « par conséquent », « dès lors »).
\item \textbf{Conclure} : la conclusion doit être exactement la thèse annoncée, ni plus, ni moins.
\item \textbf{Vérifier} : relire en contrôlant qu'aucune étape ne suppose déjà ce qu'on veut prouver (sinon : \cle{pétition de principe} / raisonnement circulaire).
\end{enumerate}
\textbf{Piège classique} : la pétition de principe consiste à prendre pour prémisse ce qu'on prétend démontrer (souvent masqué par une reformulation).
\end{methode}

\begin{exoresolu}[Construire et évaluer une démonstration]
\textbf{Énoncé.} a) Construisez une démonstration de la thèse suivante : « Le mensonge est incompatible avec la confiance nécessaire à la vie en société. » b) Le raisonnement « La liberté est le plus grand des biens parce que rien n'est plus précieux que la liberté » est-il une démonstration ? Justifiez.
\tcblower
\textbf{Solution détaillée.}
\textbf{a) Construction de la démonstration.}
\begin{enumerate}
\item \textbf{Définitions.} Le mensonge est l'affirmation délibérée d'une chose que l'on sait fausse, dans l'intention de tromper (définition augustinienne/kantienne). La vie en société repose sur des échanges (paroles, contrats, promesses, marchés) qui supposent que chacun peut croire la parole d'autrui : c'est la \cle{confiance} (fidélité à la parole donnée).
\item \textbf{Prémisse 1 (Principe anthropologique) :} Toute vie sociale organisée exige que la parole donnée soit, en principe, crue (sinon aucun marché à Mokolo, aucun contrat de travail, aucune promesse n'aurait de valeur ; la coopération devient impossible).
\item \textbf{Prémisse 2 (Analyse du mensonge) :} Or le mensonge, s'il se généralise ou devient la norme, détruit la présomption de vérité inhérente à la communication : on ne peut plus croire personne \emph{a priori}.
\item \textbf{Enchaînement (Modus Ponens) :} Si la confiance est condition nécessaire de la société (P1) et si le mensonge détruit cette confiance (P2), alors le mensonge détruit la condition de possibilité de la vie sociale.
\item \textbf{Conclusion} : Donc le mensonge est incompatible avec la confiance nécessaire à la vie en société. (On retrouve ici l'argument de Kant dans \emph{Fondements de la métaphysique des mœurs} : la maxime du mensonge universalisée se contredit elle-même, car plus personne ne croirait personne, rendant le mensonge lui-même inefficace).
\end{enumerate}
\textbf{b) Évaluation du raisonnement.} « La liberté est le plus grand des biens parce que rien n'est plus précieux que la liberté » : la prémisse (« rien n'est plus précieux que la liberté ») dit exactement la même chose que la conclusion (« la liberté est le plus grand des biens »), en d'autres termes (synonymie). Le raisonnement tourne en rond : c'est une \cle{pétition de principe} (ou raisonnement circulaire, \emph{circulus in demonstrando}). Il n'apporte aucune preuve nouvelle, aucune cause, aucune raison extérieure à la thèse.
\textbf{Conclusion.} Une démonstration authentique fait passer d'un connu certain (prémisses) à un inconnu qui en découle nécessairement (conclusion) ; dès que la conclusion est déjà contenue dans les prémisses sous une autre formulation, il n'y a plus démonstration mais tautologie.
\end{exoresolu}

\courssub{Applications : rédiger et justifier une démarche}

\begin{methode}[Traiter une situation concrète avec les outils logiques]
\begin{enumerate}
\item \textbf{Décrire la situation} en une phrase neutre (rumeur, débat familial, publicité, règlement scolaire, fake news).
\item \textbf{Extraire le raisonnement implicite} : quelles prémisses (explicites ou implicites) ? quelle conclusion ?
\item \textbf{Appliquer les opérations logiques fondamentales} :
  \begin{itemize}
  \item \emph{Définir} les termes clés (lever les équivoques).
  \item \emph{Diviser} / Classer les éléments (distinctions pertinentes).
  \item \emph{Juger} la vérité des prémisses (sources, vérification).
  \item \emph{Raisonner} : vérifier la validité de l'inférence (syllogisme, induction, analogie).
  \end{itemize}
\item \textbf{Évaluer globalement} : le raisonnement est-il valide ? les prémisses sont-elles vraies ? y a-t-il sophisme (équivoque, généralisation, ad hominem...) ?
\item \textbf{Proposer une attitude rationnelle} : vérifier les sources croisées, demander des preuves, reformuler correctement le raisonnement, suspendre le jugement si données insuffisantes.
\end{enumerate}
\textbf{Réflexe} : face à une rumeur ou une info anxiogène, toujours séparer deux questions : « Le raisonnement est-il correct (valide) ? » et « Les faits avancés (prémisses) sont-ils établis (vrais) ? ».
\end{methode}

\begin{exoresolu}[La rumeur au lycée technique de Bafoussam]
\textbf{Énoncé.} Dans un lycée technique de Bafoussam, une rumeur circule : « Le proviseur a annulé toutes les sorties pédagogiques ; or la sortie au barrage de Bamendjing est une sortie pédagogique ; donc notre sortie est annulée. » Un élève objecte : « Non, le proviseur n'a annulé que les sorties non encadrées. » En utilisant les notions de la leçon, montrez que la conclusion des élèves n'était pas assurée et dégagez la leçon logique de cette situation.
\tcblower
\textbf{Solution détaillée.}
\begin{enumerate}
\item \textbf{Structure du raisonnement initial (Syllogisme Barbara).}
  \begin{itemize}
  \item Prémisse 1 (Majeure) : « Toutes les sorties pédagogiques sont annulées » (Tout S est A).
  \item Prémisse 2 (Mineure) : « La sortie au barrage est une sortie pédagogique » (B est S).
  \item Conclusion : « La sortie au barrage est annulée » (B est A).
  \end{itemize}
  \textbf{Formellement}, ce syllogisme est \textbf{valide} : c'est le schéma correct de la 1\textsuperscript{re} figure (Barbara). La forme garantit que \emph{si} les prémisses sont vraies, la conclusion l'est nécessairement.
\item \textbf{Examen de la vérité des prémisses (Logique matérielle).} La validité ne suffit pas pour la solidité. La prémisse 1 est \textbf{fausse} : le proviseur n'a annulé que les sorties \emph{non encadrées}. La rumeur a commis une \cle{généralisation abusive} (erreur de \emph{division} : on a mal classé le genre « sortie pédagogique » en omettant l'espèce « sortie encadrée » vs « sortie non encadrée »). C'est une faute contre l'opération de division (non exhaustivité de la classification implicite).
\item \textbf{Correction de la prémisse et nouvelle déduction.}
  \begin{itemize}
  \item Prémisse 1 corrigée : « Toutes les sorties \emph{non encadrées} sont annulées » (Tout S' est A).
  \item Prémisse 2 : « La sortie au barrage est une sortie pédagogique » (B est S).
  \item Moyen terme : « sortie pédagogique » (S) vs « sortie non encadrée » (S'). \textbf{Le moyen terme a changé / n'est pas le même.} Pas de syllogisme possible.
  \item Pour conclure, il faut savoir : « La sortie au barrage est-elle encadrée ? » (B est S' ou B n'est pas S').
  \end{itemize}
\item \textbf{Leçon logique.} Une déduction n'est probante (solide) que si deux conditions sont réunies : (i) la forme est valide ; (ii) les prémisses sont vraies. La rumeur satisfaisait la première condition (forme Barbara) mais pas la seconde (prémisse majeure fausse par généralisation).
\end{enumerate}
\textbf{Conclusion.} Face à une information, l'esprit logique vérifie d'abord les faits (les prémisses, la définition des termes, la classification) avant d'accepter la conclusion, même quand le raisonnement est formellement correct : c'est l'attitude critique que la philosophie entend former.
\end{exoresolu}

\begin{methode}[Construire une argumentation de dissertation (Probatoire / Bac)]
\begin{enumerate}
\item \textbf{Analyser le sujet} : définir chaque terme clé (travail de \cle{définition} : genre + différence), repérer la présupposition éventuelle, dégager la tension (le problème).
\item \textbf{Formuler la problématique} sous forme d'une question précise, ouverte, qui commande tout le devoir (ni trop large, ni trop étroite).
\item \textbf{Construire le plan dialectique} (Thèse / Antithèse / Synthèse ou Dépassement) : chaque partie défend une thèse argumentée.
  \begin{itemize}
  \item \textbf{Argument logique} : raisonnement (déduction, induction, analogie, \emph{a fortiori}, \emph{reductio ad absurdum}).
  \item \textbf{Exemple concret} : fait historique, situation camerounaise (marché Mokolo, barrage Bamendjing, réforme éducative), référence culturelle.
  \item \textbf{Référence auteur} (si pertinent) : Aristote (logique, syllogisme), Pascal (convaincre/persuader, cœur/raison), Kant (devoir, universalisation, pétition de principe), Descartes (méthode, doute), Toulmin (modèle d'argumentation).
  \end{itemize}
\item \textbf{Rédiger chaque paragraphe} selon le schéma standard : \textbf{Idée directrice} (phrase thèse) $\rightarrow$ \textbf{Explication / Argumentation} (raisonnement) $\rightarrow$ \textbf{Exemple / Illustration} (concret, vérifiable) $\rightarrow$ \textbf{Lien avec la problématique} (transition interne).
\item \textbf{Justifier les transitions} entre parties : montrer en quoi la partie suivante répond à une insuffisance, une limite ou une objection de la précédente (démarche dialectique).
\item \textbf{Conclure} en répondant \emph{exactement} à la question posée, en synthétisant le parcours, sans ouvrir un nouveau problème.
\end{enumerate}
\textbf{Réflexe} : un paragraphe sans exemple est abstrait ; un exemple sans analyse n'est pas un argument. Citer un auteur sans expliquer sa pensée n'est pas un argument d'autorité valide.
\end{methode}

\begin{exoresolu}[Sujet de dissertation : « Faut-il toujours démontrer pour convaincre ? »]
\textbf{Énoncé.} Dégagez la problématique du sujet « Faut-il toujours démontrer pour convaincre ? » et rédigez l'introduction et la première partie du devoir.
\tcblower
\textbf{Solution détaillée.}
\textbf{Analyse du sujet.}
\begin{itemize}
\item \textbf{Démontrer} : établir une vérité par un enchaînement nécessaire de raisons à partir de principes certains (démonstration rigoureuse, au sens fort : mathématiques, logique) ou au sens large : apporter une preuve solide, des arguments décisifs.
\item \textbf{Convaincre} : obtenir l'adhésion rationnelle d'autrui (différent de persuader : adhésion affective).
\item \textbf{Toujours} : universalité de l'obligation (dans tous les domaines, en toutes circonstances). Le sujet oppose implicitement la démonstration (voie rigoureuse, idéale) à d'autres moyens d'obtenir l'assentiment rationnel (témoignage, autorité compétente, induction, arguments probabilistes, evidence immédiate).
\end{itemize}
\textbf{Problématique.} La démonstration (au sens strict de déduction nécessaire à partir de premiers principes) est-elle la \emph{seule} voie légitime et efficace pour obtenir l'adhésion de la raison, ou y a-t-il des domaines (morale, politique, histoire, vie courante, sciences expérimentales) où convaincre exige d'autres ressources que la preuve apodictique ?
\textbf{Introduction (rédaction type Bac).}
\begin{quote}
« Dans les mathématiques, nul n'accepte un théorème sans démonstration : la preuve y est la condition même de la vérité et de l'universalité. Mais dans la vie courante — au marché de Mokolo pour marchander, en politique pour voter, en morale pour juger une action, en histoire pour comprendre le passé —, nous adhérons à bien des idées sans qu'aucune démonstration au sens strict ne nous soit fournie. Faut-il le regretter comme un manque de rigueur ? La démonstration est-elle l'unique voie royale de la conviction, ou certains domaines échappent-ils par nature à l'exigence de la preuve déductive ? Nous verrons d'abord que la démonstration constitue l'idéal régulateur de la conviction rationnelle ; nous montrerons ensuite qu'elle connaît des limites structurelles (domaines, prémisses, urgence) ; nous soutiendrons enfin que convaincre autrui exige d'adapter la rigueur de la preuve à la nature du sujet, sans renoncer à l'exigence de rationalité. »
\end{quote}
\textbf{Première partie (Thèse : La démonstration, idéal de la conviction rationnelle).}
\begin{itemize}
\item \textbf{Idée directrice} : Seule la démonstration rigoureuse fonde une adhésion universelle, nécessaire et inébranlable, indépendante des opinions, des époques et des affects.
\item \textbf{Explication} : Fondée sur des définitions claires et des principes premiers (axiomes), elle enchaîne des inférences valides (syllogismes). Elle s'impose à tout esprit rationnel par sa forme même (Aristote, \emph{Analytiques postérieurs} : la science est démonstration par causes). Elle distingue le \emph{savoir} (\emph{epistémè}) de l'opinion (\emph{doxa}).
\item \textbf{Exemple concret} : Le théorème de Pythagore ($a^2+b^2=c^2$) s'impose avec la même nécessité à un élève de Garoua, un ingénieur de Tokyo ou un mathématicien de l'Antiquité, parce que sa preuve ne dépend d'aucune opinion particulière, d'aucune autorité, d'aucun affect. C'est l'archétype de la conviction partagée.
\item \textbf{Exemple d'auteur} : Descartes (\emph{Discours de la méthode}, Règle 2) : « Ne recevoir aucune chose pour vraie que je ne la connusse évidemment être telle » ; la démonstration (déduction claire et distincte) est le rempart contre l'erreur.
\item \textbf{Lien / Transition} : Si convaincre vise la vérité objective et non le simple consentement subjectif, la démonstration en est donc la forme accomplie et le modèle. Mais cette exigence mathématique peut-elle s'appliquer partout ? L'homme attend-il une preuve géométrique pour agir moralement ou voter ?
\end{itemize}
\textbf{Conclusion de l'exercice.} La problématique est dégagée, l'introduction suit le schéma accroche--définitions--problème--annonce du plan, et la première partie articule thèse, argument logique, exemple concret, référence auteur et transition : la démarche est rédigée et justifiée selon les attendus du Bac.
\end{exoresolu}

\begin{attention}[Les 5 erreurs qui coûtent des points au Bac (Probatoire / Baccalauréat Technique)]
\begin{enumerate}
\item \textbf{Confondre validité et vérité} : déclarer un raisonnement « faux » parce que sa conclusion est fausse, alors que la question porte sur la forme. \textbf{Toujours traiter séparément} : (1) La forme est-elle valide ? (2) Les prémisses sont-elles vraies ? Un raisonnement valide à prémisses fausses n'est pas « faux », il est \emph{non solide}.
\item \textbf{Affirmer le sophisme sans le nommer ni l'expliquer} : écrire « c'est un sophisme » ou « c'est fallacieux » ne suffit pas ; il faut \textbf{nommer le procédé} (\emph{ad hominem}, \emph{ad populum}, pétition de principe, équivoque, généralisation abusive, fausse analogie...) et \textbf{montrer précisément} en quoi il ne prouve pas la thèse (déplacement du débat).
\item \textbf{Confondre persuader et convaincre} : persuader n'est pas « convaincre fortement » ; c'est obtenir l'adhésion par les émotions, sans preuve rationnelle. Inverser les deux notions dans une définition ou une analyse de texte est éliminatoire en introduction ou en commentaire.
\item \textbf{Présenter une pétition de principe comme une démonstration} : répéter la conclusion dans les prémisses avec d'autres mots (synonymie, tautologie) n'est pas prouver. \textbf{Vérifier toujours} que la conclusion apporte une connaissance nouvelle par rapport aux prémisses.
\item \textbf{Illustrer sans analyser (catalogue d'exemples)} : empiler des exemples (Socrate, le marché Mokolo, la publicité, le barrage) sans dégager l'argument logique qu'ils illustrent. \textbf{Chaque exemple doit être} : introduit par une idée, décrit brièvement, \textbf{analysé} (en quoi prouve-t-il l'idée ?), et rattaché à la problématique.
\end{enumerate}
\end{attention}

\newpage

\coursec{Exercices}

\courssub{Je vérifie mes connaissances}

\begin{exercice}[QCM : notions de logique]
Pour chaque question, une seule réponse est correcte. Entoure la lettre correspondante.
\begin{enumerate}
\item Une \cle{proposition} est :
\begin{enumerate}
\item une phrase interrogative ;
\item un énoncé qui peut être vrai ou faux ;
\item un ordre donné à autrui ;
\item une exclamation de joie.
\end{enumerate}
\item Le connecteur « si \ldots alors \ldots » exprime :
\begin{enumerate}
\item la conjonction ; \item la disjonction ; \item l'implication ; \item la négation.
\end{enumerate}
\item Dans le syllogisme « Tous les hommes sont mortels ; or Socrate est un homme ; donc Socrate est mortel », la phrase « Socrate est un homme » est :
\begin{enumerate}
\item la conclusion ; \item la prémisse majeure ; \item la prémisse mineure ; \item le moyen terme.
\end{enumerate}
\item Un \cle{sophisme} est :
\begin{enumerate}
\item un raisonnement valide mais long ;
\item un raisonnement faux qui a l'apparence de la validité ;
\item une démonstration mathématique ;
\item une opinion personnelle.
\end{enumerate}
\end{enumerate}
\end{exercice}

\begin{exercice}[Vrai ou faux ? Justifie]
Réponds par vrai ou faux, puis justifie ta réponse en deux ou trois lignes.
\begin{enumerate}
\item La négation d'une proposition vraie est toujours une proposition fausse.
\item Le raisonnement inductif conduit à une conclusion absolument certaine.
\item Persuader consiste à toucher les émotions de l'auditoire, tandis que convaincre fait appel à la raison.
\item Une démonstration peut partir de prémisses fausses et aboutir à une conclusion vraie de manière valide.
\end{enumerate}
\end{exercice}

\begin{exercice}[Définitions]
Définis en deux phrases maximum chacune des notions suivantes, en donnant pour chacune un exemple tiré de la vie courante camerounaise :
\begin{enumerate}
\item la \cle{logique} ;
\item l'\cle{argumentation} ;
\item la \cle{démonstration} ;
\item le \cle{syllogisme}.
\end{enumerate}
\end{exercice}

\begin{exercice}[Tables de vérité directes]
Soient deux propositions : $p$ : « Il pleut sur Douala » et $q$ : « La route est mouillée ».
\begin{enumerate}
\item Recopie et complète la table de vérité de la conjonction $p \wedge q$ et de la disjonction $p \vee q$.
\item Construis la table de vérité de l'implication $p \Rightarrow q$.
\item À quelle condition l'implication $p \Rightarrow q$ est-elle fausse ? Donne un exemple concret.
\end{enumerate}
\end{exercice}

\courssub{Je m'entraîne}

\begin{exercice}[Formaliser un raisonnement]
On considère le raisonnement suivant, entendu dans un débat de quartier à Yaoundé : « Si le candidat gagne les élections, il devient président ; or il n'a pas gagné les élections ; donc il ne deviendra pas président. »
\begin{enumerate}
\item Pose $p$ : « le candidat gagne les élections » et $q$ : « le candidat devient président ». Formalise les deux prémisses et la conclusion.
\item Construis la table de vérité complète du raisonnement.
\item Montre qu'il existe un cas où les prémisses sont vraies et la conclusion fausse. Le raisonnement est-il valide ? Justifie.
\item Quel nom donne-t-on à ce type d'erreur logique (négation de l'antécédent) ?
\end{enumerate}
\end{exercice}

\begin{exercice}[À la chasse aux sophismes]
Voici cinq affirmations reconstituées à partir de débats publics camerounais. Pour chacune, identifie le sophisme commis et nomme-le (attaque personnelle, pente glissante, faux dilemme, généralisation hâtive, appel à la tradition).
\begin{enumerate}
\item « Ne l'écoutez pas parler d'économie : il n'a même pas réussi sa propre boutique à Mokolo. »
\item « Si on autorise les élèves à poser des questions au professeur, bientôt ils refuseront toute discipline, puis ce sera le chaos total dans le pays. »
\item « Deux choix seulement : soit tu votes pour moi, soit tu veux le malheur de ce village. »
\item « Mon voisin de taxi a été malhonnête avec moi : décidément, tous les taximen de cette ville sont des voleurs. »
\item « On a toujours célébré cette fête de cette manière ; il est donc interdit d'en changer quoi que ce soit. »
\end{enumerate}
\end{exercice}

\begin{exercice}[Construire et vérifier un syllogisme]
\begin{enumerate}
\item Construis un syllogisme valide (deux prémisses et une conclusion) sur le thème de la vie scolaire dans ton lycée technique (par exemple autour de la notion « élève assidu »).
\item Identifie le grand terme, le moyen terme et le petit terme de ton syllogisme.
\item Vérifie sa validité à l'aide d'un diagramme de Venn à trois cercles que tu traceras soigneusement.
\item Transforme ton syllogisme en syllogisme invalide en modifiant une prémisse, et explique où se situe l'erreur.
\end{enumerate}
\end{exercice}

\begin{exercice}[Démonstration par l'absurde guidée]
On veut démontrer par l'absurde la proposition suivante : « Il n'existe pas de plus grand nombre entier naturel. »
\begin{enumerate}
\item Rappelle en quoi consiste le raisonnement par l'absurde (quelles sont ses trois étapes ?).
\item Étape 1 : formule la négation de la proposition à démontrer (hypothèse de départ).
\item Étape 2 : en supposant qu'il existe un plus grand entier naturel $n$, considère le nombre $n+1$ et tire-en une contradiction.
\item Étape 3 : conclus en rédigeant la démonstration complète, en justifiant chaque étape par le nom de la règle logique utilisée.
\end{enumerate}
\end{exercice}

\courssub{Je me prépare au Bac}

\begin{exercice}[Dissertation — type Baccalauréat technique]
Sujet : \emph{« La force de l'argumentation suffit-elle à établir la vérité ? »}
\begin{enumerate}
\item Analyse le sujet : dégage les termes principaux (« force », « argumentation », « établir », « vérité »), définis-les et formule la problématique.
\item Rédige une introduction complète comportant : une amorce, la définition des termes, la problématique et l'annonce du plan.
\item Construis un plan dialectique en trois parties :
\begin{enumerate}
\item Thèse : une argumentation solide (prémisses vraies, raisonnement valide) permet d'établir la vérité (appuie-toi sur la démonstration en mathématiques et en sciences) ;
\item Antithèse : une argumentation peut être puissante sans être vraie (sophismes, rhétorique, persuasion dans les débats politiques et publicitaires) ;
\item Synthèse : la vérité exige à la fois la validité logique et la vérité des prémisses ; distingue convaincre et persuader.
\end{enumerate}
\item Rédige la conclusion de la dissertation.
\end{enumerate}
\end{exercice}

\begin{exercice}[Commentaire de texte — type Baccalauréat technique]
Texte : \emph{« Nous croyons connaître une chose au sens absolu lorsque nous croyons connaître la cause dont dépend cette chose, et savoir que cette cause est la cause de cette chose, et qu'il est impossible que la chose soit autrement. »} (Aristote, \emph{Seconds Analytiques}, adapté)
\begin{enumerate}
\item Présente brièvement l'auteur et l'œuvre dont est tiré le texte.
\item Dégage la thèse du texte : quelles sont, selon Aristote, les conditions du savoir démonstratif ?
\item Explique les expressions suivantes : « connaître la cause dont dépend cette chose » ; « il est impossible que la chose soit autrement ».
\item Montre, à l'aide d'un exemple mathématique simple (par exemple la somme des angles d'un triangle), comment une démonstration satisfait ces conditions.
\item Discussion : \emph{Tout peut-il se démontrer ?} Rédige un développement argumenté (au moins deux arguments et un exemple) sur cette question, en envisageant le cas des premiers principes indémontrables.
\end{enumerate}
\end{exercice}

\begin{exercice}[Problème de logique appliquée — type Probatoire]
Dans un lycée technique de Bafoussam, le conseil de discipline examine trois affirmations faites par des élèves. On note : $p$ : « L'élève a triché à l'épreuve » ; $q$ : « L'élève sera convoqué devant le conseil » ; $r$ : « L'élève sera sanctionné ».
\begin{enumerate}
\item Formalise les affirmations suivantes : (a) « Si l'élève a triché, alors il sera convoqué devant le conseil » ; (b) « Si l'élève est convoqué devant le conseil, alors il sera sanctionné » ; (c) « L'élève a triché ».
\item En appliquant deux fois la règle du \emph{modus ponens}, déduis la conclusion. Formalise ce raisonnement en chaîne ($(p \Rightarrow q)$, $(q \Rightarrow r)$, $p$, donc $r$).
\item Construis la table de vérité de $[(p \Rightarrow q) \wedge (q \Rightarrow r)] \Rightarrow (p \Rightarrow r)$ et montre que cette proposition est toujours vraie (tautologie).
\item Un élève déclare : « Je n'ai pas triché, donc je ne serai pas sanctionné. » Formalise ce raisonnement, montre à l'aide d'une table de vérité qu'il est invalide, et nomme l'erreur commise.
\item Rédige un court paragraphe (cinq à huit lignes) justifiant, à l'intention du conseil de discipline, la démarche correcte à suivre pour tirer une conclusion logiquement valide.
\end{enumerate}
\end{exercice}

\newpage

\coursec{Corrigés des exercices}

\begin{corrige}[Exercice 1 — QCM : notions de logique]
\begin{enumerate}
\item \textbf{Réponse b.} Une proposition est un énoncé déclaratif auquel on peut attribuer une valeur de vérité (vrai ou faux). Les questions, ordres et exclamations ne sont ni vrais ni faux : ce ne sont pas des propositions.
\item \textbf{Réponse c.} « Si \ldots alors \ldots » exprime l'\cle{implication} : $p \Rightarrow q$. La conjonction s'exprime par « et », la disjonction par « ou », la négation par « non ».
\item \textbf{Réponse c.} « Socrate est un homme » est la \cle{prémisse mineure} : elle range le petit terme (Socrate) sous le moyen terme (homme). La majeure est « Tous les hommes sont mortels » ; la conclusion est « Socrate est mortel ».
\item \textbf{Réponse b.} Un \cle{sophisme} est un raisonnement faux qui a l'apparence de la validité : il trompe par sa forme, alors que sa structure logique est défectueuse ou que ses prémisses sont inacceptables.
\end{enumerate}
\textbf{Points de vigilance :} ne pas confondre proposition (vrai/faux) et phrase (unité grammaticale) ; bien distinguer prémisse majeure (contient le grand terme) et prémisse mineure (contient le petit terme).
\end{corrige}

\begin{corrige}[Exercice 2 — Vrai ou faux ? Justifie]
\begin{enumerate}
\item \textbf{Vrai.} La négation inverse la valeur de vérité : si $p$ est vraie, alors $\neg p$ est fausse, et réciproquement. C'est le principe de la table de vérité de la négation.
\item \textbf{Faux.} L'induction généralise à partir de cas particuliers observés ; sa conclusion n'est que \emph{probable}, jamais absolument certaine. Exemple : avoir vu mille corbeaux noirs ne garantit pas que le mille-unième sera noir. Seule la déduction garantit la conclusion si les prémisses sont vraies.
\item \textbf{Vrai.} \cle{Persuader} vise l'adhésion par les émotions, les sentiments, l'imagination (rhétorique, publicité) ; \cle{convaincre} vise l'assentiment rationnel par des arguments et des preuves. La distinction est classique depuis Aristote (\emph{Rhétorique}).
\item \textbf{Vrai.} La validité ne concerne que la \emph{forme} du raisonnement, non la vérité des prémisses. Exemple : « Tous les poissons volent ; or le capitaine est un poisson ; donc le capitaine vole » est formellement valide, avec des prémisses fausses. Une démonstration \emph{solide} exige en outre des prémisses vraies.
\end{enumerate}
\textbf{Points de vigilance :} distinguer soigneusement \emph{validité} (forme correcte) et \emph{vérité} (contenu vrai) ; c'est la confusion la plus fréquente au Bac.
\end{corrige}

\begin{corrige}[Exercice 3 — Définitions]
\begin{enumerate}
\item \textbf{La logique} est la science des lois formelles du raisonnement correct ; elle étudie les conditions de validité des inférences, indépendamment du contenu des énoncés. \emph{Exemple :} au marché de Mokolo, « si ce sac coûte \num{1000} F et que j'ai \num{800} F, alors je ne peux pas l'acheter » : j'applique sans le savoir une règle logique.
\item \textbf{L'argumentation} est un discours qui vise à faire admettre une thèse à un auditoire en lui présentant des raisons (arguments). \emph{Exemple :} un élève de Terminale qui défend devant sa classe l'idée que le lycée technique prépare mieux à la vie professionnelle, en citant les stages en entreprise à Douala.
\item \textbf{La démonstration} est un raisonnement déductif rigoureux qui établit la vérité nécessaire d'une proposition à partir de prémisses vraies (axiomes, définitions, théorèmes déjà prouvés). \emph{Exemple :} démontrer en classe que la somme des angles d'un triangle vaut $180^{\circ}$, ce qui vaut pour tous les triangles, y compris ceux tracés au tableau du lycée.
\item \textbf{Le syllogisme} est un raisonnement déductif composé de deux prémisses et d'une conclusion, reliant trois termes (grand terme, moyen terme, petit terme). \emph{Exemple :} « Tous les élèves assidus réussissent ; or Amina est une élève assidue ; donc Amina réussira. »
\end{enumerate}
\textbf{Points de vigilance :} une définition doit donner le genre proche et la différence spécifique ; l'exemple doit être concret et distinct de la définition.
\end{corrige}

\begin{corrige}[Exercice 4 — Tables de vérité directes]
\begin{enumerate}
\item Table de vérité de la conjonction et de la disjonction :
\begin{center}
\begin{tabular}{|c|c|c|c|}
\hline
\rowcolor{popblueL} $p$ & $q$ & $p \wedge q$ & $p \vee q$ \\
\hline
V & V & V & V \\
\hline
V & F & F & V \\
\hline
F & V & F & V \\
\hline
F & F & F & F \\
\hline
\end{tabular}
\end{center}
La conjonction n'est vraie que si les deux propositions sont vraies ; la disjonction (ou inclusif) n'est fausse que si les deux sont fausses.
\item Table de vérité de l'implication :
\begin{center}
\begin{tabular}{|c|c|c|}
\hline
\rowcolor{popblueL} $p$ & $q$ & $p \Rightarrow q$ \\
\hline
V & V & V \\
\hline
V & F & F \\
\hline
F & V & V \\
\hline
F & F & V \\
\hline
\end{tabular}
\end{center}
\item L'implication $p \Rightarrow q$ n'est fausse que dans un seul cas : lorsque l'\emph{antécédent} $p$ est vrai et le \emph{conséquent} $q$ faux. \emph{Exemple concret :} « Il pleut sur Douala » est vrai, mais « la route est mouillée » est faux (par exemple une route couverte par un hangar) : alors l'affirmation « s'il pleut, la route est mouillée » est contredite. Si $p$ est faux, l'implication est automatiquement vraie (on ne peut pas la contredire).
\end{enumerate}
\textbf{Points de vigilance :} le cas « F $\Rightarrow$ V = V » surprend souvent ; retenir qu'une implication n'est réfutée que par un antécédent vrai suivi d'un conséquent faux.
\end{corrige}

\begin{corrige}[Exercice 5 — Formaliser un raisonnement]
\begin{enumerate}
\item Prémisse 1 : $p \Rightarrow q$. Prémisse 2 : $\neg p$. Conclusion : $\neg q$.
\item Table de vérité complète :
\begin{center}
\begin{tabular}{|c|c|c|c|c|}
\hline
\rowcolor{popblueL} $p$ & $q$ & $p \Rightarrow q$ & $\neg p$ & $\neg q$ \\
\hline
V & V & V & F & F \\
\hline
V & F & F & F & V \\
\hline
F & V & V & V & F \\
\hline
F & F & V & V & V \\
\hline
\end{tabular}
\end{center}
\item Examinons la troisième ligne : $p$ faux et $q$ vrai. Alors la prémisse 1 ($p \Rightarrow q$) est vraie, la prémisse 2 ($\neg p$) est vraie, mais la conclusion ($\neg q$) est fausse. Il existe donc un cas où les prémisses sont vraies et la conclusion fausse : le raisonnement est \textbf{invalide}. \emph{Illustration :} le candidat peut devenir président sans avoir gagné ces élections-là (par exemple en tant que vice-président succédant au président, ou lors d'une élection ultérieure).
\item Cette erreur s'appelle le \cle{sophisme de la négation de l'antécédent} : de « si $p$ alors $q$ » et « non $p$ », on ne peut rien conclure sur $q$.
\end{enumerate}
\textbf{Points de vigilance :} la forme valide correspondante serait le \emph{modus tollens} : $(p \Rightarrow q)$ et $\neg q$, donc $\neg p$. Bien vérifier chaque ligne de la table avant de conclure.
\end{corrige}

\begin{corrige}[Exercice 6 — À la chasse aux sophismes]
\begin{enumerate}
\item \textbf{Attaque personnelle} (\emph{argumentum ad hominem}) : on réfute la personne (l'échec de sa boutique) au lieu de réfuter ses arguments sur l'économie.
\item \textbf{Pente glissante} : on enchaîne des conséquences de plus en plus catastrophiques sans prouver que chaque étape entraîne nécessairement la suivante (questions $\rightarrow$ refus de discipline $\rightarrow$ chaos national).
\item \textbf{Faux dilemme} : on présente deux options comme si elles étaient les seules possibles, alors qu'il en existe d'autres (voter pour un autre candidat sans vouloir le malheur du village).
\item \textbf{Généralisation hâtive} : on tire une conclusion universelle (« tous les taximen ») à partir d'un seul cas particulier (un voisin malhonnête).
\item \textbf{Appel à la tradition} : on justifie une pratique uniquement par son ancienneté (« on a toujours fait ainsi »), ce qui ne prouve ni sa justesse ni l'interdiction de la modifier.
\end{enumerate}
\textbf{Points de vigilance :} nommer le sophisme ne suffit pas ; il faut expliquer en une phrase \emph{pourquoi} le raisonnement est défectueux.
\end{corrige}

\begin{corrige}[Exercice 7 — Construire et vérifier un syllogisme]
\begin{enumerate}
\item \textbf{Syllogisme valide :} « Tous les élèves assidus réussissent leur année ; or les élèves de la Terminale F3 de mon lycée sont assidus ; donc les élèves de la Terminale F3 de mon lycée réussiront leur année. »
\item \textbf{Termes :} le \cle{grand terme} est « réussir son année » (attribut de la conclusion) ; le \cle{petit terme} est « élèves de Terminale F3 » (sujet de la conclusion) ; le \cle{moyen terme} est « élèves assidus » (présent dans les deux prémisses, absent de la conclusion).
\item \textbf{Vérification par diagramme de Venn :} le cercle « élèves assidus » est entièrement inclus dans le cercle « ceux qui réussissent » (prémisse majeure) ; le cercle « élèves de Terminale F3 » est entièrement inclus dans le cercle « élèves assidus » (prémisse mineure). Par conséquent, le cercle « élèves de Terminale F3 » est nécessairement inclus dans le cercle « ceux qui réussissent » : la conclusion est inévitable, le syllogisme est valide.
\begin{popfigure}
\begin{tikzpicture}
\draw[popblue, thick] (0,0) circle (2.6cm);
\node[popblue] at (0,2.9) {Réussissent};
\draw[poporange, thick] (0,-0.2) circle (1.6cm);
\node[poporange] at (0,1.0) {Assidus};
\draw[popgreen, thick, fill=popgreenL] (0,-0.5) circle (0.7cm);
\node[popdark] at (0,-0.5) {\small Tle F3};
\end{tikzpicture}
\legende{Diagramme de Venn : le cercle « Tle F3 » est inclus dans « Assidus », lui-même inclus dans « Réussissent » : la conclusion est nécessaire.}
\end{popfigure}
\item \textbf{Syllogisme invalide :} « Tous les élèves assidus réussissent ; or Jean réussit ; donc Jean est assidu. » L'erreur : le moyen terme « réussir » n'est pas distribué correctement — le cercle de Jean est dans « réussissent », mais rien ne garantit qu'il soit dans « assidus » (Jean peut réussir par chance ou par talent sans assiduité). C'est le sophisme de l'\emph{affirmation du conséquent}.
\end{enumerate}
\textbf{Points de vigilance :} le moyen terme doit disparaître de la conclusion ; vérifier toujours que l'inclusion est forcée par les deux prémisses, pas seulement compatible avec elles.
\end{corrige}

\begin{corrige}[Exercice 8 — Démonstration par l'absurde guidée]
\begin{enumerate}
\item Le raisonnement par l'absurde comporte trois étapes : (1) on suppose vraie la \emph{négation} de la proposition à démontrer ; (2) on en tire par déduction rigoureuse une \emph{contradiction} (une absurdité) ; (3) on conclut que l'hypothèse de départ est fausse, donc que la proposition initiale est vraie.
\item \textbf{Étape 1 — Négation :} la proposition à démontrer est « il n'existe pas de plus grand entier naturel ». Sa négation est : « il existe un plus grand entier naturel ». Appelons-le $n$.
\item \textbf{Étape 2 — Contradiction :} supposons donc qu'il existe un entier naturel $n$ tel que tout entier naturel $m$ vérifie $m \leq n$. Or $n$ étant un entier naturel, $n+1$ est aussi un entier naturel (l'ensemble $\mathbb{N}$ est stable par addition de 1). Mais $n+1 > n$, ce qui contredit l'hypothèse selon laquelle $n$ est le plus grand. Nous obtenons une contradiction : $n$ serait à la fois le plus grand et dépassé par $n+1$.
\item \textbf{Étape 3 — Rédaction complète :} \emph{Démonstration.} Supposons, par l'absurde, qu'il existe un plus grand entier naturel $n$ (hypothèse de la négation). Considérons l'entier $n+1$ : c'est un entier naturel (définition de $\mathbb{N}$) et $n+1 > n$ (propriété de l'ordre dans $\mathbb{N}$). Donc $n$ n'est pas le plus grand entier, ce qui contredit notre hypothèse (règle de non-contradiction). L'hypothèse est donc fausse : par le raisonnement par l'absurde, il n'existe pas de plus grand nombre entier naturel. \emph{CQFD.}
\end{enumerate}
\textbf{Points de vigilance :} la contradiction doit être \emph{déduite} de l'hypothèse, pas simplement affirmée ; conclure en rejetant l'hypothèse, jamais la proposition à démontrer.
\end{corrige}

\begin{corrige}[Exercice 9 — Dissertation : « La force de l'argumentation suffit-elle à établir la vérité ? »]
\textbf{Barème indicatif (sur 20) :} introduction 4 pts ; analyse et problématique 2 pts ; développement (3 parties) 10 pts ; conclusion 2 pts ; langue et cohérence 2 pts.

\begin{enumerate}
\item \textbf{Analyse du sujet.} \emph{Force} : puissance de conviction d'un discours, sa capacité à emporter l'adhésion. \emph{Argumentation} : discours rationnel visant à faire admettre une thèse par des raisons. \emph{Établir} : fonder solidement, démontrer de manière définitive. \emph{Vérité} : conformité du discours à la réalité, ou validité démonstrative d'une proposition. \textbf{Problématique :} un discours peut-il, par sa seule puissance persuasive, garantir la vérité de ce qu'il affirme, ou la vérité exige-t-elle autre chose que la force rhétorique ?
\item \textbf{Introduction (modèle rédigé).} Dans les débats publics, en politique comme sur les marchés de nos villes, celui qui parle le mieux semble souvent avoir raison. L'argumentation, discours qui présente des raisons pour faire admettre une thèse, possède une force indéniable : celle de convaincre ou de persuader. Mais cette force suffit-elle à \emph{établir la vérité}, c'est-à-dire à garantir que la thèse défendue est conforme au réel ? Autrement dit : la puissance du discours peut-elle se substituer à la rigueur de la preuve ? Nous verrons d'abord qu'une argumentation solide, fondée sur des prémisses vraies et une déduction valide, permet effectivement d'établir la vérité ; nous montrerons ensuite qu'une argumentation puissante peut néanmoins défendre le faux ; nous conclurons enfin que la vérité exige à la fois la validité logique et la vérité des prémisses.
\item \textbf{Plan dialectique détaillé.}
\begin{enumerate}
\item \emph{Thèse.} La démonstration mathématique montre qu'un raisonnement valide à partir de prémisses vraies établit une vérité nécessaire (exemple : somme des angles d'un triangle égale à $180^{\circ}$). En sciences, l'argumentation expérimentale établit des vérités vérifiables (exemple : pasteurisation, vaccination). Aristote : la science est connaissance par les causes.
\item \emph{Antithèse.} Les sophismes ont l'apparence de la validité sans la réalité (faux dilemme, attaque personnelle). La rhétorique et la publicité persuadent sans prouver : un discours électoral émouvant peut faire adhérer à des promesses fausses. Platon, dans le \emph{Gorgias}, dénonce les sophistes qui font triompher le discours le plus fort, non le plus vrai.
\item \emph{Synthèse.} La force argumentative ne suffit pas : il faut distinguer \emph{convaincre} (par la raison, avec prémisses vraies et raisonnement valide) et \emph{persuader} (par les émotions). La vérité n'est établie que lorsque la validité formelle s'unit à la vérité matérielle des prémisses.
\end{enumerate}
\item \textbf{Conclusion (modèle rédigé).} La force de l'argumentation ne suffit donc pas à établir la vérité : elle peut tout aussi bien servir l'erreur lorsqu'elle ne s'appuie que sur l'émotion ou sur des raisonnements fallacieux. Seule une argumentation à la fois valide dans sa forme et vraie dans ses prémisses mérite le nom de démonstration. Reste ouverte une question : l'esprit humain accepte-t-il toujours la vérité démontrée, ou faut-il aussi former le cœur autant que la raison ?
\end{enumerate}
\textbf{Points de vigilance :} définir chaque terme dans l'introduction ; un exemple précis par argument ; ne pas opposer convaincre et persuader avant la synthèse ; éviter les exemples sans analyse.
\end{corrige}

\begin{corrige}[Exercice 10 — Commentaire de texte : Aristote, \emph{Seconds Analytiques}]
\textbf{Barème indicatif (sur 20) :} présentation 2 pts ; thèse 4 pts ; explication des expressions 4 pts ; exemple 4 pts ; discussion 4 pts ; langue 2 pts.

\begin{enumerate}
\item \textbf{Présentation.} Aristote (384--322 av. J.-C.), philosophe grec, disciple de Platon et maître d'Alexandre le Grand, est le fondateur de la logique formelle. Les \emph{Seconds Analytiques}, partie de l'\emph{Organon}, y exposent sa théorie de la démonstration scientifique (le syllogisme démonstratif).
\item \textbf{Thèse.} Selon Aristote, connaître véritablement (au sens absolu, \emph{épistémè}) une chose, c'est : (1) connaître la \emph{cause} dont elle dépend ; (2) savoir que cette cause est bien la cause de \emph{cette} chose (et non d'une autre) ; (3) savoir que la chose est \emph{nécessaire}, qu'elle ne peut pas être autrement. Le savoir démonstratif est donc un savoir par les causes et nécessaire.
\item \textbf{Explication des expressions.} « Connaître la cause dont dépend cette chose » signifie remonter du fait à sa raison d'être : je ne sais pas seulement \emph{que} la chose est, mais \emph{pourquoi} elle est. « Il est impossible que la chose soit autrement » exprime la \emph{nécessité} : la conclusion d'une démonstration n'est pas un fait contingent (qui aurait pu ne pas être), mais une vérité qui s'impose à tout esprit, en tout temps et en tout lieu.
\item \textbf{Exemple mathématique.} Je démontre que la somme des angles d'un triangle vaut $180^{\circ}$ en traçant la parallèle à un côté passant par le sommet opposé et en utilisant les angles alternes-internes. Je connais alors la \emph{cause} (les propriétés des parallèles et des angles) ; je sais que cette cause produit précisément ce résultat ; et je sais que tout triangle, nécessairement, vérifie cette propriété : il est impossible qu'un triangle euclidien ait une somme d'angles différente. Les trois conditions d'Aristote sont satisfaites.
\item \textbf{Discussion : tout peut-il se démontrer ?} \emph{Argument 1 :} non, car toute démonstration part de prémisses ; si chaque prémisse devait elle-même être démontrée, on tomberait dans une régression à l'infini, et rien ne serait jamais prouvé. \emph{Argument 2 :} il faut donc admettre des \emph{premiers principes indémontrables} : les axiomes (par exemple « le tout est plus grand que la partie ») et les définitions, évidents par eux-mêmes. Aristote lui-même soutient que la science repose sur des principes saisis par l'intuition intellectuelle (\emph{noûs}). \emph{Exemple :} en géométrie, les axiomes d'Euclide ne se démontrent pas ; c'est à partir d'eux que tout se démontre. \emph{Conclusion de la discussion :} tout ce qui se démontre suppose l'indémontrable ; la démonstration a des fondements qu'elle ne peut pas elle-même fonder.
\end{enumerate}
\textbf{Points de vigilance :} ne pas paraphraser le texte mais l'expliquer ; dans la discussion, envisager explicitement l'objection de la régression à l'infini ; citer correctement les trois conditions de la thèse.
\end{corrige}

\begin{corrige}[Exercice 11 — Problème de logique appliquée]
\begin{enumerate}
\item Formalisation : (a) $p \Rightarrow q$ ; (b) $q \Rightarrow r$ ; (c) $p$.
\item \textbf{Premier modus ponens :} de $p \Rightarrow q$ et $p$, on déduit $q$ (l'élève sera convoqué). \textbf{Second modus ponens :} de $q \Rightarrow r$ et $q$, on déduit $r$ (l'élève sera sanctionné). Raisonnement en chaîne :
\[ (p \Rightarrow q),\quad (q \Rightarrow r),\quad p \ \vdash\ q \ \vdash\ r \]
La conclusion est donc $r$ : « l'élève sera sanctionné ».
\item Table de vérité de $[(p \Rightarrow q) \wedge (q \Rightarrow r)] \Rightarrow (p \Rightarrow r)$ :
\begin{center}
\begin{tabular}{|c|c|c|c|c|c|c|c|}
\hline
\rowcolor{popblueL} $p$ & $q$ & $r$ & $p{\Rightarrow}q$ & $q{\Rightarrow}r$ & $(p{\Rightarrow}q)\wedge(q{\Rightarrow}r)$ & $p{\Rightarrow}r$ & Proposition \\
\hline
V & V & V & V & V & V & V & V \\
\hline
V & V & F & V & F & F & F & V \\
\hline
V & F & V & F & V & F & V & V \\
\hline
V & F & F & F & V & F & F & V \\
\hline
F & V & V & V & V & V & V & V \\
\hline
F & V & F & V & F & F & V & V \\
\hline
F & F & V & V & V & V & V & V \\
\hline
F & F & F & V & V & V & V & V \\
\hline
\end{tabular}
\end{center}
La dernière colonne ne contient que des V : la proposition est une \cle{tautologie}. Le raisonnement en chaîne (syllogisme hypothétique) est donc toujours valide.
\item Le raisonnement de l'élève se formalise : prémisse 1 (implicite) $p \Rightarrow r$ ; prémisse 2 $\neg p$ ; conclusion $\neg r$. Table de vérité :
\begin{center}
\begin{tabular}{|c|c|c|c|c|}
\hline
\rowcolor{popblueL} $p$ & $r$ & $p \Rightarrow r$ & $\neg p$ & $\neg r$ \\
\hline
V & V & V & F & F \\
\hline
V & F & F & F & V \\
\hline
F & V & V & V & F \\
\hline
F & F & V & V & V \\
\hline
\end{tabular}
\end{center}
À la troisième ligne ($p$ faux, $r$ vrai), les prémisses $p \Rightarrow r$ et $\neg p$ sont vraies, mais la conclusion $\neg r$ est fausse : le raisonnement est \textbf{invalide}. Concrètement, l'élève peut être sanctionné pour une autre faute que la tricherie. L'erreur commise est le \cle{sophisme de la négation de l'antécédent}.
\item \textbf{Paragraphe de justification.} Pour tirer une conclusion logiquement valide, le conseil de discipline doit procéder par \emph{modus ponens} : vérifier d'abord que la règle générale est établie (« si un élève a triché, alors il sera convoqué et sanctionné »), puis constater que l'antécédent est réalisé (« cet élève a effectivement triché »), et seulement alors affirmer le conséquent (« il sera sanctionné »). En revanche, il serait illogique de conclure, à partir de la seule innocence de l'élève (« il n'a pas triché »), qu'il ne sera pas sanctionné : d'autres motifs de sanction peuvent exister. La rigueur logique protège ainsi à la fois l'équité de la sanction et le droit de l'élève : on ne punit que sur preuve, et l'absence d'une faute ne prouve pas l'absence de toute faute.
\end{enumerate}
\textbf{Points de vigilance :} dans la table à trois variables, vérifier les huit lignes une à une ; ne pas confondre modus ponens (affirmation de l'antécédent, valide) et négation de l'antécédent (invalide) ; dans le paragraphe final, citer explicitement la règle logique utilisée.
\end{corrige}

\newpage

\coursec{Fiche bilan}

\begin{popfigure}
\begin{tikzpicture}[mindmap, grow cyclic, every node/.style={concept, font=\small},
root concept/.append style={concept color=popblue, font=\bfseries},
level 1/.append style={level distance=3.6cm, sibling angle=60, font=\scriptsize},
level 2/.append style={level distance=2.2cm, sibling angle=40, font=\tiny}]
\node[root concept]{Logique, argumentation et démonstration}
  child[concept color=poporange]{ node {Nature et enjeux}
    child { node {Science du raisonnement valide} }
    child { node {Penser juste, éviter les sophismes} }
  }
  child[concept color=popgreen]{ node {Opérations logiques}
    child { node {Définir, classer, diviser} }
    child { node {Jugement : quantité et qualité} }
    child { node {Raisonner} }
  }
  child[concept color=poppurple]{ node {Déduction et syllogisme}
    child { node {Majeure, mineure, conclusion} }
    child { node {Validité $\neq$ vérité} }
  }
  child[concept color=poppink]{ node {Argumentation}
    child { node {Convaincre : la raison} }
    child { node {Persuader : les émotions} }
    child { node {Thèse, arguments, exemples} }
  }
  child[concept color=popgold]{ node {Démonstration}
    child { node {Déduire de prémisses certaines} }
    child { node {Rigueur et nécessité} }
  }
  child[concept color=popblueL]{ node {Applications}
    child { node {Rédiger et justifier une démarche} }
    child { node {Dissertation et commentaire} }
  };
\end{tikzpicture}
\legende{Carte mentale de la leçon : les six branches essentielles de la logique, de l'argumentation et de la démonstration.}
\end{popfigure}

\begin{bilan}
\textbf{Définitions clés à connaître par cœur}
\begin{itemize}
\item \cle{Logique} : science qui étudie les lois et les règles du raisonnement valide, indépendamment du contenu des pensées (Aristote, fondateur).
\item \cle{Concept} : représentation générale et abstraite d'un objet ou d'une classe d'objets (ex. : « homme », « justice »).
\item \cle{Jugement} : opération qui affirme ou nie un rapport entre deux concepts (ex. : « l'homme est mortel »).
\item \cle{Raisonnement} : enchaînement de jugements permettant de passer de prémisses à une conclusion.
\item \cle{Syllogisme} : raisonnement déductif à trois propositions (majeure, mineure, conclusion) et trois termes (grand, moyen, petit).
\item \cle{Argumentation} : discours visant à faire adhérer un auditoire à une thèse ; elle \emph{convainc} par la raison ou \emph{persuade} par les émotions.
\item \cle{Démonstration} : raisonnement rigoureux qui tire nécessairement une conclusion de prémisses certaines (axiomes, définitions, théorèmes).
\item \cle{Sophisme} : raisonnement faux qui a l'apparence de la validité (ex. : l'argument d'autorité, la pente glissante).
\end{itemize}

\textbf{Repères essentiels}
\begin{center}
\begin{tabularx}{\linewidth}{|l|X|X|}
\hline
\rowcolor{popblueL} \textbf{Notion} & \textbf{Structure} & \textbf{Exemple} \\
\hline
Syllogisme & Majeure ; mineure ; conclusion & Tout homme est mortel ; or Socrate est un homme ; donc Socrate est mortel. \\
\hline
Déduction & Du général au particulier & De la loi au cas concret. \\
\hline
Induction & Du particulier au général & Des faits observés à une loi. \\
\hline
Convaincre & Preuves rationnelles & Débat scientifique, tribunal. \\
\hline
Persuader & Émotions, images, rythme & Publicité, discours politique. \\
\hline
\end{tabularx}
\end{center}

\textbf{Méthodes express}
\begin{itemize}
\item \emph{Analyser un raisonnement} : repérer la conclusion, identifier les prémisses, vérifier la validité (la forme) puis la vérité (le contenu).
\item \emph{Construire une argumentation} : thèse claire $\rightarrow$ arguments hiérarchisés $\rightarrow$ exemples concrets $\rightarrow$ réponse aux objections.
\item \emph{Rédiger et justifier une démarche} : annoncer les étapes, enchaîner par des connecteurs logiques (donc, or, cependant, par conséquent), conclure sans introduire d'idée nouvelle.
\end{itemize}
\end{bilan}

\begin{aretenir}[Checklist avant le Bac]
\begin{itemize}
\item[$\square$] Je sais définir la logique et dire à quoi elle sert.
\item[$\square$] Je distingue concept, jugement et raisonnement.
\item[$\square$] Je connais les opérations logiques fondamentales : définir, classer, diviser, juger, raisonner.
\item[$\square$] Je sais construire et analyser un syllogisme (majeure, mineure, conclusion).
\item[$\square$] Je distingue validité (forme) et vérité (contenu) d'un raisonnement.
\item[$\square$] Je différencie déduction, induction et déduction hypothétique.
\item[$\square$] Je sais distinguer convaincre (raison) et persuader (émotion).
\item[$\square$] Je reconnais les sophismes courants : argument d'autorité, généralisation hâtive, attaque personnelle.
\item[$\square$] Je sais ce qu'est une démonstration : prémisses certaines et conclusion nécessaire.
\item[$\square$] Je cite au moins un auteur (Aristote) et un exemple camerounais concret (débat, palabre, plaidoirie).
\item[$\square$] Je sais structurer une argumentation écrite : thèse, arguments, exemples, objections, conclusion.
\item[$\square$] J'utilise des connecteurs logiques pour justifier ma démarche dans une dissertation.
\end{itemize}
\end{aretenir}

\findecours

\end{document}
